\documentclass{SciPost}

\binoppenalty=10000
\relpenalty=10000

\hypersetup{
    colorlinks,
    linkcolor={red!50!black},
    citecolor={blue!50!black},
    urlcolor={blue!80!black}
}

\usepackage[bitstream-charter]{mathdesign}
\urlstyle{same}

\DeclareSymbolFont{usualmathcal}{OMS}{cmsy}{m}{n}
\DeclareSymbolFontAlphabet{\mathcal}{usualmathcal}

\fancypagestyle{SPstyle}{
\fancyhf{}
\lhead{\colorbox{scipostblue}{\bf \color{white} ~SciPost Physics Lecture Notes }}
\rhead{{\bf \color{scipostdeepblue} ~Submission }}
\renewcommand{\headrulewidth}{1pt}
\fancyfoot[C]{\textbf{\thepage}}
}

\usepackage[T1]{fontenc}
\usepackage{newunicodechar}
\newunicodechar{ž}{\v{z}}
\newunicodechar{Ž}{\v{Z}}
\newunicodechar{ñ}{\~n}
\newunicodechar{ö}{\"o}
\newunicodechar{–}{--}

\emergencystretch=2em
\tolerance=1000

\usepackage{float}
\usepackage{cleveref}
\usepackage[dvipsnames,table]{xcolor}
\usepackage[linesnumbered,ruled,vlined]{algorithm2e}

% \usepackage{showlabels}

% Math %
%%%%%%%%

\usepackage{physics}
\usepackage{bm}

\newcommand{\mc}{\mathcal}
\newcommand{\mbb}{\mathbb}
\newcommand{\msf}{\mathsf}
\newcommand{\sss}{\scriptstyle}
\newcommand{\ssss}{\scriptscriptstyle}
\newcommand{\rU}{{\rm U}}
\newcommand{\Vect}{\mathsf{Vec}}
\newcommand{\Rep}{\mathsf{Rep}}
\newcommand{\Fib}{\mathsf{Fib}}
\newcommand{\Fun}{\mathsf{Fun}}
\newcommand{\cat}{\lhd}
\newcommand{\act}{\rhd}
\newcommand{\mF}{{}^{\act\!} F}
\newcommand{\nF}{{}^{\cat\!} F}
\newcommand{\dF}{{}^{\circ\!} F}
\newcommand{\rF}{{}^{\bowtie\!} F}
\newcommand{\upup}[1]{\raisebox{0.5pt}[0pt]{$\sss{#1}$}}
\newcommand{\dodo}[1]{\raisebox{-1.2pt}[0pt]{$\sss{#1}$}}
\newcommand{\CC}[6]{\Big[{}^{\upup{#1}}_{\dodo{#4}}\, {}^{\upup{#2}}_{\dodo{#5}} \big| {}^{\upup{#3}}_{\dodo{#6}}\Big]}

\newcommand{\RCR}[3]{{\textcolor{BlueViolet}{#1}\textcolor{BrickRed}{#2}\textcolor{BlueViolet}{#3}}}
\newcommand{\RDR}[3]{{\textcolor{BlueViolet}{#1}\textcolor{black}{#2}\textcolor{BlueViolet}{#3}}}

\newcommand{\ZtZt}{\mbb Z_2\times\mbb Z_2}

\renewcommand{\vec}[1]{{\bf #1}}

% Figures %
%%%%%%%%%%%
\usepackage{graphicx}
\usepackage{tikz}
\usetikzlibrary{shapes.geometric,arrows,matrix,calc,scopes,decorations.markings,intersections,arrows.meta,decorations.pathreplacing,angles,quotes,patterns,decorations.pathmorphing}
\AddToHook{env/tikzpicture/begin}{\tracinglostchars=0}

\newcommand{\inputtikz}[1]{%
  	% \tikzsetnextfilename{#1}%
	\input{figs/#1.tex}%
}
\tikzstyle{obj1}=[line width = 0.8pt, line join=round, line cap=round]
\tikzstyle{obj2}=[line width = 0.8pt, line join=round, line cap=round, color=BrickRed]
\tikzstyle{obj3}=[line width = 0.8pt, line join=round, line cap=round, color=BlueViolet]
\tikzstyle{mpstensor}=[draw=black, line width = 0.7pt, fill=BlueViolet!14, circle, inner sep=3pt]
\tikzstyle{mpotensor}=[draw=black, line width=.7pt, fill=gray!20, rectangle, inner sep=3.8pt, rounded corners=2pt]
\tikzstyle{matrix}=[draw=black, line width = 0.7pt, fill=black!10, circle, inner sep=1.6pt]
\tikzstyle{mor}=[draw=none, fill=gray!22]

\tikzset{singlet/.style={decorate, decoration=snake}}

\tikzset{->1-/.style={decoration={
			markings,
			mark=at position #1 with {\arrow{Triangle[fill=black, width=3.5pt, length=3.5pt]}}},postaction={decorate}}}
\tikzset{-<1-/.style={decoration={
			markings,
			mark=at position #1 with {\arrowreversed{Triangle[fill=black, width=3.5pt, length=3.5pt]}}},postaction={decorate}}}
			
\tikzset{->2-/.style={decoration={
			markings,
			mark=at position #1 with {\arrow{Triangle[fill=BrickRed, width=3.5pt, length=3.5pt]}}},postaction={decorate}}}
\tikzset{->3-/.style={decoration={
			markings,
			mark=at position #1 with {\arrow{Triangle[fill=BlueViolet!80, width=3.5pt, length=3.5pt]}}},postaction={decorate}}}

\tikzset{
    mask point/.style={ transform shape, sloped, 
    minimum width=20pt, minimum height=4pt, inner sep=0pt, ultra thin, font=\tiny}
}
\tikzset{
    clip even odd rule/.code={\pgfseteorule},
    invclip/.style={clip,insert path=[clip even odd rule]{
    [reset cm](-\maxdimen,-\maxdimen)rectangle(\maxdimen,\maxdimen)
}}}
\newcommand\clipIntersection[1]{
    (#1.north west) -- (#1.north east) -- (#1.south east) -- (#1.south west) -- (#1.north west)
}
\newcommand\clipAroundNode[2]{
    \begin{pgfscope}
        \begin{scope}[overlay]
            \path[invclip] \clipIntersection{#1};#2
        \end{scope}
    \end{pgfscope}
}
\newcommand{\middleCoo}[6]{
    \path (#1) --coordinate[pos=0.5](#3) (#2);
    \path (#1) --coordinate[pos={0.5-\sp}](#4) (#2);
    \path (#1) --coordinate[pos={0.5+\sp}](#5) (#2);
}
\newcommand{\barycentre}[8]{
    \middleCoo{#1}{#2}{A}{}{}{}{#8};
    \middleCoo{#2}{#3}{B}{}{}{}{#8};
    \middleCoo{#1}{#3}{C}{}{}{}{#8};
    \path[name path=PA#3] (A) -- (#3);
    \path[name path=PB#1] (B) -- (#1);
    \path[name path=PC#2] (C) -- (#2);
    \path [name intersections={of=PA#3 and PB#1,by={#4}}];
    \path (#1) --coordinate[pos=0.84](#5) (#4);
    \path (#2) --coordinate[pos=0.84](#6) (#4);
    \path (#3) --coordinate[pos=0.84](#7) (#4);
}
\newcommand{\barycentreSq}[9]{
    \middleCoo{#1}{#2}{A1}{A2}{A3}{};
    \middleCoo{#3}{#4}{B1}{B2}{B3}{};
    \middleCoo{#1}{#3}{C1}{C2}{C3}{};
    \middleCoo{#2}{#4}{D1}{D2}{D3}{};
    \middleCoo{A1}{B1}{#5}{}{};
    \path[name path=PA2B2] (A2) -- (B2);
    \path[name path=PA3B3] (A3) -- (B3);
    \path[name path=PC2D2] (C2) -- (D2);
    \path[name path=PC3D3] (C3) -- (D3);       
    \path [name intersections={of=PA2B2 and PC2D2,by={#6}}];
    \path [name intersections={of=PA2B2 and PC3D3,by={#8}}];
    \path [name intersections={of=PA3B3 and PC2D2,by={#7}}];
    \path [name intersections={of=PA3B3 and PC3D3,by={#9}}];
}

% \usetikzlibrary{external}
% \tikzexternalize[prefix=cache/]
% \tikzexternalize

\title{张量网络讲义（Les Houches Lecture Notes on Tensor Networks）}

\begin{document}

\pagestyle{SPstyle}

\begin{center}{\Large \textbf{\color{scipostdeepblue}{
张量网络讲义（Les Houches Lecture Notes on Tensor Networks）
}}}\end{center}

\begin{center}\textbf{
Bram Vancraeynest-De Cuiper\textsuperscript{1$\star$},
Weronika Wiesiolek\textsuperscript{2} and
Frank Verstraete\textsuperscript{1,2$\dagger$}
}\end{center}

\begin{center}
{\bf 1} Department of Physics and Astronomy, Ghent University, Krijgslaan 299, 9000 Gent, Belgium
\\
{\bf 2} Department of Applied Mathematics and Theoretical Physics, University of Cambridge,
Wilberforce Road, Cambridge, CB3 0WA, United Kingdom
\\[\baselineskip]
$\star$ \href{mailto:Bram.VancraeynestDeCuiper@UGent.be}{\small Bram.VancraeynestDeCuiper@UGent.be}\,,\quad
$\dagger$ \href{mailto:fv285@cam.ac.uk}{\small fv285@cam.ac.uk}
\end{center}

\section*{\color{scipostdeepblue}{摘要}}
\textbf{\boldmath{张量网络为分类与模拟量子物质的关联相及拓扑相（topological phases）提供了一个强大的全新框架。其核心前提是：通过研究强关联物质底层的纠缠（entanglement）结构及其伴随的（广义）对称性（generalised symmetries），能够对强关联物质获得最佳理解。从本质上讲，张量网络为强关联系统中复杂的真空涨落提供了一种压缩的全息（holographic）描述，由此打破了臭名昭著的多体指数墙（exponential wall）。本讲义简要概述了该理论最重要的概念、计算与数学方面。
}}

\vspace{\baselineskip}

\vspace{10pt}
\noindent\rule{\textwidth}{1pt}
\tableofcontents
\noindent\rule{\textwidth}{1pt}
\vspace{10pt}

\section*{引言}
我们当前正在见证理解与研究量子多体物理方式的一场范式转变（paradigm shift）。在过去一个世纪的大部分时间里，\emph{量子多体问题}（quantum many-body problem）一直是理论物理与量子化学的核心问题，而与之相伴的指数墙（exponential wall）阻碍了描述强相互作用物质相（phases of matter）的进展。然而在过去的二十年中，基于\emph{纠缠}（entanglement）理论来模拟与分类这些强关联系统（strongly correlated systems）取得了显著突破，其典范代表正是\emph{张量网络}（tensor network）的概念。

从本质上讲，张量网络通过提供多体波函数的\emph{全息}（holographic）描述，绕过了指数级的量子多体墙。张量网络并非在指数级巨大的希尔伯特空间（Hilbert space）中描述波函数，而是对系统不同部分之间纠缠的路由方式进行建模。张量网络之所以取得成功，其根本原因在于：计算低能波函数任意可观测量所需的全部信息，都可以压缩至局域张量（local tensor）中，其中所谓的\emph{纠缠自由度}（entanglement degrees of freedom）构成了\emph{虚}（virtual）希尔伯特空间。能够做到这一点的根源在于纠缠熵\emph{面积律}（area law）的存在：张量网络正是对所有满足此类面积律的态构成的流形（manifold）的简洁而精确的表示。通过将定义在指数级巨大物理希尔伯特空间中的哈密顿量（Hamiltonian）投影到该低维流形上，构建用于模拟大规模量子多体系统的数值算法成为了可能。DMRG（密度矩阵重整化群，\emph{Density Matrix Renormalisation Group}）、MPS（矩阵乘积态，\emph{Matrix Product State}）与 PEPS（投影纠缠对态，\emph{Projected Entangled-Pair State}）算法正在彻底改变我们模拟原子物理前沿实验、当今量子计算机以及量子多体系统低能物理的方式。

不仅如此：自 Landau 与 Anderson 的开创性见解以来，人们便认识到多体物理的核心在于\emph{对称性}（symmetries）与\emph{对称性破缺}（symmetry breaking）。事实证明，整个系统的全局、局域及拓扑对称性，都反映在张量网络局域张量的对称性之中，而这些局域张量对整个多体波函数提供了极其有效的压缩。因此，对有能隙（gapped）\emph{物质相}、特别是\emph{拓扑相}（topological phases）的分类，便归结为对张量在一组相关对称性下不可约变换方式的分类问题。由于这些对称性作用在虚（纠缠）自由度上，并无理由要求它们必须构成一个群，从而涌现出可以用\emph{融合范畴}（fusion categories）来描述的更为一般的结合代数结构（associative structures）。事实上，张量网络提供了一种显式方法，通过 MPO（矩阵乘积算符，\emph{Matrix Product Operator}）代数及其高维类似物 PEPO（投影纠缠对算符，\emph{projected entangled-pair operators}）来表示这些\emph{广义对称性}（generalised symmetries）以及它们之间的对偶性（dualities）。在数学物理、凝聚态物理、统计物理以及高能物理中，有若干最活跃的研究方向正致力于探索这一深刻见解所带来的推论：强关联与拓扑物质相的全息描述，是唯一能够揭示其内在结构的描述途径。

\bigskip\noindent 本讲义基于 Frank Verstraete 于 2025 年 8 月在 Les Houches 物理暑期学校上以“\emph{精确可解性与量子信息}”（\emph{Exact Solvability and Quantum Information}）为题所作的五次讲座，涵盖了张量网络若干基础的计算与理论成果。本讲义并不旨在对这些主题进行详尽无遗的综述，而是作为深入探索的引导；读者亦可参考现存大量优秀的综述与讲义文献：~\cite{Verstraete2008,Schollwoeck2010,Bridgeman2016,Vanderstraeten2019,Cirac2020,Xiang2023,Banuls2022,Verstraete2023,Chen2025}。
\section{第一讲：动机与 MPS 流形}
\subsection{双指数墙与希尔伯特空间的幻象}\label{sec:Illusion}
在深入本讲义之前，认识量子多体问题所带来的真正挑战至关重要。正如人们在每一门量子力学课程中所学到的那样，希尔伯特空间的维度随构成粒子数呈指数级增长。这与经典情形其实并无太大不同——经典比特可能构成的字符串数目同样随比特数呈指数级增长。然而，由于量子力学允许态的\emph{叠加}（superposition）并由此产生\emph{纠缠}（entanglement），一个自然的问题不是希尔伯特空间由多少个基底态张成，而是希尔伯特空间中的态究竟有多少能够通过源于局域（少体）相互作用的某种物理机制被实际探索，特别是借助量子计算机进行探索。

更具体且精确地说，让我们考虑文献~\cite{Poulin2011}中提出的如下问题：在 $N$ 个量子比特的希尔伯特空间中，从一个参考态出发运行量子计算机，在随系统尺寸呈多项式增长的时间内，能够到达指数级巨大的多体希尔伯特空间的哪一部分？为了对此进行估算，我们考虑希尔伯特空间中围绕点 $\ket{\psi}$ 的所谓态的\emph{$\varepsilon$-球}（$\varepsilon$-balls），即考虑满足下式的态 $\ket{\varphi}$：
\begin{equation}
    \vert\braket{\psi}{\varphi}\vert \gtrsim 1 - \varepsilon.
\end{equation}
一个简单的几何论证表明，这类 $\varepsilon$-球所占希尔伯特空间总容积的比例按 $\mc O\big(\varepsilon^{2^N}\big)$ 标度，即仅占整个希尔伯特空间极其微小的\emph{双指数}（double-exponentially）（!）比例！由于这一标度行为，任何物理系统（包括量子计算机）要与特定的 $\varepsilon$-球产生重叠，所需时间至少随系统尺寸呈指数增长。我们由此得出结论：绝大多数态都是非物理的，对于描述物理系统完全无关紧要——希尔伯特空间只是一个（便利的）幻象，因为那些需要耗费指数级漫长时间才能制备的多粒子态，无法通过任何合理的物理过程制备出来。

\bigskip\noindent 这一引人注目的结果表明，我们在理解量子多体系统方面发生了一场范式转变，即物理波函数是那些具有低复杂度的波函数。更具体地说，我们期望量子多体系统的基态存在于一个物理上相关的低复杂度态\emph{流形}（manifold）之上，出于所有实际目的，我们希望用\emph{少量}参数来参数化该流形，即参数数量随粒子数呈多项式标度。粗略地说，流形是局部类似于平坦欧几里得空间的空间。这意味着对于流形上的每一个点，都可以关联一个线性\emph{切空间}（tangent space），其维度即为该流形的维度。
狄拉克（Dirac）的（诸多）伟大洞见之一在于：人们可以将流形解释为态的变分集合，从而完全无需离开该流形便可在其上开展物理研究。让我们将一般的态变分流形记为
\begin{equation}
    \mc M = \big\{\ket{\Psi(\vec z)} \ \big\vert\ \vec z \in \mbb C^n\big\}.
\end{equation}
我们将任意点 $\vec z$ 处的线性切空间记作 $T_\vec z \mc M$。该切空间中的一般态具有形式
\begin{equation}
    \sum_i x^i \pdv{}{z^i} \ket{\Psi(\vec z)},
\end{equation}
其中 $x^i$ 为复系数。现在假设我们希望计算某一受多体哈密顿量 $H$ 演化的态 $\ket{\Psi(\vec z)}$ 的时间演化：
\begin{equation}
    i\pdv{}{t} \ket{\Psi(\vec z)} = H\ket{\Psi(\vec z)}.
\end{equation}
此时 $i\pdv{}{t} \ket{\Psi(\vec z)}=i\sum_i\dot z^i\pdv{}{z^i} \ket{\Psi(\vec z)}$ 可以被视作 $T_\vec z \mc M$ 中的一个切向量。然而，对于一般的哈密顿量，嵌入 $\mc M$ 的环境希尔伯特空间中的向量 $H\ket{\Psi(\vec z)}$ 指向切平面\emph{之外}，因而不再属于该变分类。狄拉克的提议是在该流形上定义一个作用量，但出于我们的目的，更方便的做法是将其视作 $H\ket{\psi(\vec z)}$ 在流形切空间 $T_\vec z \mc M$ 上的正交投影。从几何上看，可以将这种投影理解为最小化如下距离：
\begin{equation}
    \min_{\{\dot z^i\}_i} \norm{i\sum_i\dot z^i\pdv{}{z^i} \ket{\Psi(\vec z)} - H \ket{\Psi(\vec z)}}.
\end{equation}
将点 $\vec z$ 处的正交投影算符记为 $\mc P_{T_\vec z}$，我们可以将其对 $H\ket{\psi(\vec z)}$ 的作用图示如下：
\begin{equation}
    \inputtikz{TangentPlane} \,\, .
\end{equation}
由此得到的\emph{含时变分原理}（time-dependent variational principle, TDVP）方程
\begin{equation}
    i\pdv{}{t} \ket{\Psi(\vec z)} = \mc P_{T_\vec z}H\ket{\Psi(\vec z)},
\end{equation}
便描述了流形 $\mc M$ 上一般多体态的一组高度非线性（源于切空间投影算符的出现）微分方程。重要的是，可以证明流形 $\mc M$ 是\emph{辛}（symplectic）流形。这特别意味着能量的期望值、对称性生成元以及其他运动常数在 TDVP 方程演化下是严格守恒的。

\bigskip\noindent 这种求解多体问题的流形方法已在众多领域获得了极其非凡的成功。事实上，毫不夸张地说，自薛定谔方程被发现以来，攻克量子多体问题的所有进展，本质上都在于寻找这类承载相关物理过程的良好态流形。\emph{矩阵乘积态}（matrix product state, MPS）及其高维推广——如二维\emph{投影纠缠对态}（projected entangled-pair state, PEPS）——便构成了这样的态类，我们稍后将对其进行介绍。但在介绍之前，不妨花些时间从流形视角讨论其他几类对某些读者而言可能更为熟悉的变分态类。
\paragraph{自由费米子与 Hartree-Fock}
在多费米子体系中，最广为人知的变分流形大概是\emph{Slater 行列式}（Slater determinant）流形~\cite{Slater1929}。Slater 行列式用 $N$ 个单粒子轨道 $\psi_i$（$i=1,\dots,N$）参数化 $N$ 费米子波函数，表示为如下形式的行列式：
\begin{equation}
    \Psi({\bf x}_1,\dots,{\bf x}_N) =
    \begin{vmatrix}
    \psi_1({\bf x}_1) & \dots & \psi_N({\bf x}_1) \\
    \psi_1({\bf x}_2) & \dots & \psi_N({\bf x}_2) \\
    \vdots & \ddots & \vdots \\
    \psi_1({\bf x}_N) & \dots & \psi_N({\bf x}_N)
    \end{vmatrix}.
\end{equation}
根据其构造，Slater 行列式在交换任意两个费米子时显式反对称，与 Pauli 不相容原理（Pauli exclusion principle）相契合。此外，请注意表示 Slater 行列式所需的参数数量随电子数\emph{线性}增长，因此 Slater 行列式仅张成希尔伯特空间中一个指数级微小的子流形。用于近似求解 $N$ 费米子薛定谔方程的 \emph{Hartree-Fock} 方法，正是定义在该态类上的自洽变分方法：通过将 TDVP 原理应用于该流形，可以获得含时与定态薛定谔方程的近似解。该方法的成功无论怎样强调都不为过。它使我们能够以极高精度模拟分子和材料，而这在传统精确方法下是完全不可行的。事实上，Hartree-Fock 方法使我们能够从第一性原理理解元素周期表结构与能带结构~\cite{Slater1935,Hartree1947}。
\paragraph{相互作用玻色子的场相干态} 与 Slater 行列式相对应的玻色体系对应物同样存在，即\emph{场相干态}（field coherent state）流形。对于玻色创生算符\footnote{在本段中，我们将破例在算符上方添加帽子以避免与标量函数混淆。} $\hat \psi^\dagger(\vec x)$，对于任意函数 $\varphi (\vec x):\mbb R^3\rightarrow\mbb C$，相干态具有形式
\begin{equation}
    \ket{\Psi[\varphi]} = \exp(\int d^3\vec x \varphi(\vec x)\hat\psi^\dagger(\vec x)) \ket{\Omega},
\end{equation}
其中 $\ket{\Omega}$ 是由 $\hat\psi(\vec x)\ket{\Omega}=0$（$\forall\vec x$）所刻画的玻色真空。回想一下，相干态是湮灭算符的精确本征态，即 $\hat\psi(\vec x)\ket{\Psi[\varphi]} = \varphi(\vec x)\ket{\Psi[\varphi]}$。现在我们将 TDVP 应用于哈密顿量
\begin{equation}
    H = \int d^3\vec x \hat \psi^\dagger(\vec x) \bigg( -\frac{1}{2m}\nabla^2 + V(\vec x)\bigg)\hat\psi(\vec x) + \frac{U}{2} \hat\psi^\dagger(\vec x)\hat\psi^\dagger(\vec x)\hat\psi(\vec x)\hat\psi(\vec x),
\end{equation}
该哈密顿量描述处于外势 $V$ 并具有接触相互作用的稀薄玻色气体。此时可以发现，TDVP 方程直接导出了（含时）\emph{Gross-Pitaevskii} 方程：
\begin{equation}
    i\pdv{}{t}\varphi(\vec x) = -\frac{\nabla^2}{2m}\varphi(\vec x) + V(\vec x)\varphi(\vec x) + U\vert\varphi(\vec x)\vert^2\varphi(\vec x).
\end{equation}
尽管这一推导非常简练，Gross-Pitaevskii 方程已被极其成功地用于描述玻色-爱因斯坦凝聚（Bose-Einstein condensate）的行为~\cite{GP}。

\bigskip\noindent 这里我们仅强调了在指数级巨大的希尔伯特空间中定义和参数化变分流形的若干最常用且最成功的方法。除此之外，我们略去了诸如\emph{耦合簇方法}（coupled cluster method）~\cite{CC}等大量后 Hartree-Fock 方法。
\subsection{量子自旋系统：单配性与阻挫}
量子晶格系统为一般的量子多体系统或量子场论提供了一种极其有用且深刻的理想化模型。在这一框架下，晶格是对空间的离散化，其顶点（某些情况下也包括边）上承载着量子自由度。从量子场论的角度看，这是规避紫外发散（UV-divergences）最可控的方式~\cite{Kogut1979,Kogut1982}。从凝聚态物理与原子物理的角度看，晶格系统为相关自由度提供了有效的（压缩的）描述。

在本讲义中，我们将关注一维或二维空间量子多体晶格模型的低能物理。将研究限制在低能扇区通常是合理的，因为在大多数实际材料中，费米温度约为 $10^4\,\text{K}$ 的量级，远高于实验的典型温度标度（若考虑电子系统，甚至室温也显得非常冷）。同样在量子场论中，当计算散射实验的截面等物理量时，基态及其低能激发态通常正是人们最感兴趣的态。通常，所考虑的希尔伯特空间是 qudit 的张量积，我们将其记作 $\mc H=(\mbb C^d)^N\equiv\bigotimes_{i=1}^N \mbb C^d$；除非另有说明，我们将考虑热力学极限（thermodynamic limit）$N\rightarrow \infty$，同时忽略有关极限顺序的细微差别。然而请注意，相应希尔伯特空间的维度作为系统尺寸的函数仍然是指数级的，因此晶格系统同样受到这种多体指数墙的困扰。

由于\emph{阻挫}（frustration）的存在，获取此类模型基态的良好近似通常极非显而易见。在此语境下，阻挫指的是基态无法同时最小化哈密顿量各项能量的特性。举例而言，考虑 $(1+1)$ 维中典型的自旋 1/2 反铁磁 Heisenberg 链：
\begin{equation}\label{eq:Heisenberg}
    H = \sum_{\expval{\msf i,\msf j}} \vec{S}_\msf i \cdot \vec{S}_\msf j.
\end{equation}
若只考虑两个自旋，基态将是由密度矩阵 $\rho=\ketbra{\phi}$ 给出的单态，其中 $\ket{\phi}=\frac{1}{\sqrt{2}}(\ket{\uparrow\downarrow} - \ket{\downarrow\uparrow})$。
对于具有 $N$ 个自旋的自旋链，若基态的所有约化密度矩阵 $\rho_{\msf i,\msf i+1}$ 都等于该单态，总能量显然将被最小化——但这显然是不可能的：\emph{纠缠单配性}（monogamy of entanglement）决定了这样的态不可能存在~\cite{Coffman1999,Osborne2006}。事实上，不存在任何三体态 $\rho_{ABC}\geq 0$，使得 $\Tr_C(\rho_{ABC}) = \Tr_B(\rho_{ABC})=\ketbra{\phi}$。因此，纠缠单配性必然导致阻挫。

纠缠单配性还可用于解释为何随着空间维度的增加，平均场理论（mean field theory）变得越来越精确。考虑一个具有大配位数的晶格而非一维自旋链，其哈密顿量同样是所有近邻自旋间反铁磁 Heisenberg 相互作用之和。直观上，在这种情况下，对于任意给定自旋，其与其他自旋的纠缠必须分散在如此多的近邻之间，以至于基态必然非常接近可分态 $\otimes_i \ket{\psi_i}$，从而使平均场理论变得精确。虽然这里的讨论略显粗略，但借助\emph{量子 de Finetti 定理}（quantum de Finetti theorem）~\cite{Raggio1989}可以使该论证变得严格。该定理刻画了在不同子系统置换下保持不变的多体密度矩阵，并在证明量子密码学协议的安全性中发挥着重要作用。

从上述讨论中我们推导出，基态不可能使其所有 $\rho_{\msf i, \msf i+1}$ 均为自旋单态。现在我们转而寻找一种在某种意义上\emph{尽可能接近}自旋单态的态。为此，人们可以尝试将\emph{共振价键态}（resonating valence bond state）作为基态的 ansatz：
\begin{equation}
    \ket{\Psi} = \ket{\Psi_+} + \ket{\Psi_-}, \quad {\rm where} \quad
    \ket{\Psi_\pm} = \prod_\msf i \big( \ket{\uparrow\downarrow}_{2\msf i,2\msf i\pm 1} - \ket{\downarrow\uparrow}_{2\msf i,2\msf i\pm 1} \big).
\end{equation}
图形上可以将该态表示为
\begin{equation}
    \inputtikz{RVB},
\end{equation}
其中波浪线代表近邻格点之间的自旋单态。事实上，$\ket{\Psi_\pm}$ 张成了 Majumdar 和 Ghosh 在文献~\cite{Majumdar1969}中引入的模型的二维基态子空间。注意到该态具有我们所期望的所有对称性：它具有 ${\rm SU}(2)$ 对称性，并具有平移、时间反演和空间反射不变性。这非常重要，因为基态（或更确切地说是基态子空间）应当继承哈密顿量的所有对称性。更重要的是，它具有非零的纠缠，因此我们预期该态的能量比平均场近似更接近真实的基态能量。尽管如此，事实证明该态的能量对于 Heisenberg 模型的真实基态能量而言仍然是一个相当差的近似，因为它大幅高估了基态能量。

至此应该清楚的是，写出具有正确对称性且能作为我们感兴趣的多体量子哈密顿量基态 ansatz 的纠缠态是极其困难的。然而，价键态暗示了一种特定的构造途径。受 Majumdar-Ghosh 模型的启发，Affleck、Kennedy、Tasaki 与 Lieb（AKLT）于 1987 年提出了一种由成对纠缠的自旋 1/2 粒子构建的态，这些粒子在每个晶格格点上被投影为一个自旋 1 粒子~\cite{Affleck1987}。他们的构造如下：从一条自旋 1/2 粒子链出发，其中每隔一对近邻粒子处于自旋单态。在每个格点 $\msf i$ 上，将最右边的粒子与下一个格点 $\msf i+1$ 上最左边的粒子结合在一起，形成一个\emph{物理}（physical）格点。构成物理格点的两个\emph{虚}（virtual）自旋 1/2 粒子随后被投影到它们的自旋 1 子空间上。该态的图形表示更为直观明了：
\begin{equation}
    \ket{\Psi_{\rm AKLT}} = \inputtikz{AKLT} \, ,
\end{equation}
其中
\begin{equation}
    \inputtikz{AKLT_proj}.
\end{equation}
正如上述作者所证明的，该态是一个与自旋 1 Heisenberg 模型密切相关的哈密顿量的基态，并且被认为属于同一个（对称保护）相（在\cref{sec:PhasesOfMatter}中将进一步讨论）。由于两个近邻格点恰好涉及一个自旋单态，构成这些物理自由度的 4 个组成自旋 1/2 粒子不可能处于耦合的自旋 2 态。因此，该态被两个自旋 1 构成的自旋 2 子空间上的投影算符所湮灭。由此，该态是 AKLT 哈密顿量的基态：
\begin{equation}\label{eq:AKLTHam}
    H = \sum_\msf i \vec{S}_\msf i \cdot \vec{S}_{\msf i+1} + \frac{1}{3} \big( \vec{S}_\msf i \cdot \vec {S}_{\msf i+1} \big)^2,
\end{equation}
其中局域两体项在相差一个常数平移和前置系数的意义下，可以被识别为自旋 2 投影算符。
\subsection{张量网络与面积律}\label{sec:AreaLaw}
在 AKLT 提出几年后的 1992 年，Fannes、Nachtergaele 和 Werner~\cite{Fannes1990} 利用量子 Markov 链的思想引入了一维\emph{有限关联态}（finitely correlated states）类，推广了 AKLT 态的构造。\emph{有限关联态}这一名称源于以下事实：这些态已被证明具有指数衰减的连通关联函数，正如稍后在\cref{sec:NormalObs}中所展示的。他们还证明了所有这些态都是有能隙无阻挫哈密顿量的基态。由文献~\cite{Delgado2004}首次倡导的现代观点是：通过从近邻格点之间最大纠缠的\emph{虚} qudit 链出发，并通过线性映射将其投影到每个格点上的\emph{物理}自旋来构造这些态。这种线性映射通常被称为\emph{投影算符}（即使它不必是幂等线性映射），这也是\emph{投影纠缠对态}（projected entangled-pair states）名称的由来，尽管该名称现在通常用于其高维推广。这些投影算符包含作为该态\emph{变分}自由度的参数。该态类与（无限且平移不变的）均匀\emph{矩阵乘积态}（MPS）类完全一致，而 MPS 正是用来刻画它们的最常用名称。让我们对此做更明确的说明。MPS 的构造从考虑近邻格点 $\msf i,\msf i+1$（$\msf i\in\mbb Z$）之间的最大纠缠态出发：
\begin{equation}
    \ket{I}_\msf i = \sum_{\alpha=1}^\chi \ket{\alpha\alpha}_{\msf i, \msf i+1}
\end{equation}
整数 $\chi$ 称为\emph{键维}（bond dimension），正如我们将看到的，它直接控制变分参数的数量。随后我们在链的每个物理格点上考虑如下线性映射：
\begin{equation}
    A = \sum_{i,\alpha,\beta} A_{\alpha\beta}^i \ket{i}\otimes\bra{\alpha}\otimes\bra{\beta} \in \mbb C^{d\chi^2},
\end{equation}
并将其作用在纠缠对上，从而张量 $A$ 所对应的态由下式给出：
\begin{equation}
    \ket{\Psi[A]} = \big(\bigotimes_{\msf i\in\mbb Z} A\big) \cdot \big(\bigotimes_{\msf i\in\mbb Z} \ket{I}_{\msf i}\big).
\end{equation}
我们称 $A$ \emph{生成}（generate）了 MPS $\ket{\Psi[A]}$。展开此式中的乘积，便直接在热力学极限下得到了更为人熟知的\emph{均匀}（uniform）矩阵乘积态表达式：
\begin{equation}\label{eq:MPS}
    \ket{\Psi[A]} = \sum_{\{i_j\}} \big(\cdots A^{i_j}A^{i_{j+1}}\cdots \big) \ket{\dots, i_j,i_{j+1},\dots}.
\end{equation}
波函数的这一表达式解释了\emph{矩阵乘积}态的名称，因为波函数振幅正是作为矩阵序列 $\{A^i\}_i$ 的乘积计算出来的。在均匀情形下，变分参数的数量为 $d\chi^2$。然而请注意，均匀 MPS 张成的态流形的维度严格小于 $d\chi^2$，因为对于任意非奇异矩阵 $X$，如下形式的变换：
\begin{equation}
    A^i \mapsto XA^iX^{-1}, \quad \forall i=1,2,\dots,d,
\end{equation}
（称为 \emph{gauge 变换}）均使态 $\ket{\Psi[A]}$ 保持不变。对该流形更精确的刻画留待下一小节讨论。

请注意，我们原本也可以在每个格点上选择不同的张量（投影算符），从而得到一般的非均匀 MPS，其键维可以在晶格的不同边上变化。通过这种方式，我们可以轻松处理具有开边界、周期性边界或扭曲边界条件的有限自旋链。需要指出的是，对于均匀情形，我们并未指定无穷远处的边界条件。对于具有\emph{单射性}（injective，下文定义）的一般矩阵乘积态，这些边界条件的选择是无关紧要的。在整个讲义中，我们几乎完全专注于热力学极限下满足该单射性性质的矩阵乘积态。为完整起见，我们在此给出有限自旋链上的一般 MPS。开链上的一般 MPS 可以写为：
\begin{equation}\label{eq:OpenMPS}
    \ket{\Psi[A]} = \sum_{\{i_j\}} \vec v_L^\dagger A^{i_1}[1]A^{i_2}[2]\cdots A^{i_N}[N]\vec v_R \ \ket{i_1,i_2,\dots,i_N}.
\end{equation}
其中向量 $\vec v_{L/R}$ 编码边界条件，方括号内的整数索引物理格点。另一方面，具有 $N$ 个格点的周期链上的平移不变态具有如下形式：
\begin{equation}
    \ket{\Psi[A]} = \sum_{\{i_j\}} {\rm Tr}\big(A^{i_1}A^{i_2}\cdots A^{i_N}\big) \ket{i_1,i_2,\dots,i_N}.
\end{equation}

\bigskip\noindent 按照处理张量网络时的惯例，我们将使用一种图形语言，其中单个 MPS 张量表示为：
\begin{equation}
    \inputtikz{MPSTensor} = \ A_{\alpha\beta}^i, \qquad \substack{i=1,2,\dots,d\\ \alpha,\beta=1,2,\dots,\chi}
\end{equation}
连接两个张量的腿表示对相应指标的收缩。遵循这些约定，我们可以将 MPS \eqref{eq:MPS} 表示为如下图形：
\begin{equation}
    \ket{\Psi[A]} = \inputtikz{MPS} \ .
\end{equation}
通常，在可以从上下文中推断出 MPS 张量时，我们不会在图中显式标出它们。

\bigskip\noindent 既然已经引入了矩阵乘积态这一变分类，下面让我们论证为何这些态能够很好地近似局域有能隙哈密顿量的基态~\cite{Verstraete2006a}。从物理上看，这源于有能隙基态所满足的纠缠熵\emph{面积律}（area law）。在一般维度下，纠缠熵的面积律意味着区域 $\mc A$ 与其补区域 $\bar{\mc A}$ 之间的两分纠缠熵（定义为约化密度矩阵 $\rho_\mc A={\rm Tr}_{\bar{\mc A}}(\rho)$ 的冯·诺依曼熵，即 $S=-{\rm Tr}(\rho_\mc A\log(\rho_\mc A))$）并不随区域 $\mc A$ 的大小（即\emph{体积}，volume）标度，而是随其边界 $\partial\!\mc A$ 的大小（即\emph{面积}，area）标度。在一维情形下，链的连通子区域的边界只是一个常数，在此设定下面积律于 2007 年由 Hastings 证明~\cite{Hastings2007}。正如我们接下来将展示的，MPS 内在地满足面积律。

为此，我们利用 MPS 的图形演算。在此演算中，让我们计算对均匀 MPS 迹化掉系统左半部分时的约化密度矩阵。这可以图示为
\begin{equation}
    \rho_R = {\rm Tr}_L\big(\ketbra{\Psi[A]}\big) = \inputtikz{MPSReducedDensity_1}\,\,,
\end{equation}
其中垂直虚线表示左（$L$）右（$R$）子系统之间的纠缠分割截面（entanglement cut）。在约化密度矩阵的计算中起着关键作用，并且也出现在如关联函数计算（见\cref{sec:NormalObs}）中的一个对象是\emph{转移矩阵}（transfer matrix）
\begin{equation}
    \mbb T_A = \sum_i A^i\otimes \bar A^i,
\end{equation}
我们在上图中也用虚线方框将其标出。让我们将其谱分解\footnote{更一般地，可能需要 Jordan 块，但这对接下来的讨论并不重要。}写为
\begin{equation}\label{eq:TransferSpec}
    \mbb T_A = \sum_{i=1}^{\chi^2} \lambda_i \ketbra{\rho_i}{\lambda_i},
\end{equation}
其中本征值按绝对值递减排列。这些本征向量显然是完全正线性映射 $\$_L(X)=\sum_i A^{i \dagger} X A^i$ 和 $\$_R(X)=\sum_i A^i X A^{i \dagger}$ 的不变子空间。若这些映射是遍历的（ergodic），量子 Perron-Frobenius 定理指出模最大的本征值将是唯一的、实且正的。此外，当重新排列为矩阵形式（具有 2 个指标）时，相应的左本征向量 $\ket{\rho_1}\equiv\ket{\rho}$ 与右本征向量 $\ket{\lambda_1}\equiv\ket{\Lambda}$ 是正定的。拥有唯一起主导作用的本征向量这一性质等价于上文提及的单射性概念。单射性的另一等价刻画如下：一个 MPS 是单射的，当且仅当存在一个有限的 $L$（称为\emph{单射长度}，injectivity length），使得矩阵集合 $\{A^{i_1}A^{i_2}\dots A^{i_L}\}$ 构成所有 $\chi\times\chi$ 矩阵的一组基。这恰恰也是保证该 MPS 是其有能隙 parent Hamiltonian 的\emph{唯一}基态所需的性质。对于一般 MPS，情况总是如此：非单射 MPS 构成的集合测度为零，因为这仅在矩阵 $A^i$ 拥有共同不变子空间时才可能发生。

为确保态被正确归一化为 1，我们总是归一化 MPS 张量 $A$，使得相应本征值 $\lambda_1$ 等于 1（参见\cref{sec:NormalObs}）。根据这些考量，可知约化密度矩阵可以表示如下：
\begin{equation}\label{eq:ReducedDensity}
    \rho_R = \inputtikz{MPSReducedDensity_2}.
\end{equation}
由于 $\rho_R$ 是一个 $d^L\times\chi$ 矩阵与一个 $\chi\times d^L$ 矩阵的乘积（其中 $L$ 为右侧格点的总数），我们得出结论：约化密度矩阵 $\rho_R$ 的秩的上界为 $\chi$。因此纠缠熵满足上界 $S \leq \log \chi$，因为这是秩为 $\chi$ 的态的最大熵。文献~\cite{Verstraete2006a,Hastings2007}中的定理确保了其逆命题亦成立：只要一个态满足纠缠熵面积律，就存在一个能忠实近似该态的 MPS。在下一讲中，我们将更详细地研究一般矩阵乘积态的纠缠性质。
\subsection{矩阵乘积算符}\label{sec:MPO}
在整个讲义中，我们还将使用线性算符的张量网络表示。均匀\emph{矩阵乘积算符}（matrix product operator, MPO）~\cite{Verstraete2004c,Pirvu2010}写为
\begin{equation}
    \sum_{\{i_k,j_k\}} \big(\cdots M^{i_k,j_k}M^{i_{k+1},j_{k+1}}\cdots \big) \ketbra{\dots, i_k,i_{k+1},\dots}{\dots, j_k,j_{k+1},\dots},
\end{equation}
图形上表示为
\begin{equation}
    \inputtikz{MPO}\,\, .
\end{equation}
类比于 MPS，向开边界条件以及非均匀情形的推广是十分直接的。

MPO 长期以来在统计物理领域得到了广泛应用，其中特定类型的 MPO 被称为转移矩阵（关于统一步骤的讨论见文献~\cite{Haegeman2016}）。与 MPS 的情况类似，当矩阵 $\{M^{ij}\}_{i,j}$ 没有公共不变子空间且张成整个矩阵代数时，我们称该 MPO 为单射的。

然而，非单射 MPO 也被广泛使用，例如用于对量子自旋链的哈密顿量进行编码。例如，让我们考虑开边界条件下的横场 Ising 模型哈密顿量：
\begin{equation}\label{eq:TFIM}
    H = -\sum_{\msf i=1}^{L-1} \sigma^Z_{\msf i}\sigma^Z_{\msf i+1} - \lambda\sum_{\msf i=1}^L \sigma^X_{\msf i}.
\end{equation}
该哈密顿量可以表示为一个具有格点无关张量\footnote{在此记号中，矩阵元 $M_{\alpha\beta}$ 是作用在物理希尔伯特空间上的算符。} $M$ 且键维为 3 的 MPO：
\begin{equation}
    M_{\alpha\beta} = \begin{pmatrix}
    \mbb 1 & -\sigma^Z & -\lambda\sigma ^X\\
    0 & 0 & \sigma^Z\\
    0 & 0 & \mbb 1
    \end{pmatrix}_{\alpha\beta},
\end{equation}
以及边界向量
\begin{equation}
\vec{v}_L=\begin{pmatrix}1\\0\\0\end{pmatrix},\qquad
\vec{v}_R=\begin{pmatrix}0\\0\\1\end{pmatrix}.
\end{equation}

注意到，由于我们在此处理的是非单射 MPO（存在不变子空间），因此必须定义边界向量。
该 MPO 是 Onsager 代数的生成元之一。
其指数形式 $\exp(itH)$ 无法再写成键维与系统尺寸无关的 MPO——除非 $\lambda=0$，此时（见习题）它可以写成一个 $\chi=2$ 的单射 MPO。

\subsection{MPS 流形}\label{sec:MPSmanifold}
既然我们已经通过纠缠对构造的视角引入了矩阵乘积态和算符，现在让我们更详细地研究均匀矩阵乘积态流形~\cite{Haegeman2013a,Haegeman2014b,Vanderstraeten2019}。为此，我们考虑具有固定键维 $\chi$ 的所有均匀 MPS $\ket{\Psi[A]}$。作为初步观察，注意到这些态并不张成整个希尔伯特空间的线性子空间，因为两个 MPS 的和通常不再是具有相同键维的 MPS。相反，它们位于一个曲面上，其局部坐标由该态的变分参数（即 $A^i_{\alpha\beta}$）给出。回顾上文，对于流形上的每一个点，我们都可以关联一个线性切空间。对于当前情形，切向量具有形式
\begin{equation}
    \ket{\Psi[B,A]} = \sum_I B^I\pdv{}{A^I}\ket{\Psi[A]},
\end{equation}
在此处及后文中，$I$ 统指 MPS 张量的物理指标与虚指标。
在图形上，这样一个切向量归结为
\begin{equation}\label{eq:TangentVector}
    \ket{\Psi[B,A]} = \sum_\msf n \inputtikz{TangentVector} \,\, ,
\end{equation}
其中对 $\msf n$ 的求和是对 $B$ 张量在无限链中所有可能位置的求和。

在这一流形上的任意给定点 $\ket{\Psi[A]}$ 处，两个切向量 $\ket{\Psi[B,A]}$ 与 $\ket{\Psi[B',A]}$ 的重叠通过下式给出一个度规或 \emph{Gram 矩阵} $G$：
\begin{equation}
    \braket{\Psi[\bar B',\bar A]}{\Psi[B,A]} = \sum_{I,J} \bar B'^I \underbrace{\braket{\pdv{}{\bar A^I}\Psi[\bar A]}{\pdv{}{A^J}\Psi[A]}{}}_{G(A,\bar A)_{I,J}} B^J.
\end{equation}
天真地看，人们可能会认为流形的维度等于张量 $B$ 中参数的数量，即 $d\chi^2$。但请注意，形式为 $B^i \rightarrow B^i + XA^i - A^iX$ 的 gauge 变换使态 $\ket{\Psi[B,A]}$ 保持不变，因此 $d\chi^2$ 过度估计了 MPS 流形的维度。文献~\cite{Vanderstraeten2019}中给出的更审慎论证表明，这些参数中只有 $(d-1)\chi^2$ 个可以被独立选取。假定 $A_L^i$ 处于所谓的\emph{左 canonical gauge}（参见\cref{sec:NormalObs}），则 $B$ 的忠实参数化形式为 $B^i=V^iX$，其中 $V^i$ 张成 $(A_L)_{\alpha\beta}^i$（解释为从 $\beta$ 到 $\alpha,i$ 的矩阵）的零空间，而 $X$ 编码 $B$ 的独立自由度。通过这种对 $B$ 的参数化，度规在局部是欧几里得的，即 Gram 矩阵等于单位阵。从计算的角度来看，该参数化非常有用，因为它极大地简化了算法的实现并提高了其数值稳定性。

其中一个此类算法便是应用于 MPS 类的 TDVP。回顾\cref{sec:Illusion}，TDVP 旨在通过将方程右侧投影到切平面上来近似求解含时薛定谔方程
\begin{equation}
    i\pdv{t}\ket{\Psi[A]} = H \ket{\Psi[A]}
\end{equation}
即：
\begin{equation}\label{eq:TDVP}
    i\pdv{t}\ket{\Psi[A(t)]} = \mc P_{A(t)}H \ket{\Psi[A(t)]},
\end{equation}
其中 $\mc P_{A(t)}$ 为切平面上的投影算符。对 TDVP 方程 \cref{eq:TDVP} 的进一步运算揭示出：TDVP 等价于 MPS 流形上由 \emph{Poisson 括号}支配的一组（半）经典运动方程。著名的密度矩阵重整化群（DMRG）算法~\cite{White1992}（在\cref{sec:DMRG}中讨论）及其所有变体都可以理解为求解该微分方程的特定分裂方法~\cite{Haegeman2016b}。回想一下，Poisson 括号是相空间上的双线性、反对称函数，且满足 Jacobi 恒等式。若我们将依赖于 $A$ 和 $\bar A$ 的任意函数 $f$ 和 $g$ 的 Poisson 括号定义如下：
\begin{equation}
    \{f,g\} = - \frac{\partial f}{\partial A_I} (G^{-1})^{IJ} \frac{\partial g}{\partial \bar A_J} + \frac{\partial g}{\partial A_I} (G^{-1})^{IJ} \frac{\partial f}{\partial \bar A_J},
\end{equation}
并记 $f(A,\bar A)=\expval{F}{\Psi[A]}$，则可以证明期望值 $f$ 的时间演化由下式决定：
\begin{equation}
    \partial_t f = i\{f,h\}.
\end{equation}
这带来两个重要推论。一方面，由 $\{h,h\}=0$ 可知，能量期望值在时间演化下是严格保持守恒的。这正如上文所预期的，意味著该流形是辛流形。另一方面，对于任何与哈密顿量对易（从而构成对称性生成元）且与切空间投影算符 $\mc P_A$ 对易的算符 $K$，可以证明 $\{k,h\}=0$，因此 $\partial_t k=0$。特别地，对于作为局域在位幺正算符张量积的对称性生成元，情况正是如此。

除 TDVP 之外，还值得一提的是，流形视角还为直接在热力学极限下进行基态优化提供了必要的框架，即所谓的\emph{变分均匀 MPS}（variational uniform MPS, VUMPS）算法，同时也为\emph{准粒子激发}（quasiparticle excitations）的构建奠定了基础。这些主题分别在\cref{sec:DMRG_VUMPS}和\cref{sec:Excitations}中讨论，但首先我们将更深入地探讨 MPS 的演算。
\subsection{第一讲小结}
第一讲的核心信息在于：完整的多体希尔伯特空间只是一种幻象；为了描述诸如量子自旋系统基态等相关的多体态，只需考虑一个捕捉了所有相关物理的低维态流形即可。对于有能隙量子自旋链这一特定情形，量子纠缠的单配性预示着纠缠熵面积律的存在。这意味著 MPS 流形与所有满足面积律的态之间存在一一对应关系。因此，我们不必聚焦于整个希尔伯特空间，而是可以遵循含时变分原理的指引将哈密顿量投影到该 MPS 流形上；所有张量网络算法都可以在这一框架下得到理解。
\newpage\section{第二讲：MPS 演算与算法}
在本节中，我们将介绍矩阵乘积态的演算。具体而言，这包括可观测量的计算，通过为 MPS 张量引入合适的 \emph{normal form}，该计算可得到极大简化。借助这些工具，我们将重点讨论用于有限链基态优化的著名\emph{密度矩阵重整化群}（density matrix renormalisation group, DMRG）算法，以及直接在热力学极限下工作的\emph{变分均匀矩阵乘积态}（variational uniform matrix product state, VUMPS）算法。随后我们研究基态之上的准粒子激发，它们被参数化为\emph{准粒子 ansatz}（quasiparticle ansatz）。首先，让我们从矩阵乘积态的更多示例开始。
\subsection{示例、单射性与 parent Hamiltonian}\label{sec:Examples}
\paragraph{直积态} 任意直积态（product state）$\ket{\Psi}=\ket{\psi}^{\otimes N}$（其中 $\ket{\psi}=\sum_i \psi_i\ket{i}$）都可以写为键维为 1 的 MPS，其形式为
\begin{equation}
    A^i_{11} = \psi_i.
\end{equation}
\paragraph{AKLT} 上文在\cref{sec:AreaLaw}中遇到的非平庸矩阵乘积态的典型范例是 Affleck-Kennedy-Lieb-Tasaki（AKLT）态。该态由如下 MPS 矩阵生成：
\begin{equation}\label{eq:AKLTState}
    A^i = \sqrt{\frac{1}{3}}\sigma^i, \quad i=X,Y,Z,
\end{equation}
其中 $\sigma^i$ 为 Pauli 矩阵，前置系数确保该态在热力学极限下归一化。

在周期性边界条件下，该态是自旋 1 哈密顿量 \eqref{eq:AKLTHam} 的唯一有能隙基态。正如后文在\cref{sec:SPT1d}中所解释的，该哈密顿量在开链上具有四重基态简并，其特征是存在无能隙边缘模。这是其\emph{对称保护拓扑}（symmetry-protected topological, SPT）序的显著标志。
\paragraph{GHZ} 另一个可以表示为键维为 2 的 MPS 的态是所谓的 Greenberger–Horne–Zeilinger（GHZ）态~\cite{Greenberger1989}：
\begin{equation}
    \ket{\Psi} = \frac{1}{\sqrt{2}} \big( \ket{0}^{\otimes N} + \ket{1}^{\otimes N} \big).
\end{equation}
其 MPS 表示为：
\begin{equation}
    A^i =
    \begin{cases}
        \frac{1}{\sqrt{2}}\begin{pmatrix} 1 & 0 \\ 0 & 0\end{pmatrix}, \quad i=0 \\
        \frac{1}{\sqrt{2}}\begin{pmatrix} 0 & 0 \\ 0 & 1\end{pmatrix}, \quad i=1.
    \end{cases}
\end{equation}
注意到 GHZ 态可以被视为表现出 $\mbb Z_2$ 对称性破缺。这反映在 MPS 矩阵是分块对角的，且这些分块在 $\prod_\msf i \sigma^X_\msf i$ 的作用下发生互换。显然，GHZ 态的这一 MPS 表示是非单射的。

这必须与 AKLT 态形成对比——AKLT 态的 MPS 矩阵张成整个 $\mbb C^{2\times 2}$ 矩阵代数，因此不能被化为分块对角形式，从而定义了一个单射 MPS。注意 AKLT 态的单射长度为 2，这可以通过直接计算加以验证。

通常情况下，每个均匀 MPS 都可以化为\footnote{在消除非对角块之后，如文献~\cite{Cirac2020}中所综述。}\emph{分块单射}（block injective）形式或\emph{标准}形式（standard form），使得所有 MPS 矩阵都具有形式~\cite{PerezGarcia2006,Cadarso2012}
\begin{equation}
    A^i = \bigoplus_{k=1}^r \mu_k A_k^i,
\end{equation}
其中所有 $\{A^i_k\}_k$（$k=1,\dots,r$）都是单射的，有时被称为\emph{单射块}（injective blocks）。需要指出的是，对于\emph{费米子} MPS，这种单射性概念需要重新审视，这是\cref{sec:Fermions}的主题。

有趣的是，每个 MPS 都可以被视为相应 \emph{parent Hamiltonian} 的基态，从而将 MPS 及其关联的 parent Hamiltonian 的研究置于同等地位。该 parent Hamiltonian 是有能隙的、局域的且无阻挫的，意味着每一项的能量都被该 MPS 最小化。它的局域项是相应连续自旋区域上局域约化密度矩阵零空间上的投影算符~\cite{PerezGarcia2006,Cirac2020}。例如，对于 GHZ 态，局域 parent Hamiltonian 项由 $\frac{1}{2}(\mbb 1 - \sigma^Z_{\msf i}\sigma^Z_{\msf i+1})$ 给出。当 MPS 为单射时，可以证明它是其 parent Hamiltonian 的唯一基态。更有趣的是，当 parent Hamiltonian 与对称群 $G$ 的在位表示对易时（如在 GHZ 态的情形中），$G$ 的作用会置换不同单射块的（一个极小不可约子集）。在那种情况下，可以找到一个极大子群 $H\subseteq G$，其作用使某一特定参考块保持不变~\cite{Schuch2011}。随后各块由 $G/H$ 中的\emph{陪集}（cosets）标记。\footnote{回想一下，（左）陪集 $gH$ 是集合 $gH=\{gh \mid h\in H\}$。所有陪集的集合记为 $G/H$，构成群 $G$ 的一个划分。}

下文在第~\ref{sec:PhasesOfMatter}讲和第~\ref{sec:GenSym}讲中，我们将再次遇到 parent Hamiltonian，将其作为分类\emph{量子物质相}（quantum phases of matter）的工具；粗略地说，量子物质相是可以通过绝热过程相互演化的哈密顿量的等价类。

\subsection{Normal form、纠缠与可观测量}\label{sec:NormalObs}
如\cref{sec:AreaLaw}所述，矩阵集合 $\{A^i\}$ 并不是给定 MPS 的唯一表示。这是因为对于任意非奇异矩阵 $X$，矩阵乘积态在如下形式的重新参数化下是不变的：
\begin{equation}
    A^i \mapsto XA^iX^{-1}, \quad \forall i=1,2,\dots,d,
\end{equation}
因为在连接相邻 MPS 矩阵的键上，矩阵 $X$ 与 $X^{-1}$ 恰好相互抵消。这种冗余通常称为 MPS 描述的 \emph{gauge} 自由度。事实证明，这种 gauge 自由度可以被充分利用，以简化特定计算并提升 MPS 算法的数值稳定性。

通过适当的 gauge 变换，MPS 矩阵可以化为所谓的\emph{左 canonical form} 或\emph{右 canonical form}，即分别通过如下条件定义：
\begin{align}
    \sum_{i} A^{i \dagger} A^i &= \mbb 1_\chi \label{eq:MPSLeftCanonical}, \\
    \sum_{i} A^i A^{i \dagger} &= \mbb 1_\chi. \label{eq:MPSRightCanonical}
\end{align}
在图形上，这些条件因此可以写为如下形式：
\begin{align}
    \inputtikz{MPSRightCanonical} \quad &= \quad \inputtikz{MPSRightIdentity}\,\, , \\
    \inputtikz{MPSLeftCanonical} \quad &= \quad \inputtikz{MPSLeftIdentity}\,\, .
\end{align}
请注意，条件 \cref{eq:MPSLeftCanonical,eq:MPSRightCanonical} 并未完全固定 MPS 矩阵，因为我们仍可以用\emph{幺正} gauge 矩阵对这些矩阵进行共轭。在后文中，我们将分别用 $A_L$ 和 $A_R$ 表示 MPS 张量 $A$ 的左 canonical form 与右 canonical form。让我们再次强调，$A$、$A_L$ 和 $A_R$ 生成的都是同一个态。将左 canonical form 变换为右 canonical form 的 gauge 变换或 \emph{pivot 矩阵}通常记为 $C$，从而由 $A_L^iC=CA_R^i$（$\forall i$）所刻画。为便于后文引用，我们还引入记号 $A_C^i =A_L^iC=CA_R^i$。若一个 MPS 可以写为如下形式，则称其处于 \emph{mixed canonical form}：
\begin{equation}\label{eq:MPSMixedGauge}
    \ket{\Psi[A]} = \inputtikz{MPSMixedgauge}\,\,.
\end{equation}
正如该图所暗示的，矩阵 $C$ 包含了 $C$ 左侧与右侧子系统之间纠缠的信息。让我们对此作更精确的阐述。

考虑一个允许张量积分解 $\mc H=\mc H_A\otimes \mc H_B$ 的希尔伯特空间，其中 $d_A=\dim(\mc H_A)$，$d_B=\dim(\mc H_B)$。借助\emph{奇异值分解}（singular value decomposition, SVD），$\mc H$ 中的任意态均可写为如下形式：
\begin{equation}\label{eq:SchmidtState}
    \ket{\Psi} = \sum_k s_k \ket{\psi_k^A}\otimes\ket{\psi_k^B},
\end{equation}
其中 $\ket{\psi_k^A}$ 与 $\ket{\psi_k^B}$ 分别是 $\mc H_A$ 与 $\mc H_B$ 的标准正交基向量，实数 $s_k$（$k=1,2,\dots,\min(d_A,d_B)$）称为\emph{奇异值}或 \emph{Schmidt 数}，并构成该态的\emph{纠缠谱}（entanglement spectrum）。通过对任一子系统求偏迹，我们可以将约化密度矩阵写为
\begin{align}
    \rho_A &= {\rm Tr}_B(\ketbra{\Psi}) = \sum_{k=1}^{d_A} s_k^2 \ket{\psi_k^A}\bra{\psi_k^A}, \\
    \rho_B &= {\rm Tr}_A(\ketbra{\Psi}) = \sum_{k=1}^{d_B} s_k^2 \ket{\psi_k^B}\bra{\psi_k^B}.
\end{align}
与 MPS \eqref{eq:MPSMixedGauge} 进行直接对比可以发现，态 $\ket{\Psi[A]}$ 的 Schmidt 谱与 $C$ 的奇异值完全一致。通过残余的幺正 gauge 自由度 $A_L^i\mapsto UA_L^iU^\dagger $，$A_R^i\mapsto V^\dagger A_R^iV$（其中 $U$ 与 $V$ 来自 SVD $C=USV^\dagger$），MPS $\ket{\Psi[A]}$ 最终可以化为式 \eqref{eq:SchmidtState} 的形式。

\bigskip\noindent MPS 的关键优势之一在于能够高效计算局域可观测量的期望值及其关联函数。这里的“\emph{高效}”是指计算代价随局域算符之间的距离呈线性增长。作为说明，让我们关注分别位于格点 $\msf i$ 与 $\msf j$ 处的局域可观测量 $O_1$ 与 $O_2$ 的两点关联函数。由于平移不变性，该期望值仅依赖于距离 $L=\vert\msf j - \msf i\vert$。按照惯例定义可观测量的期望值：
\begin{equation}
    \expval{O_1[\msf i]O_2[\msf j]}_{\Psi[A]} := \mel{\Psi[A]}{O_1[\msf i]O_2[\msf j]}{\Psi[A]},
\end{equation}
这归结为如下收缩：
\begin{equation}
    \expval{O_1[\msf i]O_2[\msf j]}_{\Psi[A]} \, = \inputtikz{MPSCorrelator_1}.
\end{equation}
假设单射性成立，从而转移矩阵具有唯一的最大本征值 1 以及相应的左右本征向量 $\Lambda$ 与 $\rho$，我们能够将该收缩重写为
\begin{equation}
    \inputtikz{MPSCorrelator_2}\,\, ,
\end{equation}
这表明该关联函数的计算代价按 $\mc O(L\chi^3)$ 标度。事实上，该张量网络的收缩是从左到右进行的，通过明智地选择收缩顺序，主要计算代价在于 $\chi\times\chi$ 矩阵的乘法（参见文献~\cite{Pfeifer2014}中确定一般张量网络最优收缩顺序的算法）。
作为该表达式的特例，取 $O_1\equiv O$ 且 $O_2\equiv\mbb 1$，我们得到局域算符的期望值：
\begin{equation}
    \inputtikz{MPSExpVal}\,\, ,
\end{equation}
以及 MPS 的模平方：
\begin{equation}
    \braket{\Psi[A]} := \inputtikz{MPSNorm}\,\, .
\end{equation}
若我们按惯例定义\emph{连通}（connected）两点关联函数：
\begin{equation}
    C(O_1[\msf i],O_2[\msf j]) := \expval{O_1[\msf i]O_2[\msf j]}_{\Psi[A]} - \expval{O_1[\msf i]}_{\Psi[A]} \expval{O_2[\msf j]}_{\Psi[A]},
\end{equation}
再次利用转移矩阵的谱分解 \cref{eq:TransferSpec}（假定其存在），可以证明对该关联函数起主导作用的贡献随 $L$ 呈\emph{指数}衰减，更精确地说：
\begin{equation}
    C(O_1[\msf i],O_2[\msf j]) \sim \lambda_2^{L-1} \sim \exp(-L/\xi),
\end{equation}
这里我们将\emph{关联长度}（correlation length）$\xi$ 定义为 $\xi = - 1/\log(\vert\lambda_2\vert)$。关联函数的这种指数衰减印证了 MPS 能够很好地描述\emph{有能隙}哈密顿量的基态物理，因此无法直接重现\emph{临界系统}中存在的关联代数衰减。然而，借助下一小节的主题——\emph{纠缠标度}（entanglement scaling）理论，依然可以实现对临界关联函数的精确近似。
\subsection{标度理论}\label{sec:scaling}
除有能隙基态物理之外，MPS 方法也是计算\emph{临界}系统性质的有力工具。利用直接在热力学极限（即无限长自旋链，参见第 \ref{sec:DMRG_VUMPS} 小节）下工作的均匀方法，可以提取描述连续相变普适类或共形场论（CFT）的临界指数与中心荷。关于该主题已有大量现有研究，最初由文献~\cite{Nishino1996,Tagliacozzo2007,Pollmann2009,Pirvu2012}引入。在此我们遵循文献~\cite{Zauner2015,Rams2018,Vanhecke2019}中发展的方法。

在二阶临界点附近，我们预期关联函数具有 Ornstein-Zernike 形式，即幂律项与指数衰减项的乘积。这可以从关联函数的 K{\"a}ll\`en-Lehmann 表示理解为连续指数项的线性叠加~\cite{Zauner2015}。若要使用 MPS 直接表示该连续谱，需要无穷大键维的 MPS。使用有限键维的 MPS 等价于对该连续谱进行离散化。我们可以通过转移矩阵谱中的能隙来量化这种“离散性”或近似程度。设 $\lambda_i$ 如前所述为转移矩阵的本征值，其中 $\lambda_1=1$，并定义 $\epsilon_i = - \ln\lambda_i$。取前 $n$ 个谱变量 $\epsilon_i$（排除第一个）的线性组合：
\begin{equation}
    \delta := \sum_{i=2}^{n} c_i \epsilon_i  \quad \text{with} \quad \sum_{i=2}^n c_i = 0.
\end{equation}
在文献~\cite{Zauner2015,Rams2018}中证明了 $\epsilon_i$ 具有动量的物理意义，因此相邻 $\epsilon_i$ 值之间的差引入了一个反长度标度 $\delta$；无论在有能隙还是无能隙情形下，当 $\chi\rightarrow\infty$ 时该标度均趋于零。类比于 Cardy 的有限尺寸标度假说~\cite{Cardy}（其中长度标度为系统尺寸），可以利用这一反长度标度 $\delta$ 为二阶相变的任意序参量 $m$ 建立如下\textit{标度假说}（scaling hypothesis）：
\begin{equation}
    m(t,\delta) = s^{-\beta/\nu} \, m\!\left( s^{1/\nu} t,\, s\delta \right).
\end{equation}
关联长度的相应标度假说为：
\begin{equation}
    \delta \xi(t,\delta) = \xi\!\left( \delta^{-1/\nu} t,\, 1 \right)
= \tilde{\xi}\!\left( \delta^{-1/\nu} t \right).
\end{equation}
由上述形式，还可以推导出纠缠熵的标度假说。利用 CFT 中纠缠熵 $S$ 的形式~\cite{Calabrese2004}，可知 $\text{exp}(\frac{6}{c}S)$ 具有长度的标度行为，因此 $S$ 的标度假说为：
\begin{equation}
    \exp\!\left( \frac{6}{c} S(t,\delta) \right)
= s \exp\!\left( \frac{6}{c} S\!\left( s^{-1/\nu} t,\, s\delta \right) \right),
\end{equation}
其中 $c$ 为中心荷。利用具有不同键维 $\chi$ 的均匀 MPS 方法进行模拟，便可通过标准优化方法拟合 $m$、$\xi$ 和 $c$。这一切的核心在于：临界点附近哈密顿量多种不同参数以及不同键维下的许多不同模拟结果可以相互关联，并且临界点与临界指数的正确数值将导致所有相应曲线发生数据塌缩。细节请参考文献~\cite{Vanhecke2019}。使用有限键维就如同在系统中打开一个能隙，其效果与在临界理论中添加一个相关微扰无法区分。通过这种方式，有限纠缠标度理论提供了模拟临界理论的最先进方法——MPS 始终具有能隙这一事实并非缺陷，而是一种特性！
\subsection{变分 MPS 算法}
\label{sec:DMRG_VUMPS}
\subsubsection{密度矩阵重整化群}\label{sec:DMRG}
\emph{密度矩阵重整化群}（density matrix renormalisation group, DMRG）算法最初由 Steven White 提出~\cite{White1992}，是一种用于寻找\emph{局域一维哈密顿量}基态的变分方法。它已被极其成功地应用于大量模型中，至今仍是 $(1+1)$ 维强关联自旋链基态优化的最先进方法~\cite{Schollwoeck2010,Verstraete2023}。

尽管 DMRG 最初是基于重整化群推导和解释的，但其现代诠释是定义在矩阵乘积态类上的变分方法，其中通过交替最小二乘法迭代寻找最优解~\cite{Verstraete2004a}。DMRG 在最大（有界）键维为 $\chi_{\rm max}$ 的所有有限 MPS 集合上最小化期望值
\begin{equation}\label{eq:min}
    \underset{\{A\}}{\rm min}\left(\frac{\expval{H}{\Psi[A]}}{\braket{\Psi[A]}}\right).
\end{equation}
所选取的 $\chi_{\rm max}$ 越大，近似效果越好。即使在计算资源适度的情况下，通过利用如 Lanczos 方法等稀疏本征求解器求解有效本征值问题（见下文），也可以轻松达到 $\chi\sim\mc O(10^3)$ 的键维；在同时利用对称性的情况下，甚至可达到 $\mc O(10^5)$~\cite{Sanz2009,Singh2009,Weichselbaum2012,Singh2013,Lootens2024,Devos2025}。

假定哈密顿量写为 MPO 形式（见\cref{sec:MPO}），式 \cref{eq:min} 归结为：
\begin{equation}
    \underset{\{A\}}{\rm min}\,\, \frac{\inputtikz{DMRG_H}}{\inputtikz{DMRG_norm}}\,\,.
\end{equation}
该算法通过提供基态的初始猜测（通常只需一个随机 MPS）启动，并对所有张量进行一系列局域顺序优化，在自旋链上来回扫描，直至达到收敛。

固定某个格点 $\msf i$，将其左侧的所有张量化为左 canonical form，右侧的所有张量化为右 canonical form，此时对张量 $A[\msf i]$ 的更新归结为将其替换为\emph{有效哈密顿量}（effective Hamiltonian）最小实本征值所对应的本征向量：
\begin{equation}
    \inputtikz{DMRG_Heff1}\,\,.
\end{equation}
在此图中，$L$ 与 $R$ 分别代表通过收缩格点 $\msf i$ 左侧与右侧的所有 MPS 和 MPO 张量所获得的（依赖于格点的）\emph{环境}（environments）。

由此，该版本的算法逐个单独更新每个格点，被称为\emph{单格点 DMRG}（one-site DMRG）。该方法的缺点是在扫描过程中没有便利的方法来增加键维。当施加对称性时，这一问题尤为严重，因为在这种情况下，算法无法在虚指标能级上探索不同的对称荷。算法的\emph{双格点}版本弥补了这一缺陷。

双格点版本的算法在每次扫描的每一步中将两个格点组合在一起，随后将组合张量 $A'[\msf i,\msf i+1]$ 更新为有效双格点哈密顿量的极小本征向量：
\begin{equation}
    \inputtikz{DMRG_Heff2}\,\, .
\end{equation}
随后利用奇异值分解将 $A'$ 分解为 $A[\msf i],A[\msf i+1]$。正是在这一步中，可以通过截断 SVD（仅保留最大的奇异值）来动态控制键维。DMRG 的高层伪代码如下所示。

\begin{algorithm}[H]
\caption{双格点 DMRG（Two-site DMRG）}
 初始化 MPS 张量 $A[\msf i]$（通常初始化为随机态）\\
\For{$k = 1$ \text{到} \texttt{num\_sweeps}}{
    \tcp{从左至右扫描}
    \For{$\msf i = 1$ \text{到} $N-2$}{
         优化并更新格点对 $\msf i, \msf i+1$ \\
         更新格点 $\msf i+1$ 处的左环境 \\
    }

    \tcp{从右至左扫描}
    \For{$\msf i = N-1$ \text{到} $2$}{
         优化并更新格点对 $\msf i, \msf i+1$ \\
         更新格点 $\msf i+1$ 处的右环境 \\
    }
}
\Return $\{A[\msf i]\}$
\end{algorithm}
\subsubsection{变分均匀矩阵乘积态}\label{sec:VUMPS}
\emph{变分均匀矩阵乘积态}（variational uniform matrix product state, VUMPS）算法是 DMRG 在无限链上的等价方法，由文献~\cite{Zauner-Stauber2018a}引入（教学性综述参见文献~\cite{Vanderstraeten2019}，另见文献~\cite{McCulloch2008}中的替代方法）。它直接在热力学极限下工作，这使得它对于模拟临界系统以及提取其标度律特别有用。当然，在无限极限下，总基态能量是无界的，因此 VUMPS 最小化的是\emph{能量密度}（energy density）。

在介绍该算法之前，我们先固定一些记号。回顾\cref{sec:NormalObs}中引入的中心 gauge MPS $A_C$ 和键张量 $C$，以及左、右 gauge $A_L, A_R$。随后我们将有效单格点哈密顿量 $H_{\rm eff}$ 以及键张量 $C$ 的有效环境 $H_C$ 定义为：
\begin{equation}
\label{eq:VUMPSEnvironments}
    H_{\text{eff}} := \inputtikz{VUMPS_EffectiveHamiltonian}\,\,, \qquad
    H_C := \inputtikz{VUMPS_EffectiveEnvironment}\,\,.
\end{equation}
此处，左环境与右环境 $L,R$ 分别是转移矩阵 $\mbb T_L^{[H]}, \mbb T_R^{[H]}$ 的不动点：
\begin{equation}
    \mbb T_L^{[H]} := \inputtikz{VUMPS_TransferMatrixHamiltonianLeft}\,\, , \qquad
    \mbb T_R^{[H]} := \inputtikz{VUMPS_TransferMatrixHamiltonianRight}\,\,.
\end{equation}
它们可以通过求解线性方程组来计算。

如文献~\cite{Zauner-Stauber2018a}中所推导，在均匀 MPS 流形中基态能量密度的变分极小值由如下方程组的同时解所刻画：
\begin{equation}
\begin{split}
    &H_{\rm eff}(A_C) \propto A_C ,\\
    &H_C(C) \propto C, \\
    &A^i_C = A_L^iC = CA_R^i.
\end{split}
\end{equation}
包含最后一个方程是为了避免必须取 $A_L=A_CC^{-1}$ 和 $A_R=C^{-1}A_C$——若 $C$ 接近奇异，这在潜在意义上是病态的。给定基态的初始猜测，VUMPS 算法迭代逼近该解直至达到收敛。在每一步中，$A_C$ 与 $C$ 的当前值通过将其替换为当前 $H_{\rm eff}$ 与 $H_C$ 具有最小实本征值的本征向量来更新。随后，通过对 $A_CC$ 和 $CA_C$ 采用奇异值分解，由 $\tilde A_C$ 与 $\tilde C$ 求得近似的 $\tilde A_L$ 与 $\tilde A_R$：
\begin{equation}\label{eq:approxALAR}
\begin{split}
    \tilde A_L &= U_LV_L^\dagger, \quad U_L\Sigma_LV_L^\dagger =\tilde A_C\tilde C,\\
    \tilde A_R &= U_RV_R^\dagger,  \quad \tilde C^\dagger \tilde A_C = U_R\Sigma_R V_R^\dagger.
\end{split}
\end{equation}
如此一来，2-范数
\begin{equation}
\begin{split}
    \epsilon_L &= \Vert\tilde A_C - \tilde A_L\tilde C\Vert_2, \\
    \epsilon_R &= \Vert\tilde A_C - \tilde C\tilde A_R\Vert_2, \\
\end{split}
\end{equation}
被最小化。

用于 MPO 哈密顿量的 VUMPS 伪代码如下所示。

\begin{algorithm}[H]
\caption{VUMPS}
    初始化 $\{A_L,A_R,C\}$（通常初始化为随机态）\\
    初始化当前精度 $\epsilon_{\rm cur}$\\
    \While{$\epsilon_{\rm cur}>\epsilon_{\rm tol}$}{
        由 $A_L, H$ 计算 $\mbb T_L^{[H]}$ \\
        由 $A_R, H$ 计算 $\mbb T_R^{[H]}$ \\
        由 $\mbb T_R^{[H]}$ 计算左环境 $L$ \\
        由 $\mbb T_L^{[H]}$ 计算右环境 $R$ \\
        由 $L, R, H$ 构造有效无限哈密顿量 $H_{\text{eff}}$ \\
        由 $L, R$ 构造有效无限环境 $H_C$ \\
        求解 $H_{\text{eff}}$ 的主导本征向量 $A_{C}$ \\
        求解 $E_{\text{eff}}$ 的主导本征向量 $C$ \\
        如 \cref{eq:approxALAR} 所示，由 $A_{C}, C$ 计算 $A_L, A_R$ \\
        设定 $\epsilon_{\rm cur}=\max(\epsilon_L,\epsilon_R)$
    }
    \Return $\{A_L,A_R,C\}$
\end{algorithm}
\subsection{准粒子激发}\label{sec:Excitations}
截至目前，我们主要关注基态以及在有限系统（\cref{sec:DMRG}）和无限系统（\cref{sec:VUMPS}）设定下对其进行变分近似的算法。对量子多体系统的全面理解还应当解释其低能激发。通常，有能隙量子多体哈密顿量的能谱由若干孤立的分支（具有局域化\emph{准粒子激发}（quasiparticle excitations）的诠释）以及一系列连续\emph{散射态}（scattering states）（本质上可以通过叠加准粒子态来构建）构成。这些准粒子表现在动力学关联函数中（可以通过例如非弹性中子散射实验进行探测），但在临界系统的情形下，还编码了其有效连续体描述的普适性质\footnote{请记住，在使用 MPS 近似临界基态时，有限键维会使激发产生能隙。}。

矩阵乘积态技术可借助所谓的\emph{准粒子 ansatz}（quasiparticle ansatz）来捕捉这些准粒子激发及其散射态，该方法最早由文献~\cite{Ostlund1995}引入，随后在文献~\cite{Haegeman2011b}中推广至一维链中的单粒子激发，并在文献~\cite{Vanderstraeten2013}中推广至两粒子态。在此，我们假设一个平移不变的局域有能隙哈密顿量 $H$。暂时我们还假设存在唯一基态以排除畴壁激发，后者将在下文讨论。我们将基态的 MPS 近似记为 $\ket{\Psi[A]}$。具有动量 $k$ 的准粒子的 ansatz 随后具有形式
\begin{equation}\label{eq:QP}
    \ket{\Psi[A,B],k}= \sum_\msf n e^{ik\msf n} \inputtikz{TangentVector}\,\,.
\end{equation}
也就是说，该 ansatz 将其中一个基态张量替换为包含该 ansatz 全部变分参数的张量 $B$，并在无限链中对 $B$ 的位置进行动量叠加。因此，它可以被视作切向量 \cref{eq:TangentVector} 的 \emph{boosted} 版本，从而使人们能够充分利用为切空间发展起来的整套完备技术。此外，值得一提的是，准粒子 ansatz 可以被诠释为 Feynman-Bijl ansatz 或为描述液氦中集体激发而发展的\emph{单模近似}（single mode approximation）的推广~\cite{Bijl1941,Feynman1953,Feynman1954}。请注意，由于有限键维的存在，我们不应将 $\ket{\Psi[A,B],k}$ 理解为描述位于张量 $B$ 处的局域粒子，而应理解为相对于均匀真空具有增加能量密度的团块。给定 $A$，对 $B$ 的优化归结为有效哈密顿量的广义本征值问题，其精神与\cref{sec:VUMPS}类似。该算法的详细阐述超出了本讲义的范围，但我们要指出的是，在其中使用与\cref{sec:MPSmanifold}所描述非常相似的对 $B$ 的高效参数化是至关重要的。

值得强调的是，基态编码了关于激发态的信息。事实上，对于一般的自伴算符，其单个本征向量通常不会揭示关于其他本征向量的任何信息。然而，哈密顿量的局域性是确保 \cref{eq:QP} 为准粒子激发提供\emph{良好} ansatz 的决定性特征。确实，在著名的 Lieb-Robinson 界~\cite{Lieb1961,Nachtergaele2006,Hastings2010}的应用中，文献~\cite{Haegeman2013b}证明了：通过用（空间上）局域算符的动量 $k$ 叠加作用于基态，可以任意好地近似每一个（精确）本征态 $\ket{\Psi_{k,\alpha}}$。因此，\cref{eq:QP} 应当被视为由支撑集截断为单个格点的上述算符作用于 MPS 基态上的极致情形。

有趣的是，在存在内部对称性的情况下，激发被组织成在对称群的（不可约）表示下变换的简并多重态。这与我们目前展示的准粒子方法并不兼容，因为在那时，围绕激发能量 $\Delta$ 的能量窗口 $[\Delta-\delta,\Delta+\delta]$ 对任意 $\delta$ 值都包含不止一个激发。为了补救这一点，可以通过用多重态算符 $\widetilde O_m$ 替代算符 $\widetilde O$ 来修改该 ansatz，这在 MPS 图像中归结为指定 $B$ 张量的荷~\cite{Zauner-Stauber2018b}。

在对称性破缺从而导致简并基态子空间的情形下，具有最低能量的激发并不具有 \cref{eq:QP} 的形式。相反，它们是构成不同基态之间界面的\emph{拓扑畴壁激发}（topological domain wall excitations）。在此语境下，“\emph{拓扑}”是指这些激发是通过一个广延算符从基态创生的，即作用于终止在畴壁位置的半无限链上的对称性算符。这可以通过修改 ansatz 在 $B$ 张量的两侧分别编码不同的基态张量来处理。在这种情况下，所有的 MPS 计算技术都同样有效，但需要注意将正确的环境相互收缩。

在\cref{fig:HeisenbergExcitations}中，我们展示了自旋 1 反铁磁 Heisenberg 模型的低能激发谱。这是通过显式施加 ${\rm SU}(2)$ 对称性获得的，并且我们用激发的自旋对其进行了标记。孤立的 $S=1$ 分支清晰可见，并在 $k=\pi$ 处达到其极小值 $\Delta\approx 0.41$，这便是著名的 \emph{Haldane 能隙}（Haldane gap）~\cite{Haldane1982,Haldane1983}。令人惊叹的是，利用准粒子 ansatz，即使仅凭借极其适度的计算资源，也能以约 10 位有效数字的精度计算该能隙。
\begin{figure}[htb]
    \centering
    \includegraphics[width=0.7\linewidth]{plots/HeisenbergExcitationSpectrum.pdf}
    \caption{利用键维为 48 的单格点准粒子 ansatz 获得的 Heisenberg 自旋 1 反铁磁体激发谱。激发按其 ${\rm SU}(2)$ 自旋标记。特别注意 \emph{Haldane} 能隙位于动量 $\pi$ 处。使用 MPSKit.jl 制作~\cite{VanDamme}。}
    \label{fig:HeisenbergExcitations}
\end{figure}
\subsection{第二讲小结}
第二讲探讨了矩阵乘积态的计算层面——如何定义 normal form，如何利用交替最小二乘法在 MPS 流形上最小化能量（DMRG），如何将这些方法推广至无限系统尺寸（VUMPS），以及如何通过将完整哈密顿量投影到 MPS 流形的切平面上来定义激发。此外，我们还触及了纠缠标度的主题——借助该理论，可以以前所未有的精度研究临界系统。

在所有这些应用中，计算代价随系统尺寸和键维均仅呈多项式标度这一事实，清楚地印证了张量网络方法切实有效地打破了指数级多体墙——从而具备解决物理与化学中核心问题的巨大潜力。
\newpage\section{第三讲：物态}
\label{sec:PhasesOfMatter}
\subsection{有能隙物态、纠缠与对称性}
量子多体物理中的一个基本挑战是确定给定模型的\emph{相图}（phase diagram）。如果两个模型之间存在一条由有能隙、有界强度、局域 Hamiltonian 构成的光滑插值路径，则称它们属于同一个相。从务实的角度来看，这等价于存在一个将它们的基态相互映射的\emph{次线性深度量子线路}（sublinear-depth quantum circuit）~\cite{Bravyi2006,Hastings2010}。不同的相之间由相变所分隔；在二阶相变的情形下——正如上文在\cref{sec:scaling}中所遇到的——通常由共形场论（conformal field theory, CFT）来描述~\cite{DiFrancesco1997}。对于具有两个参数 $h_{1,2}$ 的模型，虚构的相图示意如下。

相变的 \emph{Landau 范式}（Landau paradigm）指出，每一种物态都与一种对称性破缺模式一一对应~\cite{Landau1937}。最著名（经典）的例子自然是水结成冰的相变，在此过程中连续欧几里得对称性破缺为其 230 个子群之一，这些子群也被称为\emph{空间群}（space groups）。此外，该范式指出，这种对称性的破缺可以通过在对称群下带荷的局域序参量（order parameter）来诊断。反之，正如 Kadanoff 和 Ceva 所论证的，\emph{对称}相可以通过\emph{乱序算符}（disorder operators）来识别~\cite{Kadanoff1970}。自诞生以来，Landau 范式已在诸多方面得到了推广。

首先，人们逐渐认识到，量子系统在将其对称性破缺至某一子群后，其剩余对称性可以在其纠缠自由度上以不同的方式实现。如果我们要求 Hamiltonian 之间的绝热路径与对称性对易，那么这些可区分的纠缠模式就对应于不同的相。出于这个原因，此类\emph{对称保护拓扑}（symmetry-protected topological, SPT）相可以自然地借助张量网络来描述和分类~\cite{Chen2010,Chen2011,Schuch2011,GarreRubio2022}。

本讲的目标是利用 MPS 框架，概述在存在群对称性的一维量子物态的分类。正如在\cref{sec:NormalObs}中所指出的，MPS 中的对称性破缺表现为基态 MPS 中存在不同的单射块（injective blocks）。这些块由子群 $H\subseteq G$ 的陪集 $G/H$ 标记，并在群作用下相互置换。因此，挑战在于对实现在每个此类单射块上的\emph{对称相}进行分类。为此，我们将首先推导 MPS 的\emph{基本定理}（fundamental theorem），该定理指出：一个 MPS 在某一对称性下不变，当且仅当其局域张量在该对称性下平庸变换。这一定理生动地体现了张量网络将态的\emph{全局}性质编码在\emph{局域}张量之中的核心思想。这将引导我们进入\emph{射影表示}（projective representations）的领域，其对应的上同调分类提供了一个\emph{拓扑指数}（topological index），用于对不同的对称相进行分类，即两个单射 MPS 属于同一个相，当且仅当它们的拓扑指数相同。值得注意的是，如果不对系统施加任何对称性，则所有一维玻色系统均属于同一个相。对于费米子而言，情况更为复杂，由于始终存在的费米子宇称对称性，可能存在非平庸的相；这将是\cref{sec:Fermions}的主题。

我们在此也预示后续各讲的目标。随着分数量子霍尔效应及其他二维\emph{拓扑物态}（topological phases of matter）的发现，人们清楚地认识到，即使系统不表现出任何对称性，它仍然可以属于非平庸的量子物态。这些物态的特征在于：依赖于底空间拓扑的基态简并度、有能隙的\emph{任意子激发}（anyonic excitations），以及在某些情况下的手征边缘模（chiral edge modes）。它们同样不能被任何局域序参量所探测。正如我们将在下一讲中所讨论的，这些相的稳定性源于系统中的纠缠是\emph{长程}（long range）的。特别地，在\emph{非手征拓扑序}（non-chiral topological order）的情形下，纠缠模式也是基态简并度的根源所在，因为人们已经认识到，纠缠自由度表现出由\emph{融合范畴}（fusion category）所描述的弦状“对称性”。我们将在下文的\cref{sec:QuantumDoubles}中遇到一类特别简单的拓扑序，并在\cref{sec:SN}中给出一般非手征情形的微观构造。纠缠自由度上这些\emph{广义对称性}（generalised symmetries）的存在尤其意味着，描述边缘模动力学的拓扑序边界 Hamiltonian 也会继承这些对称性。因此，我们自然倾向于为这些广义对称性也考虑一个 Landau 范式，这也将是最后一讲的主题之一。
\begin{equation}
    \inputtikz{PhaseDiagram}
\end{equation}
\subsection{MPS 基本定理}\label{sec:FundamentalTheorem}
正如前文所预期的，MPS 的\emph{基本定理}（fundamental theorem）\cite{Perez2008}在分类一维对称有能隙物态以及 MPS 的诸多其他应用中起着至关重要的作用。我们将应用该定理的如下具体形式：两个均匀\emph{单射}（injective）矩阵乘积态在任意具有周期性边界条件的系统上（与系统长度无关）表示相同的态，当且仅当这两个 MPS 的张量在虚指标层面上通过一个非奇异 gauge 变换相关联：
\begin{equation}\label{eq:FundTh}
\boxed{
    \ket{\Psi[A]} \sim \ket{\Psi[B]} \iff A^i = e^{i\chi} XB^iX^{-1}, \quad\forall i
    }.
\end{equation}
注意，我们在不失一般性的前提下假设两个 MPS 具有相同的键维。当两个 MPS 都处于 canonical form 时，$X$ 是幺正的。

作为前面几讲中所引入的若干关键概念的应用，我们在此给出一个证明。

首先注意到，$\impliedby$ 方向是平凡的，因为 $X$ 与 $X^{-1}$ 的所有出现均在虚指标层面上相互抵消。为了证明 $\implies$ 方向，我们在不失一般性的前提下假设 $\ket{\Psi[A]}$ 处于右 canonical form，而 $\ket{\Psi[B]}$ 处于左 canonical form。现在考虑如下混合转移矩阵（mixed transfer matrix）：
\begin{equation}
    \mbb T_{A,B} = \inputtikz{MPSMixedTransfer} \,\, .
\end{equation}
由于 $\vert\braket{\Psi[A]}{\Psi[B]}\vert=1$，混合转移矩阵的主本征值具有单位模长~\cite{Cirac2016}。我们将其对应的右本征矢量记为 $X$，从而满足 $\mbb T_{A,B}X = e^{i\phi} X$。现在考虑如下收缩：
\begin{equation}
    e^{-i\phi}\,\inputtikz{FundTh_1}\,.
\end{equation}
首先注意，一方面，借助 $\mbb T_{A,B}X = e^{i\phi} X$，该表达式等于 $\Tr(X^\dagger X)$。另一方面，将该表达式解释为关于对角切口的两个矢量的内积，并应用 Cauchy-Schwarz 不等式，我们得到
\begin{equation}
    \Tr(X^\dagger X) = e^{-i\phi} \inputtikz{FundTh_1} \leq \left\lvert\, \inputtikz{FundTh_2} \,\right\rvert^{1/2} \times \hspace{10pt} \left\lvert\,\inputtikz{FundTh_3} \,\right\rvert^{1/2} = \Tr(X^\dagger X),
\end{equation}
其中我们利用了 $A$ 与 $B$ 分别处于右和左 canonical form 的事实。因此等号显然成立，这蕴含着 $A^iX = e^{i\chi}XB^i$；又因 $X$ 是满秩的，由此即得\cref{eq:FundTh}。
\subsection{一维玻色物质的 SPT 分类}\label{sec:SPT1d}
让我们考虑一个具有唯一基态的有能隙 Hamiltonian $H$，并假设它与一个对称性 $G$ 对易：对于表示 $\{U(g),g\in G\}$，满足 $[U(g)^{\otimes L},H]=0$。由于 $H$ 是有能隙的且具有唯一基态，基态必须继承该对称性。这保证了该基态可由一个单射 MPS 忠实表示。因此，该 MPS $\ket{\Psi[A]}$ 在 $G$ 作用下必须平庸变换，即：
\begin{equation}
    U(g)^{\otimes N} \ket{\Psi[A]} \simeq \ket{\Psi[A]}, \quad \forall g \in G.
\end{equation}
直接应用基本定理\cref{eq:FundTh}表明，MPS 张量本身在对称群下必须平庸变换，即：
\begin{equation}\label{eq:gMPS}
    \forall g \in G,\, \exists X_g,\chi_g: \, \sum_j U(g)_{ij} A^j = e^{i\chi_g} X_gA^iX_g^\dagger
\end{equation}
在此应当注意，gauge 矩阵和相位本身均依赖于群元素 $g$。现在让我们考虑由 $g$ 和 $h$ 标记的对称性算符相继作用在该态上的情形。我们可以通过两种不同方式来实现这一点：要么直接作用 $U(gh)$ 并应用一次基本定理，要么依次作用 $U(g)$ 和 $U(h)$ 并应用两次基本定理。令这两个结果相等可得：
\begin{equation}
    e^{i\chi_{gh}}X_{gh} A^i X_{gh}^\dagger = e^{i(\chi_g+\chi_h)}X_gX_h A^i X_h^\dagger X_g^\dagger,\quad\forall i,g,h.
\end{equation}
由于 $A^i$ 是单射的，从而张满整个矩阵代数，这仅在 $\chi_{gh}=\chi_g +\allowbreak \chi_h$ 且 $ X_gX_h = \psi(g,h) X_{gh}$（其中 $\psi(g,h)$ 为 $\rU(1)$ 相位）时才可能成立。前者意味着 $\{e^{i\chi_g},g\in G\}$ 构成了 $G$ 的一维表示。注意，在将有限数目的 $M$ 个格点合并成块（blocking）之后，我们获得了一个放大的原胞，对于该原胞，相位 $e^{iM\chi_g}=1$ 消失。因此我们可以说这一相位\emph{在分块下是不稳定的}，在后续相分类中我们将不再对其做进一步考虑。另一方面，满足如下关系式的矩阵集合 $\{X_g,g\in G\}$：
\begin{equation}\label{eq:ProjRep}
    X_gX_h = \psi(g,h) X_{gh}, \quad \forall g,h\in G,
\end{equation}
（对于某些相位 $\psi(g,h)\in\rU(1)$）构成了 $G$ 的一个所谓\emph{$\psi$-射影表示}（$\psi$-projective representation）。射影表示通过在矩阵乘法中允许射影相位 $\psi(g,h)$，推广了普通的\emph{线性表示}（linear representations）。由于矩阵乘法的结合律，相位 $\psi(g,h)$ 必须满足
\begin{equation}\label{eq:2cocycle}
    \psi(g,h)\psi(gh,k) = \psi(g,hk)\psi(h,k), \quad \forall g,h,k\in G,
\end{equation}
这一自洽性方程通常称为\emph{二重上循环方程}（second cocycle equation）。对于任意函数 $\phi:G\rightarrow\rU(1)$，该条件具有形如 $\psi(g,h)=\frac{\phi_{gh}}{\phi_g \phi_h}$ 的平庸解。在群上同调的语境下，这种平庸的上循环称为\emph{上边缘}（coboundary）。为了对相进行分类，我们将希望考虑上循环方程\emph{模去}上边缘的解。事实上，从基本定理可以清楚看出，对于任意函数 $\phi$，我们都可以自由地重新定义 $X_g\mapsto \phi_gX_g$。为此，注意到上循环与上边缘分别构成阿贝尔群，分别记为 $Z^2(G,\rU(1))$ 与 $B^2(G,\rU(1))$。由于 $B^2(G,\rU(1))$ 显然是 $Z^2(G,\rU(1))$ 的正规子群，我们可以考虑商群 $H^2(G,\rU(1)):=Z^2(G,\rU(1))/B^2(G,\rU(1))$，即 $G$ 在 $\rU(1)$ 上的\emph{第二上同调群}（second cohomology group）。

注意，在上循环方程\cref{eq:2cocycle}两边取对数并将未知相位 $\psi(g,h)$ 收集到矢量 $\bm{\psi}$ 中后，它可以写为一个模 $2\pi$ 满足的线性方程组：
\begin{equation}
    \Omega \bm{\psi} = \vec{0}, \mod 2\pi.
\end{equation}
计算第二上同调群及相应的代表元上循环，可以通过将 $\Omega$ 化为 \emph{Smith normal form} 来完成，这是具有整数系数的矩阵的一种矩阵分解，形式为 $\Omega=P\Lambda R$，其中 $P$ 和 $R$ 均为幺模整数矩阵，而 $\Lambda$ 为对角矩阵。详细解释参见\cref{sec:GroupCohomology}。主要观察是，第二上同调群总是有限群，且具有如下形式：
\begin{equation}\label{eq:H2}
    H^2(G,\rU(1)) \simeq \mbb Z_{d_1}\times\mbb Z_{d_2}\times\dots\times\mbb Z_{d_r},
\end{equation}
其中 $d_i$ 是 $\Lambda$ 的对角元。$H^2(G,\rU(1))$ 的群结构在物理上是通过\emph{堆叠}（stacking）两条 $G$-对称自旋链并将 $G\times G$ 的对角子群作为整体的 $G$ 对称性来实现的。

人们可以证明，如果两个 MPS 在某一 $G$ 对称性下以对应于 $H^2(G,\rU(1))$ 中不同上同调类的射影表示变换，则这两个 MPS 必定表示处于不同相中的基态~\cite{Schuch2011}。换言之，不存在一条可以连接二者而不穿过相变的对称绝热路径。在那类情形中，不同的上同调类有时被称为构成了实现这一转变的\emph{拓扑阻碍}（topological obstruction），而代表元上循环有时被称为表征该相的\emph{拓扑指数}（topological index）。如果我们允许破缺对称性的路径，那么所有 MPS 都将再次属于同一个平庸相。特别地，由于\cref{eq:H2}，对称相的数目总是有限的。该证明超出了本讲义的范围，但让我们指出，该论证严重依赖于 MPS 的 parent Hamiltonian 构造（见\cref{sec:NormalObs}）。

\bigskip\noindent 任何射影表示的一个重要性质是，当它对应于非平庸上同调类时，其矩阵维度必然严格大于 1。例如，考虑具有非平庸第二上同调群的最小群，即 $\mbb Z_2\times \mbb Z_2$，$H^2(\mbb Z_2\times \mbb Z_2,\rU(1))\simeq \mbb Z_2$。如果我们将其群元素用二元组 $(a,b)$（$a,b=0,1$）表示，则一个非平庸射影表示由 Pauli 矩阵给出如下：$X_{(a,b)}=(\sigma^X)^a(\sigma^Z)^b$，且该表示不能被进一步约化。这一事实对于处于非平庸 SPT 相中的 MPS 的纠缠具有重要推论，并且还保证了无能隙对称保护边缘模的存在~\cite{Pollmann2009}。

事实上，纠缠谱（参见\cref{sec:NormalObs}）的特征在于其简并度至少为该类中最小不可约射影表示的维度。这一点在自旋 1 反铁磁 Heisenberg 模型中得到了极为优美的体现。在\cref{fig:HeisenbergEntanglement}中绘制了前 48 个 Schmidt 值，能谱明显组织成至少具有二重简并的多重态。在这一情形下，保护对称性是 ${\rm SO}(3)$，其第二上同调群确实具有非平庸元素：$H^2({\rm SO}(3),\rU(1))\simeq \mbb Z_2$。${\rm SO}(3)$ 的所有不可约表示均具有奇数维，因此这些简并显然源自 ${\rm SU}(2)$ 的半整数不可约表示，后者确实构成了 ${\rm SO}(3)$ 的射影表示。事实上，该相也受到 ${\rm SO}(3)$ 的 $\mbb Z_2\times \mbb Z_2$ 子群的保护，因此我们甚至可以用破缺部分 ${\rm SO}(3)$ 对称性的项来微扰 Hamiltonian，只要这些项仍与 $\mbb Z_2\times \mbb Z_2$ 对易，系统就不会离开该 SPT 相。
\begin{figure}[htb]
    \centering
    \includegraphics[width=0.7\linewidth]{plots/HeisenbergEntanglementSpectrum.pdf}
    \caption{热力学极限下键维为 48 的自旋 1 Heisenberg 模型基态的纠缠谱。使用 MPSKit.jl 绘制~\cite{VanDamme}。}
    \label{fig:HeisenbergEntanglement}
\end{figure}

正如在\cref{sec:Examples}中所预期的，SPT 相的另一个标志是，当在开链上考虑时，存在受对称性保护的无能隙\emph{边缘模}（edge modes）。矩阵乘积态（及其高维推广，见\cref{sec:PEPS}）的一个特征是，这些边缘模的张量积基底由边界上的纠缠自由度提供。事实上，假设我们有一个键维为 $\chi$ 的均匀 MPS $\ket{\Psi[A]}$。我们可以在具有 $N$ 个格点的开链上考虑如下 $\chi^2$ 个态：
\begin{equation}
    \ket{\Psi[A]} = \sum_{\{i_j\}} (A^{i_1}A^{i_2}\cdots A^{i_N})_{\alpha\beta} \ \ket{i_1,i_2,\dots,i_N}.
\end{equation}
所有态 $\ket{\Psi[A],\alpha,\beta}$（$\forall \alpha,\beta=1,\dots,\chi$）对于 parent Hamiltonian 都具有相同的能量，因此表面上看我们拥有 $\chi^2$ 个简并的边缘模。然而，通过微扰 Hamiltonian，其中某些边缘模可以被打开能隙。尽管如此，只要我们限制在保持对称性的微扰上，边缘模就会在虚指标层面上按照射影表示进行变换：
\begin{equation}
    U(g)^{\otimes N}\ket{\Psi[A],\alpha,\beta} = \sum_{\alpha',\beta'} (X_g)_{\alpha\alpha'} (\bar X_g)_{\beta\beta'}\,\,\ket{\Psi[A],\alpha',\beta'},
\end{equation}
因此，边缘模至少保持二重简并。

\bigskip\noindent 至此，我们总结了\emph{玻色} $G$-对称有能隙一维相的分类。对于给定的对称群 $G$，有能隙相与在基态子空间中表示\emph{未破缺}对称性的子群 $H\subseteq G$，以及表示残存 $H$ 对称性的 SPT 相的上同调类 $[\psi]\in H^2(H,\rU(1))$ 之间存在一一对应。

顺便指出，此类数对 $(H,[\psi])$ 也表征了 $\Vect_G$ 上的（不可分解）\emph{模范畴}（module categories），其中 $\Vect_G$ 是 $G$-分阶向量空间的\emph{融合范畴}（fusion category）（参见\cref{sec:cat}）。因此，在 $\Vect_G$-模范畴与有能隙 $G$-对称相之间存在一一对应。当然，对于群对称性的简单情形，这一观察完全是多余的，但在最后一讲中我们将论证，对于一般的融合范畴对称性，这种对应关系依然成立。
\subsection{时空对称性}
上述讨论集中在 Hamiltonian 的内部对称性上，但在稍作修改后同样适用于时空对称性。在处理这些对称性时，必须解决两个在内部对称性情形下不存在的问题。首先，注意诸如平移对称性和空间反演等某些时空对称性是以非局域方式起作用的。另一方面，诸如时间反演或反演对称性等\emph{改变定向}（orientation-reversing）的对称性是以反幺正方式起作用的，因此依赖于所选取的基底。正如我们下文所见，后者对 SPT 相的上同调分类有着深刻的影响。尽管如此，MPS 仍然提供了一个合适的框架，在此框架下由空间对称性保护的对称相分类可以被明确导出~\cite{Pollmann2009,Vancraeynest2022}。

让我们说明类似于上述的推理如何导出由时间反演对称性保护的相的分类。通常，时间反演对称性表现为 $UK$，其中 $U$ 是幺正操作，在自旋 1/2 的情形下正如 Wigner 所证明的那样为 $(\sigma^Y)^{\otimes N}$~\cite{Wigner1960}，而 $K$ 是复共轭。为简便起见，我们在此假设 $U=\mbb 1$。用复共轭作用在单射 MPS $\ket{\Psi[A]}$ 上并应用基本定理，可得
\begin{equation}
    \bar A^i = e^{i\chi}QA^iQ^{-1}.
\end{equation}
随后再次施加时间反演对称性，并利用 MPS 的单射性，可知：
\begin{equation}\label{eq:TRProjective}
    Q\bar Q = \pm \mathbb 1.
\end{equation}
我们由此获得了一个对应于该符号的非平庸拓扑指数，从而给出时间反演 SPT 相的 $\mbb Z_2$ 分类。作为例证，考虑 AKLT 态\cref{eq:AKLTState}。容易验证在该情形下 $Q=\sigma^Y$，$Q\bar{Q}=-\mathbb 1$，因此 AKLT 态关于时间反演对称性属于一个非平庸 SPT 相。

注意，由于 $Q$ 上的复共轭，\cref{eq:TRProjective} 中的 $Q$ 不能被解释为 $\mbb Z_2$ 的非平庸射影表示。因此，\cref{eq:TRProjective} 中的相位 $\pm 1$ 也不能被解释为 $H^2(\mbb Z_2,\rU(1))$ 的非平庸上循环（后者实际上是平凡的）。因此，我们需要射影表示和上循环的推广概念，以适应某些群元素以反幺正方式起作用的对称群。容易验证，当一个一般对称群 $G$ 具有反幺正元素时，gauge 矩阵 $X_g$ 按照如下规则相乘：
\begin{equation}
    X_g\beta_g(X_h) = \psi(g,h) X_{gh},
\end{equation}
其中当 $g$ 为反幺正元素时 $\beta_g(\bullet)$ 表现为复共轭，否则表现为恒等映射。显然，当所有群元素均以幺正方式起作用时，该式等价于\cref{eq:ProjRep}；当对称群为 $\mbb Z_2$ 时间反演时，该式等价于\cref{eq:TRProjective}。同样容易验证，此时关于 $\psi(g,h)$ 的自洽性条件写作
\begin{equation}
    \psi(g,h)\psi(gh,k) = \psi(g,hk)\beta_g(\psi(h,k)), \quad \forall g,h,k\in G.
\end{equation}
相应地，上边缘的概念也必须稍作调整。综合起来，这意味着在存在反幺正对称性的情况下，SPT 相的上同调分类必须做出调整，以考虑非平庸的群作用 $\beta_g(\bullet)$。相应的上同调群记为 $H^2(G,\rU(1)^\beta)$，且通常不同于 $H^2(G,\rU(1))$。重要的是，该群可以通过前述并在\cref{sec:GroupCohomology}中详述的相同方法进行计算。与 AKLT 态的例子相对应，我们确实有 $H^2(\mbb Z_2,\rU(1)^\beta)=\mbb Z_2$。
\subsection{费米子 MPS 与对称相}\label{sec:Fermions}
到目前为止，我们仅关注了\emph{玻色}矩阵乘积态。本节的目标是概述将 MPS 框架推广至\emph{费米子系统}的方法，并强调由此产生的一些技术复杂性，在此紧密遵循文献~\cite{Bultinck2016}，亦参见文献~\cite{Kapustin2016,Mortier2024}。

\bigskip\noindent 至关重要的是\emph{超向量空间}（super vector space）的概念，即允许如下形式直和分解的向量空间：
\begin{equation}
    V = V^0 \oplus V^1,
\end{equation}
其中称 $V^0$（分别对应 $V^1$）中的矢量具有\emph{费米子偶}（fermion even）（分别对应\emph{奇}（odd））宇称。在不失一般性的前提下，我们假设基底矢量 $\ket{i}$ 具有确定的宇称，记为 $\lvert i \rvert = 0,1$。超向量空间此外还赋予了一个\emph{次数化张量积}（graded tensor product），在张量积重新排序下具有如下行为：
\begin{equation}
    \ket{i_1} \otimes_\mathfrak{g} \ket{i_2} \mapsto (-1)^{\lvert i_1 \rvert \lvert i_2 \rvert} \ket{i_2} \otimes_\mathfrak{g} \ket{i_1}.
\end{equation}
这一性质是费米子粒子交换统计的数学体现。\emph{超迹}（supertrace）则定义为
\begin{equation}
    \mc C (\ket{i}\otimes_\mathfrak g \bra{j}) = (-1)^{\lvert i \rvert \lvert j \rvert} \delta_{i,j}.
\end{equation}

一个均匀费米子 MPS（fMPS）由一个\emph{偶宇称张量}构造而成：
\begin{equation}
    A = \sum_{i,\alpha,\beta}A_{\alpha\beta}^i \ket{\alpha} \otimes_\mathfrak{g} \ket{i} \otimes_\mathfrak{g} \bra{\beta},
\end{equation}
即对于所有指标均满足 $\lvert i \rvert + \lvert \alpha \rvert + \lvert \beta \rvert = 0 \pmod 2$。根据所生成的态具有偶费米子宇称还是奇费米子宇称，构造略有不同。
\paragraph{偶宇称态}
偶宇称态由张量 $A$ 构造如下：
\begin{equation}\label{eq:EvenfMPS}
    \ket{\Psi}_e = \mc C_v\big(A\otimes_\mathfrak g A \otimes_\mathfrak g\dots \otimes_\mathfrak g A\big).
\end{equation}
在此表达式中，$\mc C_v$ 表示对虚指标求超迹。现在我们可以将该 fMPS 重写为更显式的形式。注意，一般偶宇称张量的矩阵元具有如下形式：
\begin{equation}
    A^i =
    \begin{pmatrix}
        B^i & 0 \\
        0 & C^i
    \end{pmatrix}, \,\, \lvert i \rvert = 0, \quad
    A^i =
    \begin{pmatrix}
        0 & D^i \\
        F^i & 0
    \end{pmatrix}, \,\, \lvert i \rvert = 1.
\end{equation}
定义\emph{宇称矩阵}（parity matrix）
\begin{equation}
    \mc P = \begin{pmatrix}\mbb 1 & 0\\ 0 & -\mbb 1\end{pmatrix},
\end{equation}
并仔细计算表达式\cref{eq:EvenfMPS}，便得到：
\begin{equation}
    \ket{\Psi[A]}_e = \sum_{\{i_j\}} {\rm Tr}\big(\mc PA^{i_1}A^{i_2}\dots A^{i_N}\big) \ket{i_1,i_2,\dots,i_N}.
\end{equation}
\paragraph{奇宇称态}
奇宇称态的构造与偶宇称态的构造类似，但需要在超迹中插入一个额外的矩阵：
\begin{equation}
    \ket{\Psi}_o = \mc C_v\big(Y\otimes_\mathfrak gA\otimes_\mathfrak g A \otimes_\mathfrak g\dots \otimes_\mathfrak g A\big).
\end{equation}
引入矩阵 $Y$ 的要求归根结底源于如下事实：要获得奇宇称态，必须引入扭扭边界条件（twisted boundary conditions）。我们将在第四讲的第~\ref{sec:GenSym}~节中遇到类似的讨论。

与前一种情形类似的推导使我们能够将奇宇称 fMPS 重写为
\begin{equation}
    \ket{\Psi[A]}_o = \sum_{\{i_j\}} {\rm Tr}\big(YA^{i_1}A^{i_2}\dots A^{i_N}\big) \ket{i_1,i_2,\dots,i_N}.
\end{equation}
随后施加平移不变性表明，$Y$ 可以选取为如下形式：
\begin{equation}\label{eq:fMPSY}
    Y =
    \begin{pmatrix}
        0 & -\mbb 1 \\ \mbb 1 & 0
    \end{pmatrix},
\end{equation}
并且在此基底中，MPS 矩阵写作：
\begin{equation}
    A^i =
    \begin{pmatrix}
        B^i & 0 \\
        0 & B^i
    \end{pmatrix}, \,\, \lvert i \rvert = 0, \quad
    A^i =
    \begin{pmatrix}
        0 & B^i \\
        -B^i & 0
    \end{pmatrix}, \,\, \lvert i \rvert = 1.
\end{equation}
与偶宇称情形相反，该 fMPS 仅由单个块 $B^i$ 确定。

注意，由于奇宇称 fMPS 矩阵与 $Y$ 对易，MPS 张量具有一个共同的不变子空间。然而，与玻色情形相反，当 $B^i$ 张满整个 $\mbb C^{\chi/2\times \chi/2}$ 矩阵代数时，奇宇称态 $\ket{\Psi[A]}_o$ 定义了一个不可约矩阵乘积态。关于费米子分阶 MPS 的 canonical form 的详细技术讨论与证明，请参阅文献~\cite{Piroli2020}的附录 D。在物理上，其深层原因可以追溯到费米子宇称的超选择定则（superselection rule）：奇宇称态与偶宇称态之间不存在叠加。

奇宇称 fMPS 态的一个显著范例是无自旋费米子 Kitaev 链的基态，其 Hamiltonian 是配对项的简单求和~\cite{Kitaev2000}：
\begin{equation}
    H = -i\sum_{\msf i=1}^{N-1} \gamma_{2\msf i}\gamma_{2\msf i+1},
\end{equation}
其中 Majorana 模满足 $\gamma_\msf i=\gamma^\dagger_\msf i$ 及 $\{\gamma_\msf i, \gamma_\msf j\}=2\delta_{\msf i,\msf j}$。平移不变的 fMPS 基态被发现具有奇宇称，其 MPS 矩阵为
\begin{equation}\label{eq:KitaevTensor}
    A^0 =
    \begin{pmatrix}
        1 & 0 \\
        0 & 1
    \end{pmatrix}, \quad
    A^1 =
    \begin{pmatrix}
        0 & 1 \\
        -1 & 0
    \end{pmatrix}.
\end{equation}
该模型以展现 Majorana 边缘模（Majorana edge modes）而著称。事实上，在开边界条件下，该态的形式为
\begin{equation}
    \sum_{\{i_j\}} \big(A^{i_1}A^{i_2}\dots A^{i_N}\big)_{\alpha\beta} \ket{i_1,i_2,\dots,i_N}.
\end{equation}
由此容易检验，由于张量的结构，选择边界条件 $\alpha=\beta=0$ 或 $\alpha=\beta=1$ 会给出同一个基态。类似地，边界条件 $\alpha=0,\beta=1$ 与 $\alpha=1,\beta=0$ 给出第二个基态，最终导致二重基态简并，这种简并无法通过在 Hamiltonian 中添加额外的局域相互作用来消除。

\bigskip\noindent 与玻色情形类似，当施加全局内部对称性或时空对称性时，人们可以从费米子矩阵乘积态中提取出某些离散拓扑不变量。一个著名的例子是相互作用时间反演对称费米子 SPT 相。在这种情况下，存在三个取值于 $\{0,1\}$ 的不同拓扑指数，其中一个与 fMPS 是否具有 Majorana 边缘模有关，另外两个与 fMPS 张量在时间反演下的局域变换性质有关。正如文献~\cite{Bultinck2016}中所严谨论证的，可以证明在相应 fMPS 的堆叠下，恢复了 Fidkowski 和 Kitaev 的 $\mbb Z_8$ 分类~\cite{Fidkowski2010}。
\subsection{第三讲小结} 第三讲探讨了对张量施加额外对称性时 MPS 的流形。对称张量网络的核心工具是 \emph{MPS 基本定理}，它指出任何\emph{全局}对称性都必须反映在\emph{局域} MPS 张量的对称性中。有趣的是，这些\emph{虚}对称性可以\emph{射影地}表示。按不同射影表示变换的对称 MPS 之间，任何保持对称性的绝热路径都必然经过无能隙点。因此，具有 $G$-对称 Hamiltonian 的一维量子自旋链的所有玻色有能隙物态，都与所有可能的子群 $H\subseteq G$ 及其由第二群上同调 $H^2(H,\rU(1))$ 标记的可能射影表示一一对应。后者对所谓的 \emph{SPT 相}进行了分类。如果不考虑对称性，那么所有 MPS 都可以绝热地相互转化，因此这些相仅在施加对称性时才受到保护。我们证明了这一分类可以推广到时空对称群的情形。有趣的是，当考虑超选择定则（例如费米子的宇称超选择定则）时，会出现真正的拓扑物态（即不需要对称性保护的相）。
\newpage\section{第四讲：MPO 代数：从代数 Bethe ansatz 到 G-单射 PEPS}
现在让我们将注意力转向可积模型（integrable models），这是本讲习班的核心主题之一。事实证明，张量网络提供了一种非常简明的方式，可以用矩阵乘积态与算符来重新表述\emph{代数 Bethe ansatz}（algebraic Bethe ansatz, ABA）。特别是，我们将看到构成其可积性核心的 \emph{Yang-Baxter 方程}（Yang-Baxter equation），可以被理解为基本定理（\cref{sec:FundamentalTheorem}）的一个推论。尽管张量网络技术并未向我们揭示关于 Bethe ansatz 的任何根本全新的物理，但本讲的见解将对由\emph{融合范畴}（fusion categories）所描述的\emph{广义}或\emph{范畴对称性}（categorical symmetries）的 MPO 表示论发挥关键作用，而这正是最后一讲的主题。

在讨论完代数 Bethe ansatz 之后，我们将在\cref{sec:PEPS}中重点讨论 MPS 的二维推广，即\emph{投影纠缠对态}（projected entangled-pair states, PEPS）。我们将猜想基本定理的一种推广形式，并讨论 MPO 对称性在二维 SPT 与拓扑相及其边缘模的张量网络表示中所起的作用。
\subsection{磁振子态}\label{sec:magnon}
张量网络可以用作 Hamiltonian 解析解的 ansatz，正如针对 AKLT 模型\cref{eq:AKLTState}所展示的那样。然而，本讲所关注的范例是一类著名的精确可解模型，称为\emph{可积系统}（integrable systems）。粗略地说，可积量子系统是具有广延数目的守恒荷（conserved charges）——即与 Hamiltonian 对易的局域算符——的系统，其精确可解性正归功于此。可积系统可以用称为 \emph{Bethe ansatz} 的 ansatz 波函数来求解（参见文献~\cite{Gaudin2014}）。该方法有两种表述形式，分别称为\emph{坐标 Bethe ansatz}（coordinate Bethe ansatz）与\emph{代数 Bethe ansatz}（algebraic Bethe ansatz），它们产生同一组解。这两种表述形式都可以使用张量网络写出。

为了具体起见，我们考虑自旋 $1/2$ XXZ Hamiltonian
\begin{equation}
    H_{XXZ}=\sum_{\msf i=}^L \sigma^X_{\msf i}\sigma^X_{\msf i+1}+\sigma^Y_\msf i\sigma^Y_{\msf i+1}+\Delta \sigma^Z_\msf i\sigma^Z_{\msf i+1}
\end{equation}
在周期性边界条件（即 $L+1\equiv 1$）下。求解 Bethe 方程时用作基底的\emph{磁振子态}（magnon states）可以很容易地写为 MPS。例如，\emph{周期性}边界条件下的 1-磁振子态可以写为\emph{开}边界条件 MPS（参见\cref{eq:OpenMPS}），其左边界矢量为 $\vec v_L^\dagger = \bra{0}$，右边界矢量为 $\vec v_R = \ket{1}$，MPS 张量为：
\begin{equation}
\begin{split}
A^0(k) = 
\begin{pmatrix}
    e^{ik} & 0 \\
    0 & 1
\end{pmatrix}, \quad
A^1(k) = \begin{pmatrix}
    0 & e^{ik} \\
    0 & 0
\end{pmatrix} \\
\implies \ket{\Psi_1[A]} = \sum_{\msf n=1}^L e^{ik\msf n} \ket{0_0 ...0_{\msf n-1} 1_\msf n 0_{\msf n+1}...0_{L}}
\end{split}.
\label{AAAA}
\end{equation}
注意，这正是在\cref{sec:Excitations}中在 MPS 流形上将激发态构造为具有动量的切矢量的方式。一般 n-磁振子态可以在文献~\cite{Alcaraz2004,Murg2012,Haegeman2016}中找到，并涉及一个键维随系统中粒子数呈指数级增长的 MPS。
\subsection{作为 MPS 基本定理的 Yang-Baxter 方程}\label{sec:YangBaxterFromMPO}
然而，还有一种完全不同的方法，即借助代数 Bethe ansatz（ABA）导出精确本征态。在用张量网络表述这两种方法之前，并不清楚这两种方法彼此之间有何关联；现在人们已经清楚，正如由 MPS 基本定理所必然要求的那样，这两种方法通过一个非平庸的 gauge 变换相关联~\cite{Katsura2010}。在代数 Bethe ansatz 中，Hamiltonian 的所有本征态都可以通过依赖于连续\emph{谱参数}（spectral parameter）$\mu_i$ 的 MPO 相乘来构造：
\begin{equation}\label{eq:BetheAnsatzMPS}
    \inputtikz{YB_ABA}.
\end{equation}
此处，所有 MPO 张量均具有如下形式：
\begin{equation}
    A_{\alpha\beta}^{ij}(\mu) =\mu\delta_{\alpha\beta}\delta_{ij}+\delta_{\alpha i}\delta_{\beta j}= \mu \inputtikz{YB_A_1} + \inputtikz{YB_A_2}.
    \label{AA}
\end{equation}
容易检验，每个 MPO 都会向系统中增加一个磁振子，因为左边界为 $\bra{0}$，右边界为 $\ket{1}$。系数 $\mu_i$ 是磁振子动量的函数，可以通过求解所谓的 \emph{Bethe 方程}（Bethe equations）来确定。ABA 的一个关键要素是所有这些 MPO 均相互对易。

其背后的逻辑如下。我们从一个二维统计力学模型的转移矩阵出发，该转移矩阵可以写为单射 MPO $T(\mu)$ 的形式，由与式~\ref{AA}~中相同的张量 $A(\mu)$ 确定，但此时具有周期性边界条件。事实证明，该系统可积当且仅当所有这些 MPO 与其谱参数无关地相互对易：
\begin{equation}
        \forall \lambda, \mu: [T(\lambda), T(\mu)] = 0.
\end{equation}
在此情形下，我们可以调用 MPS 基本定理\cref{eq:FundTh}来证明存在一个将这些乘积相互转化的 gauge 变换 $R$~\cite{Haegeman2016}。事实上，注意我们可以将 MPO 的两条腿合并为一条腿，从而可以将 MPS 基本定理用于秩为 4 的 MPO。具体而言，存在一个满足如下方程的交织算符（intertwiner）张量 $R(\lambda, \mu)$：
\begin{equation}\label{eq:YB_intertwiner}
    \inputtikz{YB_intertwiner_1} \,=\, \inputtikz{YB_intertwiner_2}\,\,.
\end{equation}
该 $R$-张量是由对易 MPO 给出的全局对称性的局域实现。

现在让我们取 $R$-张量的三个副本。由于 MPO 代数的结合律，这些 $R$-张量可以通过两种不同方式融合在一起，从而导出如下方程：
\begin{equation}
    \inputtikz{YB_1} \,=\, \inputtikz{YB_2}\,\,.
\end{equation}
这正是 Yang-Baxter 方程——因此 MPS 基本定理以非常简明的方式重现了这一结果。给定任意可积转移矩阵，相应的交织算符因此必须满足 Yang-Baxter 方程。可以检验，\cref{AA} 中的张量正是这些解之一。因此，寻找可积模型的一种构造性方法是将思路反转过来：首先寻找 Yang-Baxter 方程的一个解（这通常极其困难！），从而得出张量 $R(\lambda,\mu)$ 的显式表示。下一步，将转移矩阵 $T$ 作为\cref{eq:YB_intertwiner}的解求解出来。令人惊叹的是，通过将 MPO 张量识别为 $A_{\alpha \beta}^{ij}(\mu) := R_{j \beta}^{i \alpha}(\mu, 0)$，该方程等价于 Yang-Baxter 方程本身。因此，Yang-Baxter 代数的这一所谓\emph{基本表示}（fundamental representation）允许构造单参数族对易转移矩阵。注意，这并不是求解\cref{eq:YB_intertwiner}的唯一 MPO 族。人们可以构造同一代数的不同表示——在后文中，我们将遇到类似的情形，届时我们将构造同一对称系统的不同表示。作为最后一步，现在可以用相同的张量 $A$ 但采用不同的边界条件，来构造\cref{eq:BetheAnsatzMPS}中的磁振子 MPO。

因此，张量网络表述的核心原理与优势如下：\emph{全局}守恒律和对易关系可以写成涉及\emph{局域}自由度（即 MPS 和 MPO 张量）的方程，而这些方程通常更容易求解。

在给出 MPO 表述之后，还可以恢复可积 Heisenberg Hamiltonian。所有权重彼此相等的 6-顶点模型（相应的可积统计力学模型）的转移矩阵恰好具有\cref{AA}的形式。根据构造，具有这些张量 $A(\mu)$ 的所有 MPO $T(\mu)$ 都将相互对易，因此该 MPO 的导数是唯一确定的。Heisenberg Hamiltonian 确实是该转移矩阵的对数导数；我们将证明留作一个（非常有趣的）练习（详尽讨论参见文献~\cite{Murg2012}）：
\begin{equation}
    H = T^{-1}(0) \frac{{\rm d}T(\mu)}{{\rm d}\mu}\vert_{\mu=0} = \sum_{\msf i}\quad\inputtikz{YB_Ham}\,.
\end{equation}

总之，两种 Bethe ansatz 表述（坐标和代数）在张量网络中都有自然的表示。张量网络与 Bethe ansatz 显然共享着相同的基因，但张量网络已被证明具有更广泛的适用性——它们同样适用于不可积系统和更高维度。
\subsection{投影纠缠对态}\label{sec:PEPS}
矩阵乘积态向二维的自然推广是所谓的\emph{投影纠缠对态}（projected entangled-pair states, PEPS）~\cite{Verstraete2004b}。与一维情形类似，平移不变的 PEPS 由单个张量构造而成，在正方晶格的情形下，该张量具有一个物理指标和四个虚指标：
\begin{equation}
    A = \sum_{\substack{i,\alpha,\beta\\\gamma,\delta}} A^i_{\alpha\beta\gamma\delta} \ket{i}\otimes\bra{\alpha}\otimes\bra{\beta}\otimes\bra{\gamma}\otimes\bra{\delta} \in \mbb C^{d\chi^4}.
\end{equation}
该张量如图所示：
\begin{equation}
    \inputtikz{PEPSTensor} = \ A^i_{\alpha\beta\gamma\delta}, \qquad \substack{i=1,2,\dots,d\\ \alpha,\beta,\gamma,\delta=1,2,\dots,\chi},
\end{equation}
从而整个 PEPS 波函数可以表示为
\begin{equation}
    \ket{\Psi[A]} = \inputtikz{PEPS}\,\, .
\end{equation}
注意，PEPS 自然满足面积律（area law），因为一段连续自旋区域的纠缠熵随该区域边界上的键数（即其周长）线性增长。因此，PEPS 预期为局域有能隙 Hamiltonian 的低能态构成一个良好的变分族，而长期以来人们一直猜想这些低能态满足面积律。反之，与 MPS 一样，可以证明每个 PEPS 都是某个局域无阻挫（frustration free）有能隙 parent Hamiltonian 的精确基态。如何将这种构造推广到开边界条件、任意晶格或更高维度也是显而易见的。类似地，人们可以将 MPO 推广到二维情形。此类算符被称为\emph{投影纠缠对算符}（projected entangled-pair operators, PEPO）。

人们付出了相当大的努力来将 DMRG 和 VUMPS 算法推广到 PEPS 的情形。虽然原则上将流形图像推广到 PEPS 是直截了当的，但其关于键维的实际计算复杂度增长远差于一维情形。尽管如此，这些 PEPS 算法仍在二维强关联晶格系统（例如 Hubbard 模型）中提供了最前沿的变分结果，提出更好的算法是一个非常活跃的研究领域。关于最前沿的结果，我们推荐参阅文献~\cite{Verstraete2008,Vanderstraeten2016,Corboz2016,Liao2019,Vanderstraeten2022}。在这些讲义中，我们将侧重于 PEPS 的理论。

与矩阵乘积态的情形相反，目前尚不存在完全通用的基本定理，能够为两个不同的 PEPS 张量 $A$ 和 $B$ 在所有系统尺寸下生成相同的态提供充分必要条件。然而很明显，使其成立的一个充分条件是存在一个矩阵乘积算符，按照如下方式\emph{交织}（intertwines）PEPS 张量：
\begin{equation}
    \inputtikz{PEPSPullingThrough_1} \,\,=\,\,
    \inputtikz{PEPSPullingThrough_2}\,\,,
\end{equation}
以及所有其他穿透（pulling through）MPO 的方式，例如：
\begin{equation}
    \inputtikz{PEPSPullingThrough_3} \,\,=\,\,
    \inputtikz{PEPSPullingThrough_4}\,\,.
\end{equation}
在缺乏严格证明的情况下，我们称其为 PEPS 的基本\emph{猜想}（conjecture）。假设周期性边界条件，对称性作用的实现方式是：从 PEPS $\ket{\Psi[A]}$ 出发，成核一个由交织 MPO 构成的小泡，将其穿透整个晶格并与自身湮灭，最终得到 PEPS $\ket{\Psi[B]}$。正如我们将在本讲义后续部分所看到的，大类 PEPS 恰好是按照这一方案来实现其对称性的。在某些情况下，该 MPO 只是一个直积（键维为 $1$），从而虚对称性仅作为直积作用在虚腿上。

在接下来的各节中，我们将讨论若干表示拓扑物态的特别有趣的 PEPS 波函数。拓扑波函数可以使用张量网络形式轻松构造，这显然是它们受到如此广泛研究的主要原因之一。此外，这些张量网络表示将长程纠缠与作用在虚自由度上的（MPO）对称性等同起来。这也是该基本猜想如此重要的原因：拓扑序由非平庸的 MPO 对称性所表征。
\subsection{二维 SPT：案例研究}\label{sec:SPT2d}
具有非平庸 MPO 对称性的 PEPS 的典型范例是 Chen、Liu 和 Wen 在文献~\cite{Chen2011}中引入的所谓 \emph{CZX 模型}（其张量网络表示参见文献~\cite{Buerschaper2013,Williamson2016,Molnar2017}）。他们的模型是由 $\mbb Z_2$ 对称性保护的非平庸 SPT 的一个例子，该对称性在 PEPS 的物理层面上以单格点方式 $U_g^{\otimes N}$ 起作用。在 PEPS 张量的虚指标层面上，该对称性通过键维为 2 的 MPO 实现为：
\begin{equation}\label{eq:PEPS_CZX}
    \inputtikz{PEPS_CZX_1} = 
    \inputtikz{PEPS_CZX_2}\,\, ,
\end{equation}
其中 $g\in\{0,1\}$。由 $g$ 索引的 MPO 构成了 $\mbb Z_2$ 的非局域表示，其图形表示为
\begin{equation}\label{eq:GroupMPO}
    \inputtikz{MPOGroupStacked} =\,
    \inputtikz{MPOGroupgh}\,\, .
\end{equation}
乘法\cref{eq:GroupMPO}可以通过满足如下条件的交织张量在局域上实现\footnote{$g=0$ 对应的 MPO 不是单射的，但不能分解为单射块。这可以通过底层的非半单代数结构来解释~\cite{Liu2025,FloridoLlinas2025}。因此，无法直接通过调用基本定理来找到交织算符。}~\cite{Chen2011,Williamson2016,Liu2025}：
\begin{equation}\label{eq:GroupIntertwiner}
    \inputtikz{GroupMPOIntertwiner_1} = \inputtikz{GroupMPOIntertwiner_2}\,\, .
\end{equation}
回想在\cref{sec:SPT1d}中我们曾论证，在一维 SPT 相的情形下，存在一个取值于第二上同调群 $H^2(G,\rU(1))$ 的拓扑指数。该指数完全由构成对称性射影表示的虚矩阵的乘法规则所捕获。现在很自然地会想知道二维 SPT 相对应的拓扑指数是什么，以及它是如何编码在 MPO 代数中的。事实证明，如果考虑表示 $\mbb Z_2$ 对称性的三个 MPO 的乘法，那么上面定义的融合张量可以通过两种不同方式进行收缩。这两种方式必须相互兼容，因此必然存在一个依赖于三个群元素的相位 $\omega$，满足
\begin{equation}\label{eq:GroupMPOFusion}
    \inputtikz{GroupMPOPentagon_1}
    = \omega(g,h,k) \,\,
    \inputtikz{GroupMPOPentagon_2}\,\, .
\end{equation}
矩阵乘法的结合律表明 $\omega$ 必须满足
\begin{equation}
    \omega(g,h,k)\omega(g,hk,l)\omega(h,k,l) = \omega(gh,k,l)\omega(g,h,kl).
\end{equation}
该条件被称为 \emph{3-上循环方程}（3-cocycle equation）（亦参见\cref{sec:GroupCohomology}）。注意，我们可以给融合张量自由地乘以一个相位 $\varphi(g,h)$，在此重新定义下，3-上循环按照如下规则变化：
\begin{equation}\label{eq:3cocycle}
    \omega(g,h,k) \mapsto \omega(g,h,k) \frac{\varphi(g,h)\varphi(gh,k)}{\varphi(g,hk)\varphi(h,k)}.
\end{equation}
模去此类重新定义后，我们可以得出结论：指数 $\omega$ 取值于第三上同调群 $H^3(G,\rU(1))$。该群可以通过\cref{sec:GroupCohomology}中给出的方法进行计算，并且特别地它也是有限群。对于上述考虑的 $G=\mbb Z_2$ 情形，有 $H^3(\mbb Z_2,\rU(1))\simeq \mbb Z_2$，因此恰好存在一个由 $\mbb Z_2$ 对称性保护的非平庸 SPT 相。相应非平庸类的代表元 3-上循环为
\begin{equation}
    \omega(g,h,k) = (-1)^{ghk},
\end{equation}
这可以直接从 CZX 模型的融合张量的\cref{eq:GroupIntertwiner}中推导出来。正如在量子自旋链的情形中一样，具有不同虚对称性/3-上循环的两个 PEPS 波函数无法在不闭合能隙的情况下通过保持对称性的绝热路径相互转化。
\subsection{边缘 Hamiltonian 与异常}
当 PEPS 定义在开面片 $\mc A$ 上时，出现了一个值得注意的特征。在这种情况下，未收缩的键定义了一个维度为 $\chi^{\lvert\partial \mc A\rvert}$ 的希尔伯特空间，其中 $\lvert\partial\! \mc A\rvert$ 是与边界相交的连键数目。图形表示为：
\begin{equation}
    \inputtikz{PEPS_open}\,\,.
\end{equation}
这些边缘自由度为相应 parent Hamiltonian 的边缘模定义了一个张量积结构，这与我们在\cref{sec:SPT1d}中识别一维情形下 SPT 边缘模的方式完全一致。当 PEPS parent Hamiltonian 随后受到微扰时，这会在这些边界自由度上诱导出一个\emph{有效} Hamiltonian，并有可能驱动它们经历相变~\cite{Yang2014}。有趣的是，在 SPT 的 PEPS 表示情形下，这些有效边界 Hamiltonian 继承了体态的非局域 MPO 对称性。有关若干示例及其数值研究，请参阅文献~\cite{Williamson2016,Roose2018}。

给出一个群的 MPO 表示，反过来提出了对称（边缘）Hamiltonian 的相图会是什么样的问题。正如我们现在将要展示的，这个问题的答案强烈取决于与 MPO 代数相关联的 3-上循环是否平庸。正如文献~\cite{Chen2011,Cirac2020}中所证明的，一个非常优美且直接的张量网络论证表明，只要 $\omega$ 在上同调上是非平庸的，与之对易的 Hamiltonian 就必须要么自发破缺对称性，要么是无能隙的。在这种情况下，称上循环 $\omega$ 代表了对称性的\emph{异常}（anomaly）。这一观察与著名的 Lieb-Schultz-Mattis 定理~\cite{Lieb1961,Oshikawa2000,Hastings2003}高度平行。论证过程如下。

在该论证中，我们假设任意群 $G$ 的一个\emph{单射} MPO 表示。由于单射性，我们现在可以调用基本定理来证明存在局域实现群乘法的融合张量，如\cref{eq:GroupIntertwiner}所示。现在假设按\cref{eq:GroupMPOFusion}定义的 3-上循环是非平庸的，并假设存在一个由单射 MPS 表示的唯一的、平庸有能隙的对称基态。在图形上，这归结为：
\begin{equation}
    \inputtikz{GroupMPOMPS} \,\, = \,\,
    \inputtikz{GroupMPS} \,\, .
\end{equation}
随后应用 MPS 基本定理表明，必然存在一个满足如下条件的三价\emph{作用张量}（灰色）：
\begin{equation}
    \inputtikz{GroupMPSIntertwiner_1} \,\, = \,\,
    \inputtikz{GroupMPSIntertwiner_2} \,\, .
\end{equation}
与上一节类似，如果我们在态上作用两个对称 MPO，则融合 MPO 张量随后作用在 MPS 上的结果，在相差一个相位的意义下等于相继作用这些 MPO。具体而言：
\begin{equation}
    \inputtikz{GroupMPSIntertwinerRecoupling_1} \,\, = \,\, \lambda(g,h)
    \inputtikz{GroupMPSIntertwinerRecoupling_2} \,\, .
\end{equation}
如果我们现在考虑三个此类 MPO 在我们的 MPS 上的作用，并考虑在局域重新耦合它们的两种不同方式，我们便可以推导出如下条件：
\begin{equation}
    \frac{\lambda(g,h)\lambda(gh,k)}{\lambda(h,k)\lambda(g,hk)}=\omega(g,h,k),
\end{equation}
但左边可以被识别为一个上边缘，这反过来意味着 3-上循环在上同调上是平庸的。由于这与 $\omega$ 是非平庸的假设相矛盾，我们因此必须得出结论：不存在对称单射 MPS，并且每个具有 MPO 对称性的一维 Hamiltonian 都必须是对称性破缺的或无能隙的。注意，无论该 Hamiltonian 编码的是二维 SPT 相的边缘模还是真实的量子自旋链，该论证均成立。
\subsection{拓扑序：量子双重模型与管代数}\label{sec:QuantumDoubles}
引人注目的是，除 SPT 相之外，在高于一维的空间维度中还存在不受任何对称性保护的鲁棒量子物态。称此类相展现出\emph{内禀拓扑序}（intrinsic topological order）~\cite{Wen1989,Witten1988,Zeng2015,Wen2016}，其特征在于仅依赖于底层空间流形拓扑（球面、环面等）的基态简并度、具有\emph{任意子}编织统计的有能隙激发，以及根据其任意子内容决定的对纠缠熵面积律的普遍负修正~\cite{Kitaev1997,Kitaev2005a,Levin2006}。值得注意的是，与不存在对称性紧密相关的是，这些相不能由局域序参量来表征，从而超出了对称性破缺的传统 Landau 范式的范畴~\cite{Landau1937,Wegner1971,Bravyi2006}。尽管如此，存在基态简并度的事实清楚地表明某种对称性被破缺了。事实证明，这种对称性纯粹作用在虚自由度上。

回想短程纠缠的 SPT 序的稳定性源于与全局对称性下纠缠自由度的变换性质相关联的拓扑指数的存在（\cref{sec:SPT1d} 和 \ref{sec:SPT2d}）。另一方面，内禀拓扑序的鲁棒性源于其长程纠缠~\cite{Chen2010}。在某种意义上，表征拓扑序因此归结为对长程纠缠的等价类进行分类。类似地，这些拓扑相之间的相边界可以从纠缠性质的根本性改变来理解~\cite{Haegeman2015,Marien2016}。正如在 PEPS 中那样能够直接处理这些纠缠自由度，这启示我们可以进一步将该问题约化为对产生\emph{全局}拓扑态的\emph{局域} PEPS 张量进行分类：内禀拓扑序系统的分类等价于构造具有 MPO 对称性的张量。

事实上，用 PEPS 的语言来说，拓扑序的存在可以追溯到与基态 PEPS 张量对易的\emph{虚} MPO 对称性的存在：
\begin{equation}\label{eq:PEPSPullingThrough}
    \inputtikz{PEPS_TopoPulling_1} \,=\, \inputtikz{PEPS_TopoPulling_2}\,\,,
\end{equation}
注意，与 SPT 情形相反，在物理自由度上没有任何作用。因此，这些 MPO 在虚指标层面上是完全可形变的——即\emph{拓扑的}。这也意味着底层的对称性结构可以比群所描述的结构丰富得多。正如我们将在第~\ref{sec:GenSym}~讲中所看到的，对称代数由一个\emph{融合范畴}（fusion category）编码。由于这些 MPO 对物理自由度不敏感，此类态对物理层面上的（微小）局域微扰具有鲁棒性~\cite{Marien2016}。然而，\emph{任意子凝聚}（anyon condensation）提供了一种在足够强的微扰下驱动系统走向平庸相的机制~\cite{Bais2008,Burnell2017}。

人们猜想，非手征拓扑序完全由其任意子内容以及融合规则与编织统计所表征。在数学上，这些数据被封装在一个称为 \emph{Drinfeld 中心} $\mc Z(\mc D)$ 的\emph{模张量范畴}（modular tensor category）中~\cite{Etingof2015,Kong2017,Bultinck2015,Kong2022}。该论断的完整范畴以及一般拓扑序的微观实现将是\cref{sec:SN}的主题。在本节中，我们将重点关注一类称为（无扭）\emph{量子双重模型}（quantum doubles）~\cite{Dijkgraaf1989,Roche1990,Kitaev1997}的拓扑序的构造，作为对拓扑序的 PEPS 表示及其若干性质的更为温和的引介。

\bigskip\noindent Kitaev 在文献~\cite{Kitaev1997}中引入的量子双重模型是由有限群 $G$ 构造的一族模型，方法是在二维晶格的边上赋予 $\lvert G \rvert$ 维 qudit，并服从上述文献中详述的对易投影算符 Hamiltonian。它们是无扭 Dijkgraaf-Witten 拓扑场论的晶格 Hamiltonian 实现~\cite{Dijkgraaf1989}。对于群 $\mbb Z_2$，该模型约化为\emph{复环面码}（toric code）。这些量子双重模型的基态允许非常简单的重整化群不动点 PEPS 表示，称为 \emph{$G$-单射}（G-injective），或者更精确地称为 \emph{$G$-等距}（G-isometric）PEPS。$G$-等距的概念本质上抓住了仅在群 $G$ 作用下不变的虚自由度子空间上可逆的 PEPS 张量，是 MPO 单射性的一种具体体现。

在局域物理变换的意义下，所有 $G$-等距 PEPS 都可以约化为我们现在将要描述的普适形式。在本节中，我们假设二维环面上的正方晶格。我们用 $L$ 表示 $G$ 的\emph{左 regular 表示}（left regular representation），其作用为 $L_g\ket{h} = \ket{gh}$。然后定义投影算符 $1/\lvert G \rvert\sum_g L_g\otimes L_g\otimes L_g^\dagger\otimes L_g^\dagger$，它构成了基态 PEPS 张量，图形表示为：
\begin{equation}\label{eq:PEPSGTensor}
    \frac{1}{\lvert G \rvert}\sum_{g\in G}\,\,\, \inputtikz{PEPS_G_1}\,\,.
\end{equation}
在此图中，灰色方框代表 $L_g$，横线表示逆群元素（回想 $L_g^\dagger = L_{\bar g}$）。注意，最内侧的四条腿共同构成了物理层面。由于式~\ref{eq:PEPSGTensor}~中的求和遍历所有群元素，该 PEPS 张量对于如下由群元素标记的 MPO 满足穿透条件\cref{eq:PEPSPullingThrough}：
\begin{equation}\label{eq:MPOGroupRegular}
    \inputtikz{MPOGroupg}:= \bigotimes_\msf i L_g,
\end{equation}
其中虚线用于强调键维为 1，即每个 MPO 均表现为直积。

PEPS 张量\cref{eq:PEPSGTensor}仅给出了该模型的一个基态。正如文献~\cite{Schuch2010}中所展示的，所有其他基态都可以通过用环绕环面任一或两个不可缩回圈的算符作用在该态上而幺正地获得。这些算符的广延性质反映了基态的局域不可区分性。等价地，可以通过在虚指标层面上用 MPO 环绕环面的两个圈并在其交叉点处用一个需要求解的张量 $\mc P_Z$ 将它们连接起来，从而获得所有基态。在此过程中，我们要求基态均具有平移不变性，并在热力学极限下保持可归一化和正交~\cite{Zhang2011}。利用 MPO \cref{eq:MPOGroupRegular} 的融合张量，我们可以将这样的态展开如下：\footnote{注意，此时的融合张量仅仅强制群乘法，即 $\sum_{g,h,k}\ketbra{k}{g,h}=\sum_{g,h}\ketbra{gh}{g,h}$。}
\begin{equation}
    \inputtikz{PEPS_G_idempotent_1} \,\, = \sum_{g,h} \mc (\mc P_Z)_{g,h}\,\,\inputtikz{PEPS_G_idempotent_2}\,\,.
\end{equation}
这些态的平移不变性由穿透条件显式保证。正如文献~\cite{Sahinoglu2014,Bultinck2015,Williamson2017}中所论证的，MPO 单射 PEPS 中的一组基态基底与 \emph{Ocneanu 管代数}（tube algebra）~\cite{Ocneanu1993,Ocneanu2001}的\emph{极小中心幂等元}（minimal central idempotents）——或者等价地，不可约表示——一一对应。事实上，在实践中 Drinfeld 中心 $\mc Z(\mc D)$ 通常就是通过管代数计算的。对于量子双重模型的特定情形，相关的管代数由如下形式的闭合 MPO 张成：
\begin{equation}
    T_x^g := \inputtikz{PEPS_G_Tube}\,\, ,
\end{equation}
可将其视为从外侧指标作用到内侧指标。它们的乘法定义为将它们相互套叠。在代数上这归结为：
\begin{equation}\label{eq:GTubesMult}
    T^g_x\circ T_y^h = \delta_{x,hyh^{-1}} T^{gh}_x.
\end{equation}
该乘法定义了一个通常称为群 $G$ 的\emph{（无扭）量子双重模型}（quantum double）$\mc D(G)$ 的代数~\cite{Roche1990}。值得注意的是，管代数赋予了一个 \emph{star} 结构，即交换乘积顺序的反线性对合。对于当前情形：$\big( T_x^g\big)^\star := T^{g^{-1}}_{gxg^{-1}}$。
因此，管代数被赋予了 \emph{$C^\star$-代数}的结构。根据著名的 Artin-Wedderburn 定理，每个 $C^\star$-代数都可以分解为单矩阵代数。前面提到的极小中心幂等元便充当投射到每个单块上的投影算符。

让我们显式地计算它们。为此，对于固定的 $x\in G$，首先关注所有满足 $g\in Z_x=\{g\in G \mid gx=xg\}\subseteq G$ 的管 $T_x^g$ 是非常有益的。这些管构成了稳定子群 $Z_x$ 的一个表示，因为
\begin{equation}
    T^{g_1}_x\circ T^{g_2}_x = T^{g_1g_2}_x, \quad \forall g_1,g_2\in Z_x,
\end{equation}
正如直接由\cref{eq:GTubesMult}所给出的那样。令 $\mu_x$ 表示 $Z_x$ 的一个维度为 $d_{\mu_x}$ 且相应特征标为 $\chi^{\mu_x}$ 的不可约表示。算符
\begin{equation}
    \mc P_{(x,\mu)} := \frac{d_{\mu_x}}{\lvert G \rvert} \sum_{k\in Z_x} \chi^{\mu_x}(k) T_x^k,
\end{equation}
构成了 Hermitian 投影算符，即\footnote{利用特征标的 Schur 正交性 $\sum_k \chi^{\mu_x}(k)\chi^{\nu_x}(k^{-1}m) = \lvert Z_x \rvert/d_\mu \chi^\mu(m)$。}
\begin{equation}
    \mc P_{(x,\mu_x)} \mc P_{(y,\nu_y)} = \delta_{x,y}\delta_{\mu_x,\nu_x} \mc P_{(x,\mu_x)}.
\end{equation}
这些 $\mc P_{(x,\mu)}$ 是\emph{单}幂等元（simple idempotents）。中心的幂等元随后通过对共轭类 $C_x = \{gxg^{-1} \mid g\in G\}$ 中的所有 $x$ 额外求和而获得：
\begin{equation}
    \mc P_{(C_x,\mu_x)} \equiv \sum_{y\in C_x} \mc P_{(y,\mu_x)}.
\end{equation}
这重现了文献~\cite{Kitaev2005a}的结果，即量子双重模型的环面基态由共轭类 $C_x$ 以及 $C_x$ 的代表元 $x$ 的稳定子子群 $Z_x$ 的不可约表示所标记。

重要的是，任意子激发也由管代数的极小中心幂等元所标记~\cite{Lan2013,Hu2015}。在 PEPS 层面上，可以通过将基态张量中的两个替换为支集包含在标记该激发的相应幂等元的像中的张量，来创建一对任意子激发，详细内容请参阅文献~\cite{Bultinck2017}。

这里考虑的（无扭）量子双重模型可以在许多方面进行推广。一方面，可以考虑该模型的\emph{有扭}（twisted）版本，其中 MPO 对称性构成 $G$ 的一个反常表示，其异常取值于第三上同调群 $H^3(G,\rU(1))$，正如在\cref{sec:SPT1d}中所描述的那样。对此，我们推荐参阅文献~\cite{Buerschaper2013,Bultinck2016,Williamson2017}。

在\cref{sec:SN}中，我们将讨论 Levin 和 Wen 在文献~\cite{Levin2004}中首次构造的所谓弦网模型（string-net models）的 PEPS 表示。这是一类以球面融合范畴为输入并实现所有非手征拓扑序 $\mc Z(\mc D)$ 的模型。在这些模型中获得由 $\mc Z(\mc D)$ 所描述的环面基态和任意子激发的一种方法，恰好是通过与 $\mc D$ 相关联的管代数，这与上述构造完全一致。对于输入 $\mc D=\Vect_G^\omega$，相应的 Levin-Wen 模型与有扭量子双重模型相重合。
\subsection{第四讲小结} 
第四讲介绍了三个新主题：构成特殊类量子自旋链可积性基础的 MPO 形式的对称性；MPS 向 PEPS 的高维推广，以及参数化有能隙高维晶格模型低能态的相应流形；以及满足所谓\emph{穿透对称性}（pulling through symmetry）的 MPO 对称性，该对称性表征了此类系统的不同有能隙（拓扑）相。关于可积性，我们论证了 Yang-Baxter 方程是 MPS 基本定理的一个推论。如果这种观点能够引向更一般的可积模型类别（例如满足非平庸范畴对称性的模型，参见第五讲），那将是令人振奋的。

我们展示了二维晶格系统中的 SPT 相由其结合子编码了 3-上循环的非平庸 MPO 对称性所表征，从而导出了二维 SPT 相的分类。然而，在不涉及全局对称性的情况下，PEPS 的穿透对称性为在二维系统中实现拓扑物态开辟了全新的途径。我们以量子双重模型的形式展示了内禀拓扑序晶格系统的存在，并证明了这些系统相对于纠缠自由度表现出对称性破缺。这是第一个启示：如果用底层 PEPS 表示的对称性来描述多体物态，Landau 范式便可以得到拯救——不同的相在根本不同的方式下实现张量的对称性。
\newpage\section{第 5 讲：广义对称性（Generalised symmetries）}\label{sec:GenSym}
\subsection{动机（Motivation）}
在第~\ref{sec:PhasesOfMatter}讲中，我们对具有群对称性的一维量子点阵模型的相图进行了分类。我们还看到，在一维玻色子系统中，对称性的存在对于产生非平凡的相图是\emph{必要}的。回顾 \cref{sec:Fermions}，对于费米子的情况，由于超选择定则（superselection rule）的存在，费米子系统可以表现出本征拓扑序（topological order），其情形更为微妙。近年来，学术界投入了大量的研究精力来理解超出群对称性范畴的对称性系统~\cite{Gaiotto2014,Thorngren2019,Aasen2020,Komargodski2020,McGreevy2022,Moradi2022,Freed2022,Shao2023,Schafer-Nameki2023,Moradi2023,Bhardwaj2023}。正如我们在第~\ref{sec:PhasesOfMatter}讲引言中所预示的，由融合范畴（fusion category）描述的广义对称性，作为非手征二维拓扑序的\emph{虚拟对称性}出现，因而也表现为其边界哈密顿量的对称性，或者等价地说，表现为描述其边缘模的有效理论的对称性。研究这些广义对称性如何影响系统的相图催生了\emph{广义 Landau 范式}，在此范式下，不同的物态被发现与打破广义对称性的不同方式相对应~\cite{Thorngren2019,Moradi2022,Bhardwaj2023,Lootens2024,Chen2025}。对于非群对称性，我们不能简单地说对称性破缺到了某个子群，因此我们需要探讨：一个系统如何（部分地）破缺一个广义对称性？在展现不同广义对称性破缺方式的系统之间是否存在映射？本讲将对这些问题进行深入阐述。

为什么张量网络是描述此类范畴对称性的理想工具？一个关键的观察在于，为了编码超出局域单点（on-site）形式的对称性，我们需要在相邻自由度上以关联方式起作用的算符。此外，我们希望能够在任意尺寸的系统上定义这些算符，因此我们需要局部地编码这些关联对称性。MPO 正是用于此目的的完美语言。它们由单个局域张量参数化，可以在任意数量的格点上定义，并通过其虚拟自由度在相邻物理自由度之间传递信息。基于这些考虑，我们常常交替使用\emph{广义对称性}、\emph{范畴对称性}与 \emph{MPO 对称性}这几个术语。

事实上，对一般 MPO 对称性的研究表明，在单射性（进一步在 \cref{sec:MPORepresentationTheory} 中讨论）等温和条件下，其底层结构等价于融合范畴。简言之，融合范畴是由（简单）对象构成的半单结合代数结构，非常类似于群是其元素（其简单对象）的结合乘法规则。然而，与群的情形不同的是，融合范畴的简单对象的乘法可能导致简单对象的直和。一个典型范例是融合范畴 $\Rep(G)$，其简单对象是群的不可约表示，其乘法规则由不可约表示张量积的直和分解决定，例如 $1/2\otimes 1/2=0\oplus 1$。融合范畴理论在数学与数学物理中是研究得非常透彻的课题，而 MPO 恰好构成了这些融合范畴在点阵上\emph{表示}的基石，正如矩阵在自旋系统中表示单点群对称性一样。

在最后一讲中，我们将首先研究范畴对称性的 MPO 表示论，并展示\emph{模范畴}（module category）如何标记不同的表示。随后我们将给出三个具体的应用。首先，我们将探讨这些 MPO 对称性如何在 Levin-Wen 弦网（string-net）模型中产生，该模型为所有非手征二维拓扑序提供了微观实现。然后我们提供两种方法来构造这些弦网模型的\emph{边缘理论}。在 \cref{sec:SC} 中，我们将展示这些 PEPS 表示的所谓\emph{奇异关联函数}（strange correlator）如何导出临界模型的二维\emph{经典}配分函数的构造~\cite{Vanhove2018}，从而推广了 Aasen、Mong 和 Fendley 的构造~\cite{Aasen2016}。这种奇异关联函数构造的主要优点在于，它清晰地揭示了临界点阵模型的\emph{拓扑}特征与\emph{几何}或\emph{动力学}特征之间的区别：(2+0) 维配分函数继承了来自 (2+1) 维弦网的 MPO 对称性，因而提供了 TFT/CFT 对应在点阵上的对应物。这种构建具有范畴对称性的边缘理论的全息方法，可以理解为 Kong、Wen 和 Zheng 的范畴构造~\cite{Kong2015}在点阵上的显式实现，它为现在广泛称为 \emph{SymTFT} 或\emph{拓扑全息}（topological holography）的理论提供了坚实的基础~\cite{Apruzzi2021,Chatterjee2022,Freed2022,Delcamp2024,Chen2025}。最后，我们将提供一种系统的方法来构造局域 (1+1) 维 MPO 对称哈密顿量，以及针对广义对称性的广义对偶性（duality）或 gauge 过程。最后，我们讨论范畴对称性的广义 Landau 范式及其若干物理推论。

\subsection{MPO 表示论（MPO representation theory）}\label{sec:MPORepresentationTheory}
在本节中，我们将遵循与 \cref{sec:YangBaxterFromMPO} 类似的逻辑。假设有一组表示融合代数的单射 MPO，我们推导出融合张量满足的自洽方程。随后我们展示如何利用这些方程的解来构造这些 MPO 和融合张量的表示。我们将说明如何从描述对称性的融合范畴上的\emph{模范畴}所编码的数据中构造出不同的表示。

\bigskip\noindent 具体而言，考虑一组标记为 $a,b,c,\dots$ 的有限单射 MPO，并施加周期性边界条件。我们感兴趣的是结构常数 $N_{ab}^c$ 与周期链格点数无关的 MPO 代数：
\begin{equation}
\label{eq:FusionRules}
    \inputtikz{MPOsStacked} = \sum_{c} N_{ab}^{c} \inputtikz{MPOc}\,\, .
\end{equation}
在此处及下文中，求和遍历代数中出现的所有不同 MPO。其中存在一个标记为 $\mbb 1$ 的单位元，满足 $N_{a\mbb 1}^{a'}=N_{\mbb 1 a}^{a'}=\delta_{a,a'}$。非负整数张量 $N$ 定义了一个\emph{含幺结合融合}结构，特别意味着结构常数满足 $\sum_d N_{ab}^dN_{dc}^e = \sum_d N_{bc}^d N_{ad}^e$。利用基本定理~\cref{eq:FundTh}，我们发现通过将以下\emph{融合张量}穿过 MPO 张量，可以在局域上实现 \cref{eq:FusionRules}：
\begin{equation}
\label{eq:IntertwinerHorizontal}
    \inputtikz{MPOIntertwiner_1} \,= \,\,
    \raisebox{4.5pt}{\inputtikz{MPOIntertwiner_2}}\,\,,
\end{equation}
其中重数 $i=1,2,\dots,N_{ab}^c$ 标记线性无关的融合张量。类似于 \cref{sec:SPT2d}，MPO 乘法的结合律意味着存在一组称为 \emph{$F$-符号}或\emph{结合子}的\emph{幺正}矩阵，它们编码基底变换
\begin{equation}\label{eq:MPOPentagon}
    \raisebox{-5.5pt}{\inputtikz{MPOPentagon_1}} = \sum_{f,k,l}\big(F^{abc}_{d}\big)_{e,ij}^{f,kl}\,\,\raisebox{-5.5pt}{\inputtikz{MPOPentagon_2}}.
\end{equation}
要求四个 MPO 的乘积（可以以两种不同方式重新耦合）具有结合性，揭示了这些 $F$-符号服从一个通称为\emph{五边形方程}（pentagon equation）的自洽条件。示意形式上，五边形方程是如下形式的多元线性方程组：
\begin{equation}\label{eq:pentagon}
    FF = \sum FFF,
\end{equation}
或者展开写成分量形式为：
\begin{equation}\label{eq:pent}
    \sum_o \left(F^{fcd}_e\right)^{h,no}_{g,lm} \left(F^{abh}_e\right)^{i,pq}_{f,ko} = \sum_{j,rst} \left(F^{abc}_g\right)^{j,rs}_{f,kl} \left(F^{ajd}_e\right)^{i,tq}_{g,sm} \left(F^{bcd}_i\right)^{h,np}_{j,rt}.
\end{equation}
本质上，由 $N$-符号指定的融合规则连同五边形方程的幺正解构成了\emph{幺正融合范畴} $\mc C$ 中编码的数据。\footnote{注意到五边形方程与通过类似推理应用于群对称性而在 \cref{eq:3cocycle} 中导出的 3-上循环极为相似。事实上，一个有限群连同一个 3-上循环的选择定义了一个称为 $\Vect_G^\omega$ 的融合范畴（如 \cref{sec:cat} 所综述），其 $F$-符号恰与 $\omega$ 重合。}重要的是，对于任意合法的融合结构，在重数指标的基变换意义下，五边形方程的解的数量是有限的。

给定五边形方程的一个解，我们如何找到对应于 \cref{eq:FusionRules} 意义下的 MPO 表示？完全契合 \cref{sec:YangBaxterFromMPO} 的逻辑，我们可以通过将 $F$-符号等同于 \cref{eq:IntertwinerHorizontal} 中定义的融合张量（的分量）以及 MPO 张量来找到表示。在此情况下，\cref{eq:MPOPentagon} 与 \cref{eq:IntertwinerHorizontal} 均等价于五边形方程 \eqref{eq:pent}~\cite{Sahinoglu2014,Bultinck2015,Williamson2017,Molnar2022}。具体来说，非零分量的分配如下\footnote{我们略去了涉及量子维数的常规预因子~\cite{Bultinck2017,Lootens2020}。}：
\begin{equation}\label{eq:SingleLineLabels}
    \inputtikz{PEPS_C_single_labels} = \big(F^{abc}_d\big)_{e,ij}^{f,kl}= 
    \inputtikz{MPO_single_labels}\,.
\end{equation}
请注意，每个指标都由一个对象三元组 $a,b,c$ 和取值为 $1,2,\dots,N_{ab}^c$ 的重数指标标记。特别地，为了使分量非零，所有三元组必须同时满足融合规则。此外，这些张量的结构使得一些对象在各指标之间共享。张量的这种内在结构启示我们将其描绘为\emph{三线张量}（triple line tensor）：
\begin{equation}
    \inputtikz{PEPS_C_triple}\equiv\inputtikz{PEPS_C_single_labels} ,\quad
    \inputtikz{MPO_single_labels}\equiv\inputtikz{MPO_triple}\,.
\end{equation}
这些张量的缩并是通过识别串联张量链上的简单对象、配对指标并对其求和来进行的，详见文献~\cite{Bultinck2017,Lootens2020}。

由于融合规则的存在，MPO 所作用的希尔伯特空间通常并非简单的张量积空间，而是由融合规则所施加的具有局域\emph{约束}的希尔伯特空间。在这种意义下，此处构造并在下文推广的 MPO 构成了在\emph{任意子链}（anyonic chain）背景下所考虑的拓扑对称性的一种表示和推广~\cite{Feiguin2006,Trebst2008,Gils2013,Buican2017,Huang2021,Jones2024}。例如，原始文献~\cite{Feiguin2006}中考虑的模型就是由 Fibonacci 融合范畴 $\Fib$ 构造的（见 \cref{sec:cat}）。

\bigskip\noindent 上述讨论仅提供了 $\mc C$ 的一种 MPO 表示。正如文献~\cite{Lootens2020}中所展示的，给定融合范畴 $\mc C$ 上的任意\emph{模范畴}都提供了在点阵上实现该对称性的一种不同方式。简言之，一个 $\mc C$-模范畴 $\mc R$ 包含简单对象 $A,B,\dots$ 的有限集合，以及 $\mc C$ 在 $\mc R$ 上的\emph{作用}，可表示为 $a\act A=\sum_{B} N_{aA}^{B}B$。值得注意的是，该 $N$-符号不同于 $\mc C$ 的 $N$-符号，因为它依赖于 $\mc R$ 的两个对象和 $\mc C$ 的一个对象。模范畴还提供了一组幺正的 \emph{$\mF$-符号}或\emph{（左）模结合子}，它们与 $\mc C$ 的 $F$-符号一起服从如下形式的\emph{混合}五边形方程：
\begin{equation}\label{eq:ModulePentagon}
    \mF\mF = \sum F\mF\mF.
\end{equation}
给定一个 $\mc C$-模范畴及其对应的模结合子，我们可以通过将其分量等同于 $\mF$-符号来构造融合张量——完全遵循 \cref{eq:SingleLineLabels} 并如下文 \cref{eq:TripleLineTensors} 所详述——从而借助 \cref{eq:ModulePentagon} 使得穿透条件 \cref{eq:IntertwinerHorizontal} 得到满足。值得注意的是，每个融合范畴 $\mc C$ 都在其自身上定义了一个模范畴，称为 \emph{regular 模范畴}。在此情况下，我们找回前一段落中的 MPO 表示。在解释 MPO 张量的求值之前，让我们通过一个简单例子来建立对该框架的直觉。

\paragraph{示例：$\mc C = \Vect_G$} 形式上，一个（无反常）有限群对称性由 $G$-分次向量空间的融合范畴 $\Vect_G$ 描述。简单对象对应于群元素，融合规则由群乘法给出，即 $N_{g,h}^k = \delta_{gh,k}$。$\Vect_G$ 上的模范畴与对偶对 $(H,\psi)$ 一一对应，其中 $H\subseteq G$ 是在共轭意义下的子群，$\psi$ 是子群 $H$ 的取值于 $\rU(1)$ 的 2-上循环。细节请参阅 \cref{sec:cat}。让我们聚焦于 $H=G$ 且 $\psi$ 为平凡的情形。对应的模范畴记作 $\Vect$，导出如下 MPO：
\begin{equation}\label{eq:MPOonsite}
    \inputtikz{MPOGroupg} \,\,= \bigotimes_\msf i u_\msf i(g).
\end{equation}
这里的 $u_\msf i$ 是 $G$ 的（不一定不可约的）幺正表示，并且可以在每个格点上独立选择。MPO 的群结构源自 $u_\msf i(g)u_\msf i(h)=u_\msf i(gh)$，$\forall g,h\in G$，$\forall \msf i$。虚线强调了 MPO 的键维为 $1$。值得注意的是，MPO 作用在张量积希尔伯特空间上这一事实，是选取模范畴 $\Vect$ 的直接结果。对于 $\Vect_G$ 上的其他允许模范畴，如上所述，MPO 通常作用在受约束的希尔伯特空间上。

\bigskip\noindent 注意在局域单点群对称性 \cref{eq:MPOonsite} 的情形下，人们可以把所有由群元素 $g\in G$ 索引的 MPO 张量 $u_\msf i(g)$ 同时化为块对角形式。各块对应于表示 $u_\msf i$ 中所包含的不可约表示。更一般地，我们可以将 MPO 张量的集合视为由块 $a,b,c,\dots$ 标记的算符值矩阵，并在\emph{垂直}方向上对由它们张成的代数进行块分解。关键的洞见在于，这些单射块的标签及其融合规则，即\footnote{注意这些 MPO 在垂直方向延伸——如图形中虚线所示——但在水平方向并不延伸。}
\begin{equation}
    \inputtikz{MPOsStackedVertical_1} = \sum_\gamma N_{\alpha\beta}^\gamma \,\, \inputtikz{MPOsStackedVertical_2}\,\,,
\end{equation}
由另一个融合范畴 $\mc D$ 所支配，该范畴通常\emph{不同于}对称性范畴 $\mc C$。它的简单对象将用希腊字母 $\alpha,\beta,\gamma,\dots$ 标记。给定 $\mc C$ 与 $\mc R$，$\mc D$ 是完全确定的。换句话说，一旦固定了对称性范畴以及通过所选模范畴 $\mc R$ 在点阵上表示它的方式，$\mc D$ 就被唯一确定。在技术上，$\mc D$ 对应于 $\mc C$ 关于 $\mc R$ 的 \emph{Morita 对偶}（亦见 \cref{sec:cat}），记作 $\mc D=\mc C_\mc R^\star$，且等价地有 $\mc C =\mc D^\star_\mc R$。$\mc C$ 与 $\mc D$ 的 Morita 等价特别保证了 $\mc C$ 与 $\mc D$ 的 \emph{Drinfeld 中心}是等价的：$\mc Z(\mc C)\simeq\mc Z(\mc D)$。在 \cref{sec:SN} 中我们将看到，$\mc Z(\mc D)$ 编码了关于弦网基态区与任意子激发的全部信息，其中弦网的拓扑虚拟对称性由 $\mc C$ 描述。在 \cref{sec:SC} 和 \cref{sec:dualities} 中，Drinfeld 中心的元素将标记弦网边界理论的\emph{拓扑扇区}或\emph{超选择区}。

作为一个融合范畴，$\mc D$ 带有其自身的一组 $F$-符号，我们记作 $\dF$。它们满足类似于 \cref{eq:pentagon} 的五边形方程。Morita 对偶 $\mc D$ 的构造方式使得 $\mc R$ 也是 $\mc D$ 上的模范畴，此时 $\mc D$ 从\emph{右侧}作用于 $\mc R$。因此再次存在对应的 $F$-符号，记作 $\nF$，用以编码 $\mc D$ 在 $\mc R$ 上的作用的结合性。这些符号满足我们遇到的第四个五边形方程，现在涉及 $\mc D$ 的结合子：
\begin{equation}\label{eq:Pent2}
    \nF\nF = \sum \dF\nF\nF.
\end{equation}
除了同时作为 $\mc C$ 和 $\mc D$ 的模范畴之外，$\mc R$ 还是一个（可逆的）\emph{$(\mc C,\mc D)$-双模范畴}（见 \cref{sec:cat}）。与其双模结构相关联的是最后一组记作 $\rF$ 的 $F$-符号，服从
\begin{align}
    \rF\mF &= \sum\mF\rF\rF, \label{eq:Pent3}\\
    \nF\rF &= \sum\rF\rF\nF \label{eq:Pent4}.
\end{align}
MPO 张量的分量等于 $\rF$-符号，而 $\nF$-符号恰好是\emph{垂直}方向融合张量所求得的值。汇总如下：
\begingroup
\allowdisplaybreaks
\begin{align}
    \inputtikz{PEPS_C_triple_mod}&\equiv\inputtikz{PEPS_C_single_labels_mod} = \big(\mF^{abC}_A\big)_{c,ij}^{B,kl}, \notag\\
    \inputtikz{PEPS_D_triple}&\equiv\inputtikz{PEPS_D_single_labels}=\big(\nF^{A\alpha\beta}_B\big)_{C,ij}^{\gamma,kl},\label{eq:TripleLineTensors}\\
    \inputtikz{MPO_triple_mod}&\equiv\inputtikz{MPO_single_labels_mod}=\big(\rF^{aA\alpha}_D\big)_{C,ij}^{B,kl}.\notag
\end{align}
\endgroup
注意在这些三线张量中，模范畴中的对象始终位于外侧的\emph{环圈}上。然后可以验证垂直穿透条件
\begin{equation}\label{eq:IntertwinerVertical}
    \inputtikz{MPOIntertwiner_3} =
    \inputtikz{MPOIntertwiner_4}
\end{equation}
借助 \cref{eq:Pent4} 得到满足，而 \cref{eq:IntertwinerHorizontal} 等价于 \cref{eq:Pent3}。这些融合张量进而通过图 \cref{eq:MPOPentagon} 的旋转版本关联起来，该图涉及这些融合张量和 $\dF$-符号，等价于 \cref{eq:Pent2}。

在上面考虑的示例 $\mc C=\Vect_G$、$\mc R=\Vect$ 中，融合范畴 $\mc D$ 为 $\Rep(G)$，其简单对象是群 $G$ 的不可约表示。对应的 $\dF$-符号等于 $G$ 的 6j-符号，而 $\nF$-符号等于 \emph{Clebsch-Gordan 系数}。这些 Clebsch-Gordan 系数将 MPO 张量 $u_\msf i(g)$ 分解为其不可约块，而这些块中的数据编码在 $\rF$ 中，其值由 $G$ 的不可约表示给出。在一般（非群）情形下，$\nF$-符号在下文将被解释为此处构造的 MPO 对称性的\emph{广义} Clebsch-Gordan 系数。与此同时，我们将把 \cref{eq:IntertwinerVertical} 解释为\emph{广义 Wigner-Eckart 定理}，表明 MPO 对称张量是由这些广义 Clebsch-Gordan 系数构建的。

\bigskip\noindent 综上所述，展现出如下完整图像。给定由融合范畴 $\mc C$ 指定的对称性，可通过可逆 $(\mc C,\mc D)$-双模范畴 $\mc R$ 构造其 MPO 表示。该结构以五个 F-符号 $(F,\mF,\rF,\nF,\dF)$ 的形式编码结合子数据，它们满足六个五边形方程。双模结合子 $\rF$ 可用于构造融合范畴 $\mc C$ 的单射 MPO 表示，其张量在局域上通过分别求值为 $\mF$ 和 $\nF$ 的张量在水平方向 \cref{eq:IntertwinerHorizontal} 和垂直方向 \cref{eq:IntertwinerVertical} 融合。两者分别满足涉及 $F$ 和 $\dF$ 的重新耦合恒等式。

\subsection{应用 1：弦网基态与交织算符（string-net ground states \& intertwiners）}\label{sec:SN}
\emph{弦网（string-net）模型}最早由 Levin 和 Wen 在文献~\cite{Levin2004}中构造。它们构成了一类重整化群（RG）不动点点阵模型，其输入为一个（球形）融合范畴 $\mathcal{D}$，并实现由其 Drinfeld 中心 $\mc Z(\mc D)$ 给出的非手征拓扑序~\cite{Kirillov2011,Hu2015,Kawagoe2024}。这意味着，其基态与有能隙激发重现了编码在 $\mc Z(\mc D)$ 中的拓扑区和任意子统计。\footnote{等价地，它们可以被视为 Turaev--Viro--Barrett--Westbury 拓扑场论的哈密顿量实现~\cite{Turaev1992,Barrett1993,Kirillov2011}。}在 \cref{sec:QuantumDoubles} 中遇到的量子双重模型在选取 $\mc D=\Vect_G$ 时被重现。名称\emph{弦网}源于该模型描述有色弦的凝聚相这一解释，其中允许的弦类型及其相互作用编码在 $\mc D$ 中。哈密顿量是每个面元（plaquette）上的局域对易投影算符之和，从而可以直接导出其基态的 PEPS 表示。

正如文献~\cite{Lootens2020}所示，对于输入 $\mc D$ 上的每一个模范畴 $\mc R$，都存在对应的基态 PEPS 表示。这一观察推广并统一了早期文献中隐含假定 regular 模范畴的情形~\cite{Verstraete2006b,Buerschaper2009,Gu2009,Williamson2017,Bultinck2017}，以及在量子双重模型的具体情形中选取 $\Vect$ 作为 $\Vect_G$ 上的模范畴的情形~\cite{Schuch2010}。从这个角度来看，弦网因而是一个具有特定类型\emph{虚拟}对称性的 PEPS。构成底层 Morita 对偶融合范畴 $\mathcal{C}$ 表示的 MPO 是其对称性，并且它们纯粹作用在张量网络的虚拟自由度上。根据上面讨论的广义 Wigner-Eckart 定理，PEPS 张量如 \cref{eq:TripleLineTensors} 中所示求值为 $\nF$，其物理层指标取值为 $(\alpha\beta\gamma,k)$。

类似于 \cref{sec:QuantumDoubles}，不同的环面基态区和任意子激发（均由 $\mc Z(\mc D)$ 中的元素标记）可以通过\emph{管代数}（tube algebra）来理解。管代数是由 $\mc C$ 中的 MPO 构造的，其中心幂等元与 $\mc Z(\mc C)$ 中的元素一一对应~\cite{Williamson2017,Bultinck2017}。回顾由 $\mc R$ 的可逆性保证了 $\mc Z(\mc D)$ 与 $\mc Z(\mc C)$ 等价。事实上，如果我们考虑环面上的弦网，可以通过将上一节的 MPO 对称性缠绕在其不可收缩环路周围来获得不同的拓扑区。在它们的交叉点处，放置一个张量，其分量由中心幂等元决定。在弦网 PEPS 中插入与对象 $Z\in\mc Z(\mc C)$ 相关联的中心幂等元可以描绘为：
\begin{equation}
    \sum_{\substack{a_\alpha,b_\beta,c_\gamma\\i,j}} (\mc P_Z)_{a_\alpha,b_\beta,c_\gamma}^{i,j} \quad \inputtikz{String-net_PEPS} \,\, .
\end{equation}
请注意，紫色张量与 \cref{eq:IntertwinerVertical} 的融合张量重合，因此中心幂等元的具体位置是无法检测到的。在此，$a_\alpha$ 代表 $a_\alpha\equiv(a,\alpha)$，其中 $a$ 是 $\mc C$ 中的简单对象，$\alpha$ 标记 $a$ 在 $Z$ 中的重数。

给定基于 $\mc R$ 的 PEPS 表示，可以通过一类实现模范畴变换的新 MPO \emph{交织算符}（intertwiner），在局域上将其转换为基于 $\mc R'$ 的表示。特别地，当 $\mc R'=\mc D$ 时，这些 MPO 交织算符由 $\mc R$ 中的对象标记，并由模结合子构造而成\footnote{\label{fn}在一般情况下，需要模范畴之间的\emph{（$\mc D$-）模函子}概念，它们组成记作 $\Fun_\mc D(\mc R,\mc R')$ 的范畴。值得注意的是，$\Fun_\mc D(\mc R,\mc R')$ 可以被赋予可逆 $(\mc D_\mc R^\star,\mc D_{\mc R'}^\star)$-双模范畴的结构。}。模范畴的改变暗示了这些不同 PEPS 表示的边缘哈密顿量之间的等价性。稍后在 \cref{sec:dualities} 中，我们将看到对应于不同表示 $\mc R,\mc R'$ 的边缘哈密顿量通过\emph{对偶变换}联系起来，或者等价地说，通过对其（部分）MPO 对称性实施 gauging 来联系。上述 MPO 交织算符充当了这些对偶模型之间的交织算符角色，并且在充分考虑边界条件时，可以被提升为等距算符（isometry）。

此外，张量网络图像还可以用于显式构造弦网模型的\emph{有能隙边界}。这超出了本讲义的范围，我们建议读者参阅文献~\cite{Lootens2020}了解细节，但在此提一下，在该情形下模范畴同样作为标记不同边界条件的工具出现。总之，由 Kitaev 和 Kong 在文献~\cite{Kitaev2011}中发展的关于弦网畴壁与边界的完整图像，可以重新用张量网络语言表述，使得这些结构非常适合在计算机上进行模拟。

\subsection{应用 2：经典配分函数作为奇异关联函数（strange correlators）}\label{sec:SC}
掌握了弦网基态的张量网络表示之后，我们现在可以讨论一般 MPO 形式体系的第二个应用，即所谓的\emph{奇异关联函数}（strange correlator）构造。该构造最初出现在文献~\cite{You2013}中，作为一种通过考察矩阵元 $C(\vec r,\vec r') = \mel{\Omega}{\phi(\vec r)\phi(\vec r')}{\Psi}$ 来探测非平凡 SPT 相的诊断手段，其中 $\ket{\Psi}$ 是待考察的状态，$\phi$ 是局域可观测量，而 $\ket{\Omega}$ 表示精心挑选的短程纠缠平凡态。对于表现出 SPT 序的许多一维和二维态 $\ket{\Psi}$，已经证明量 $C(\vec r,\vec r')$ 要么随距离 $|\vec r -\vec r'|$ 代数衰减，要么收敛到一个常数。因此，该方法提供了根据体基态波函数识别 SPT 序的方法。

这种构造的一个自然推广在于选择表现出本征拓扑序的状态 $\ket{\Psi}$。聚焦于二维非手征拓扑序的情形，我们因而可以通过 \cref{sec:SN} 的弦网基态与适当状态 $\ket{\Omega}$ 的重叠来构造奇异关联函数。这种方法由文献~\cite{Vanhove2018}开创。正如我们在上面所见，弦网基态的 PEPS 表示具有纯粹作用于虚拟能级的精确 MPO 对称性。因此，由这些状态构造的奇异关联函数继承了这些对称性，而与所选状态 $\ket{\Omega}$ 无关。然而，对于许多有趣的范例，$\ket{\Omega}$ 可以简单地选择为直积态。这样，很自然地将这种二维张量网络解释为带有拓扑缺陷线/对称性的\emph{经典统计模型}\footnote{然而请注意，取决于所选的 $\ket{\Omega}$，经典模型可能具有非正的 Boltzmann 权重，要求这些权重为正通常需要对状态 $\ket{\Omega}$ 进行微调。}的\emph{配分函数}。示意形式上，在 $\ket{\Omega}$ 为（均匀）直积态的情形下，配分函数采取如下形式：
\begin{equation}
    Z(\{\beta_i\}) := \braket{\Omega[\{\beta_i\}]}{\Psi}= \,\inputtikz{SC_1}= \,\inputtikz{SC_2}\,\, ,
\end{equation}
其中 $\{\beta_i\}_i$ 统指定义状态 $\ket{\Omega}$ 的参数，从而间接决定经典模型的 Boltzmann 权重。对于参数 $\{\beta_i\}_i$ 的适当取值，该配分函数实际上可以描述一个临界模型。在这种情况下，MPO 对称性可以解释为 CFT 连续拓扑缺陷在点阵上的遗存。因此，奇异关联函数构造提供了共形场论与拓扑场论之间对应关系的显式点阵表示，这种对应最初由 Witten~\cite{Witten1988}、Moore 和 Seiberg~\cite{Moore1988}构想，随后由 Fr\"ohlich、Fuchs、Runkel 和 Schweigert 在文献~\cite{Frohlich2004,Frohlich2006,Fuchs2002,Fuchs2003,Fuchs2004}中进一步发展。

\paragraph{示例：硬六边形模型} 也许最简单的示例是基于六角点阵上的 Fibonacci 融合范畴的弦网模型。该融合范畴包含两个记作 $\{\mbb 1,\tau\}$ 的简单对象，唯一的非平凡融合规则为 $\tau\otimes\tau=\mbb 1\oplus\tau$。我们将考虑的奇异关联函数是一个直积态，它将弦网的所有物理自由度投影到对象 $\tau$ 上——或者更确切地说，在 \cref{eq:TripleLineTensors} 的记号中，投影到 $\phi\cdot(\tau\tau\tau, 1)$，其中 $\phi=\frac{1+\sqrt 5}{2}$ 是黄金比例。此时环圈上的自由度扮演了经典二元变量的角色，并受制于相邻环圈不能同时处于 $\mbb 1$ 态的约束。很自然地将该模型解释为面元上的经典粒子气体，使得标记为 $\mbb 1$ 的环圈表示相应面元上存在粒子，而 $\tau$ 表示其不存在。这一所谓的\emph{硬六边形}（hard hexagon）模型在 1980 年由 Baxter 极其出色地精确求解，并被证明展现出两个相态~\cite{Baxter1980}。对于较大的逸度值，该模型表现出有序相，其特征在于完全占据六角点阵面元的三个子点阵之一；而对于低逸度，气体是稀疏且均匀的。这些相态被逸度的一个临界值 $z_c$ 处的二阶相变所分隔，其标度极限由 $c=4/5$ 的极小 Potts 模型描述~\cite{DiFrancesco1997}。

再次聚焦于奇异关联函数，构造配分函数的张量具有以下非零分量：
\begin{equation}
    \inputtikz{SCFib_1} = -1, \qquad \inputtikz{SCFib_2} = \phi^{5/6}.
\end{equation}
手头有了上述张量，随后很容易推导出配分函数可以写为
\begin{equation}
    Z = \sum_{n=0}^{N/3} z_c^n \ g(n,N),
\end{equation}
其中 $N$ 表示面元的总数，$g(n,N)$ 是组合因子，等于在服从上述约束的前提下在 $N$ 个面元中分布 $n$ 个粒子的方式总数~\cite{Fidkowski2009}。简单的代数计算表明 $z_c=\phi^{5}$，这恰好是硬六边形模型的临界逸度。因此，这个简单的奇异关联函数立即给出了临界逸度，而无需诉诸 Baxter 所采用的更为高深的方法。

从转移矩阵中还可以提取有限尺寸 CFT 能谱。上文我们已经强调了管代数在获取量子双重与弦网模型的拓扑超选择区方面的重要作用。如文献~\cite{Petkova2000,Aasen2020,Vanhove2018}所详述，奇异关联函数的转移矩阵可以补充由对称融合范畴的对象所标记的对称扭曲边界条件。这些边界条件在拓扑意义下保持点阵平移对称性（在差一个幺正变换的意义下）。在它们存在的情况下，管代数可用于将转移矩阵的能谱分解为各个拓扑区。对于 Fibonacci 的情形，存在 4 个标记为 $0,\tau,\bar\tau,\tau\bar\tau$ 的区。其对应的能谱展示在文献~\cite{Vanhove2018}中，并在其中解释了与 Potts 极小模型的初级场的对应关系。

\bigskip\noindent 在前一节中，我们论证了弦网基态的不同张量网络表示对应于输入融合范畴上的模范畴。巧合的是，对于本节的案例研究，我们并没有这种自由度。事实证明，模范畴的选择指定了环面上 CFT 的\emph{模不变性}（modular invariant）以及其在更高亏格曲面上的配分函数~\cite{Moore1988,DiFrancesco1997,Fuchs2002}。所谓的\emph{对角情形}对应于选择 regular 模范畴。改变该模范畴则被称为 \emph{orbifolding} 或 \emph{gauging}，在文献~\cite{Vanhove2021a}中进行了数值研究。

\bigskip\noindent 关于使用奇异关联函数形式体系的临界配分函数的张量网络表示，存在大量文献。我们建议读者参阅文献~\cite{Freedman2011,Vanhove2018,Lootens2019,Brehm2021,Vanhove2021a,Vanhove2021b,Zeng2022,Ji2024,Shen2025,Hung2025}以获取更多案例研究、部分解析结果以及开边界条件和不可定向流形上的能谱。此外，在最近的工作中，提出了完全从纯拓扑数据构造临界（可积）奇异关联函数的系统方法~\cite{Fendley2020,Hung2025}。其吸引力是显而易见的：奇异关联函数形式体系提供了直接构造具有精确范畴缺陷的临界点阵系统的方法。拓扑部分（编码在弦网中）与几何部分（编码在重叠态 $\ket{\Omega}$ 中）完全分离，从而使理论中的异常变得非常清晰。核心挑战之一是构造对应于新 CFT 的临界点阵模型——Haagerup 临界点阵模型~\cite{Huang2021,Vanhove2021b,Hung2025}是这种理论的第一个可能范例。

\subsection{应用 3：对偶性（dualities）}\label{sec:dualities}
在上一节中，我们构造了具有 \cref{sec:MPORepresentationTheory} 中分类和构造的 MPO 对称性的二维经典配分函数。在最后一个应用中，我们将构造与这些对称性对易的 (1+1) 维哈密顿量，并研究其对偶性。在某种意义上，这些模型可以被视为 \cref{sec:SN} 中弦网的\emph{边缘哈密顿量}。

\bigskip\noindent 对于群的局域单点表示情形，\emph{Clebsch-Gordan 系数}是对称张量的构建基石。事实上，\emph{Wigner-Eckart 定理}指出，对称张量可以分解为 Clebsch-Gordan 系数与\emph{约化矩阵元}的乘积。值得注意的是，后者仅依赖于不可约表示标签，而基态的耦合则完全由 Clebsch-Gordan 系数刻画~\cite{Eckart1930,Wigner1960}。这一结果不仅在微扰论和光谱学（通常在这些背景下引入它们）中处于核心地位，而且构成了对称张量网络算法的基石~\cite{Sanz2009,Singh2009,Weichselbaum2012,Singh2013,Lootens2024,Devos2025}。受到这一视角的启发，对于 MPO 对称性 $\mc C$ 的更一般情形，我们可以从 \cref{eq:IntertwinerVertical} 推导出一般的局域算符由融合张量构造如下：
\begin{equation}\label{eq:LocalOp}
    h_{\msf i,\msf i+1}^\mc R = h(\alpha,\alpha',\beta,\beta',\gamma,i,i')\inputtikz{MPOLocalOp}.
\end{equation}
在此，$h$ 是可以自由选取的复数，且可能依赖于所指出的标签。重要的是，这些系数与张量的内部虚拟指标无关，因此可以解释为约化矩阵元的推广。这些算符与 MPO 对称性的对易性源自
\begin{equation}
    \inputtikz{MPOLocalOp_1} =
    \inputtikz{MPOLocalOp_2} = 
    \inputtikz{MPOLocalOp_3} \,\, ,
\end{equation}
该式对 $\mc D$ 和 $\mc C$ 的所有指出简单对象均成立。由 \cref{eq:LocalOp} 出发，局域对称哈密顿量可构造为 $H=\sum_{\msf i} h_{\msf i,\msf i+1}^\mc R$。注意我们在此限制于两体算符，但通过取~\cref{eq:LocalOp} 的乘积同样可以容易地获得作用在更多格点上的对称算符。因此，局域算符的构造 \cref{eq:LocalOp} 将传统的 Wigner-Eckart 定理推广到了 MPO 对称性~\cite{Bridgeman2022}。

到目前为止，我们隐式假定了无限长链。完全按照 \cref{sec:SC} 的思路，我们也可以在闭合周期边界条件和对称扭曲边界条件下定义这些哈密顿量~\cite{Petkova2000,Aasen2020,Lootens2022}。对于后者，在两个格点之间的键上引入了一个额外的自由度，并且穿过该键的局域哈密顿量项相应地被修改，详情参阅文献~\cite{Lootens2022}。包含对称扭曲边界条件的重要性将在下文变得明朗。

\bigskip\noindent 为了说明表达式 \cref{eq:LocalOp}，考虑上文 \cref{eq:Heisenberg} 中遇到的自旋 1/2 反铁磁 Heisenberg 模型。通过调用（传统的）Wigner-Eckart 定理，局域哈密顿量项 $\vec{S}_\msf i \cdot \vec{S}_{\msf i+1}$ 可以从 ${\rm SU}(2)$ Clebsch-Gordan 系数构建。为了在当前框架中重新表述这一点，我们需要选择 $\mc D=\Rep({\rm SU}(2))$（${\rm SU}(2)$ 的表示范畴），并取 $\mc R=\Vect$。\footnote{严格来说，$\Rep({\rm SU}(2))$ 是张量范畴，而不是要求具有有限个简单对象的融合范畴。偏好更多数学严谨性的读者可以考虑将此形式体系应用于 ${\rm SU}(2)$ 的有限子群。}显式展开：
\begin{equation}
    \vec{S}_\msf i \cdot \vec{S}_{\msf i+1} = -\inputtikz{MPOLocalHamSU2} = \sum_{\substack{i,j\\i'\!,j'}} \CC{{\bf 2}}{{\bf 2}}{{\bf 0}}{i'}{j'}{1} \CC{{\bf 2}}{{\bf 2}}{{\bf 0}}{i}{j}{1} \,\, |i',j'\rangle\langle i,j|,
\end{equation}
相差一个加性常数。这里，${\bf 2}$ 和 ${\bf 0}$ 分别表示 ${\rm SU}(2)$ 的二重态和单态，而 $\CC{{\bf 2}}{{\bf 2}}{{\bf 0}}{\bullet}{\bullet}{1}$ 代表关联于
\begin{equation}
    \CC{{\bf 2}}{{\bf 2}}{{\bf 0}}{\bullet}{\bullet}{1} : {\bf 2} \otimes {\bf 2} \rightarrow {\bf 0}: \ket{i}\otimes\ket{j} = \CC{{\bf 2}}{{\bf 2}}{{\bf 0}}{i}{j}{1} \ket{1},
\end{equation}
的（实数）Clebsch-Gordan 系数，其中 $i,j,1$ 表示相应基础表示空间的基向量。因此，该哈密顿量项本质上将两个相邻的自旋 1/2 投影到单态表示上，而在张量积 $\frac{1}{2}\otimes\frac{1}{2}$ 中出现的自旋 1 表示被选定的参数 $h$ 赋予了零权重。

\bigskip\noindent 让我们考虑所有局域对称算符 \eqref{eq:LocalOp} 的集合，以及它们的有限乘积和线性组合。这构成了一个代数——有时被称为\emph{键代数}（bond algebra）~\cite{Cobanera2009,Cobanera2011,Moradi2022,Moradi2023,Jones2024}——由算符基 $\{h_A^\mc R\}$ 的结构常数决定：
\begin{equation}
    h_A^\mc R h_B^\mc R  = \sum_C f_{AB}^C(\dF) h_C^\mc R .
\end{equation}
至关重要的是，\emph{结构常数} $f_{AB}^C(\dF)$ 仅依赖于融合范畴 $\mathcal{D}$ 的结合子 $\dF$，也就是说，它们完全独立于模范畴 $\mathcal{R}$ 的选择。特别地，该观察意味着哈密顿量的能谱（在不计简并度的情况下）也与 $\mc R$ 无关。对于固定的 $\mc D$ 和 \cref{eq:LocalOp} 中系数 $h$ 的选择，改变模范畴 $\mc R\rightarrow\mc R'$ 归结为一次\emph{对偶变换}（duality transformation）。因此，$\mc D$ 上的不同模范畴（或等价地对称性范畴 $\mc C$）组织了一个庞大的对偶性网络。尽管这一框架看似抽象，但大多数具有物理意义的对偶性（包括 Kramers-Wannier 对偶、Kennedy-Tasaki 对偶和 Jordan-Wigner 变换）都自然地融入其中，正如后文进一步阐明的。在整篇讨论中，我们经常将感兴趣的（虚构）模型称为\emph{原始}模型，并将其与它的\emph{对偶}模型进行比较。对偶模型在以下几个方面可能看起来与原始模型大不相同。

如上所述，如果 $\mc R$ 拥有多于一个简单对象，局域对称算符 \cref{eq:LocalOp} 及其 MPO 对称性通常作用在受约束的希尔伯特空间上。因此，一个模型的希尔伯特空间通常与其任意对偶模型的希尔伯特空间完全不同，它们的维数甚至不必匹配。这与此处考虑的这类对偶性可以通过对原始 MPO 对称性的（部分）实施 gauging 来实现密切相关，从而受约束的希尔伯特空间作为（部分）所施加的 Gauss 约束被化解的\emph{有效}希尔伯特空间出现~\cite{Haegeman2014a,Bhardwaj2017,Seifnashri2025,Vancraeynest2025b}。因此，不同对偶模型的全局对称性通常也不同，因为它们分别编码在 $\mc D^\star_\mc R$ 和 $\mc D^\star_{\mc R'}$ 中。换句话说，给定初始对称性和对偶变换，对偶对称性完全被确定。模型之间存在对偶性还意味着一个理论中的每个可观测量都等价于对偶理论中的一个\emph{对偶}可观测量。然而，只有该理论中的\emph{对称}局域可观测量才被映射到对偶理论中的对偶局域对称算符，而带荷算符（例如局域序参量）则被映射到非局域的\emph{弦序参量}（string order parameter）。

张量网络处理对偶性的一大亮点在于，我们可以用（非可逆）矩阵乘积算符构建对偶模型之间显式的\emph{交织算符}（intertwiner）~\cite{Lootens2021b}。在哈密顿量项的层面上，它们的作用遵循：
\begin{equation}\label{eq:DualInter}
    h_{\msf i,\msf i+1}^{\mc R'} \circ D = D\circ h^{\mc R}_{\msf i,\msf i+1},
\end{equation}
由此可知，用 $D$ 作用在 $H^\mc R$ 的本征态上，将导致具有完全相同能量的 $H^{\mc R'}$ 的本征态。

正如文献~\cite{Lootens2021b}中所证明的，这样的交织算符仅依赖于范畴 $\mc D,\mc R$ 和 $\mc R'$，而不依赖于所研究的具体模型。事实上，这些 MPO 交织算符正是 \cref{sec:SN} 中提到的将与模范畴 $\mc R,\mc R'$ 相关联的 $\mc D$ 弦网模型的不同 PEPS 表示相互转换的 MPO 交织算符。

\cref{eq:DualInter} 的另一个推论是，给定一个交织算符 $D$，我们可以将 $D$ 从\emph{右侧}乘以初始对称性 $\mc C=\mc D^\star_\mc R$ 的对称算符，并从\emph{左侧}乘以 $\mc D^\star_\mc R$ 中的对偶对称算符。这样，不同的交织算符自身组织在原始对称性范畴与对偶对称性范畴之间的可逆双模范畴（我们将记作 $\mc M$）中。由于 $\mc D^\star_{\mc R'} =\mc C^\star_\mc M$，该构造是自洽的。\footnote{这源于 \cref{fn} 中的观察，即对偶算符由组织在 $\Fun_\mc D(\mc R,\mc R')$ 中的 $\mc D$-模函子标记。}

对于一对对偶模型，我们可以将这些结构汇集在如下三角形图中：
\begin{equation}
    \inputtikz{TriangleDiagram_1}.
\end{equation}
在此图中，每个顶点代表一个融合范畴。顶部的顶点——$\mc D$——编码由局域对称算符生成的键代数，而另外两个顶点——$\mc D_\mc R^\star$ 和 $\mc D_{\mc R'}^\star$——编码全局对称性。连接的边代表可逆双模范畴。对角线上的两条边对应于 $\mc R$ 和 $\mc R'$，编码任一模型的自由度。底部的边——$\mc M$——编码 MPO 交织算符。

$\mc R$ 与 $\mc R'$ 的可逆性特别保证了原始对称性 $\mc D_\mc R^\star$ 与对偶对称性 $\mc D_{\mc R'}^\star$ 的 Drinfeld 中心是等价的。在此背景下，Drinfeld 中心中的对象标记理论的\emph{拓扑区}。这些区对应于一个边界条件（破缺了部分对称性）连同幸存对称性的对称荷。通过仔细研究对偶算符与这些区之间的相互作用，非可逆对偶算符可以被提升为幺正算符~\cite{Lootens2022}。在对偶变换下，拓扑区通常发生置换，正如 Kramers-Wannier 对偶所示。

\paragraph{示例：Kramers-Wannier 对偶} 让我们展示该形式体系如何重现横场 Ising 模型~\cref{eq:TFIM} 著名的 Kramers-Wannier（KW）对偶~\cite{Kramers1941}。为此，考虑哈密顿量
\begin{equation}\label{eq:TFIM_per}
    H_{\rm TFIM}^{\Vect_{\mbb Z_2}} = -\sum_{\msf i=1}^{L} \sigma^Z_{\msf i}\sigma^Z_{\msf i+1} - \lambda\sum_{\msf i=1}^L \sigma^X_{\msf i},
\end{equation}
采用周期性边界条件，即 $L+1\equiv 1$。该哈密顿量显然与由 $\prod_\msf i \sigma^X_\msf i$ 生成的 $\mbb Z_2$ 对称性对易。键代数对应于融合范畴 $\mc D=\Vect_{\mbb Z_2}$，模范畴为 $\mc R=\Vect_{\mbb Z_2}$。Kramers-Wannier 对偶可以被看作是对 $\mbb Z_2$ 对称性实施 gauging，或者更抽象地说是改变模范畴 $\mc R=\Vect_{\mbb Z_2}\rightarrow\mc R'=\Vect$，对应于 $\mc M=\Vect$。（唯一的）对偶 MPO 可以描绘为：
\begin{equation}\label{eq:KW_MPO_per}
    D_{\rm KW} = \inputtikz{KW_MPO_per}\,\,,
\end{equation}
其中我们用虚线显式指明了周期性边界条件，在此处和下文中用{\color{BlueViolet}紫色}描绘对偶 MPO。注意两侧的物理腿平移了半个格点。我们再次强调，该对偶 MPO 对每个具有 $\mbb Z_2$ 对称性的模型都是相同的，因而并非特定于 Ising 模型。

MPO 张量是在 $\sigma^Z$ 和 $\sigma^X$ 本征基下的 GHZ 张量，显式读作：
\begin{equation}
    \raisebox{5pt}{\inputtikz{KW_tensor_1_a}} = \ket{000} + \ket{111},\qquad \raisebox{5pt}{\inputtikz{KW_tensor_2}} = \ket{+\!+\!+} + \ket{-\!-\!-}\,\,.
\end{equation}
交织算符的所有（全局）性质均可从这些局域张量的对称性质中推导出来。特别地，注意到第一个张量在任意两条腿上用 $\sigma^Z$ 作用，以及在所有腿上同时用 $\sigma^X$ 作用下是不变的：
\begin{equation}
\begin{gathered}
    \raisebox{5pt}{\inputtikz{KW_tensor_1_a}} = \raisebox{7pt}{\inputtikz{KW_tensor_1_c}} =\\  \inputtikz{KW_tensor_1_e} = \inputtikz{KW_tensor_1_d} = \inputtikz{KW_tensor_1_b}
\end{gathered}\,\,.
\end{equation}
第二个张量具有相同的对称性，只需互换 $\sigma^Z$ 和 $\sigma^X$。因此，该 MPO 按照如下规则映射局域哈密顿量项：
\begin{equation}\label{eq:KW_mapping}
    \sigma^Z_{\msf i}\sigma^Z_{\msf i+1} \rightarrow \sigma^Z_{\msf i+\frac{1}{2}}, \quad
    \sigma_\msf i^X \rightarrow \sigma^X_{\msf i-\frac{1}{2}}\sigma^X_{\msf i+\frac{1}{2}},
\end{equation}
从而对偶哈密顿量为
\begin{equation}\label{eq:TFIM_dual}
    H_{\rm TFIM}^{\Vect} = -\sum_{\msf i=1}^{L} \sigma^Z_{\msf i+\frac{1}{2}} - \lambda\sum_{\msf i=1}^L \sigma^X_{\msf i-\frac{1}{2}}\sigma^X_{\msf i+\frac{1}{2}}.
\end{equation}
该哈密顿量具有由 $\prod_\msf i \sigma^Z_{\msf i+1/2}$ 生成的对称性，对应于对称性范畴 $\Rep(\mbb Z_2)=(\Vect_{\mbb Z_2})_\Vect^\star$，符合该框架的预测。映射 \cref{eq:KW_mapping} 还表明原始模型中的序参量被映射到对偶模型中的无序参量，反之亦然：
\begin{equation}
    \sigma^Z_\msf i \sigma^Z_\msf j \rightarrow \prod_{\msf k=\msf i}^{\msf j-1}\sigma^Z_{\msf k+\frac{1}{2}},\quad  \prod_{\msf k=\msf i}^{\msf j} \sigma^X_\msf k \rightarrow \sigma^X_{\msf i-\frac{1}{2}} \sigma^X_{\msf j+\frac{1}{2}}.
\end{equation}

MPO \cref{eq:KW_MPO_per} 具有一个零空间，包含所有 $\mbb Z_2$ 奇宇称态：$\prod_\msf i\sigma_\msf i^X\ket{\Psi}=-\ket{\Psi}$。类似地，其像空间仅由在对偶对称性下为偶宇称的状态组成。为了在任一模型中访问非平凡的荷区，我们需要调整 MPO 以适应所有对称扭曲边界条件。对于当前情况，这归结为周期性或反周期性边界条件。在哈密顿量层面上，改变边界条件归结为翻转其中一个最近邻相互作用项的正负号，习惯上选取格点 $L$ 和 $1$ 之间的项。

在指定了原始模型的边界条件和荷区之后，对偶模型的拓扑区由 \cref{tab:KW} 中的映射给出~\cite{BenTov2014,Radicevic2018,Li2023,Lootens2022}。等价地，在指定了两个模型中的边界条件之后，可访问的荷区被完全确定。特别请注意拓扑区的置换：在 KW 对偶下，边界条件和荷区本质上发生了互换。对于这些区中的每一个，都存在如下形式的相应 MPO 交织算符：
\begin{equation}
    \inputtikz{KW_MPO_bc}\,\,,
\end{equation}
其中 $\alpha$ 是矩阵 $\{\mbb 1,\sigma^X,\sigma^Y,\sigma^Z\}$ 之一的占位符，亦如 \cref{tab:KW} 所示。
\begin{table}[htb]
    \centering
    \begin{tabular}{cc||cc||c}
        \shortstack{边界\\条件} & \shortstack{荷\\区} & \shortstack{对偶边界\\条件} & \shortstack{对偶荷\\区} & \shortstack{边界\\矩阵} \\
        \hline
        P & $+1$ & P & $+1$ & $\mbb 1$\\
        P & $-1$& AP & $+1$ & $\sigma^Z$\\
        AP & $+1$ & P & $-1$ & $\sigma^X$\\
        AP & $-1$ & AP & $-1$ & $\sigma^Y$
    \end{tabular}.
    \caption{Kramers-Wannier 对偶下拓扑区的映射。此处，\emph{(A)P} 代表\emph{（反）周期性}。}
    \label{tab:KW}
\end{table}

\bigskip\noindent 注意当 $0\leq \lambda <1$ 时，原始模型 \cref{eq:TFIM_per} 处于对称破缺相，而当 $\lambda>1$ 时处于对称相；在 Kramers-Wannier 对偶下，这些相在对偶模型中发生了互换。反转这一逻辑，对于磁场强度 $\lambda$ 的任意取值，总存在一个模型（\cref{eq:TFIM_per} 或 \cref{eq:TFIM_dual}）是（完全）对称破缺的。我们可以推测这一观察普遍成立：通过适当的对偶映射，每个相都可以被映射到相对于对偶对称性完全对称破缺的相。Kennedy-Tasaki 对偶的例子为这一推测提供了进一步的证据。

\paragraph{示例：Kennedy-Tasaki 对偶} \emph{Kennedy-Tasaki}（KT）对偶最初被设计用来阐明自旋 1 反铁磁 Heisenberg 模型的基态物理，该模型现在被理解为展现出非平凡的 $\ZtZt$ SPT 序。该对偶作为开链上 Heisenberg 模型与自发破缺 $\ZtZt$ 对称性模型之间的非幺正映射被引入~\cite{Kennedy1992,Oshikawa1992}。最近，它已根据 $\ZtZt$ 对称性的扭曲 gauging 得到了重新解释~\cite{Lootens2021b,Li2023,ParayilMana2024,Seifnashri2024,Vancraeynest2025a}。等价地，这被描述为输入 $\Vect_{\ZtZt}$ 上的模范畴改变 $\mc R=\Vect\rightarrow\mc R'=\Vect^\psi$，其中 $\psi$ 是 $\ZtZt$ 的非平凡 2-上循环。因此，更准确的理解是将该对称性视为由 $\Rep(\ZtZt)$ 中的特征标标记，该范畴作为融合范畴等价于 $\Vect_{\ZtZt}$。可以证明 $\mc M=\Rep^\psi(\ZtZt)$，因此给定 $\ZtZt$ 的 $\psi$-射影表示 $\pi$，MPO 由具有如下分量的张量构造：
\begin{equation}\label{eq:KT_MPO}
    \inputtikz{KT_MPO} = \delta_{g,h} \pi(g)_{\alpha\beta}.
\end{equation}
长期以来，如何将 KT 从开链推广到闭链一直不甚明确。MPO 框架提供了一个清晰的答案：只需取 MPO 的迹即可将其定义在周期性边界条件上。此外，通过为 MPO \cref{eq:KT_MPO} 选择合适的边界条件，可得到开链上的原始 KT 变换。注意 MPO 张量在其物理腿上是对角的，因此 KT 不改变荷区，而只是置换边界条件。

在此我们不将其应用于 Heisenberg 模型，而是聚焦于 AKLT 态 \cref{eq:AKLTState}，它也是一个非平凡的 $\ZtZt$ SPT。回顾 AKLT 态可以表示为矩阵为 $A^i = \sigma^i/\sqrt 3$ 的 MPS。为明确起见，如果我们取通过 $\pi((a,b))=(\sigma^X)^a(\sigma^Z)^b$ 定义的射影表示，我们发现用该对偶作用在 AKLT 态上将给出一个矩阵为 $A^i=\sigma^i\otimes\sigma^i$ 的 MPS。该 MPS 显然是可约的，因为所有 MPS 矩阵都对易。事实上，利用 gauge 变换
\begin{equation}
    U = \frac{1}{\sqrt 2}
    \begin{pmatrix}
    1 & 0 & 0 & 1\\
    0 & 1 & 1 & 0 \\
    0 & 1 & -1 & 0 \\    
    1 & 0 & 0 & -1
    \end{pmatrix},
\end{equation}
该 MPS 可以约化为四个直积态之和：
\begin{equation}
\begin{split}
    U(\sigma^X\otimes\sigma^X)U^\dagger &= {\rm diag}(1,1,-1,-1), \\
    U(\sigma^Y\otimes\sigma^Y)U^\dagger &= {\rm diag}(-1,1,-1,1), \\
    U(\sigma^Z\otimes\sigma^Z)U^\dagger &= {\rm diag}(1,-1,-1,1).
\end{split}
\end{equation}
这说明对偶算符完全负责 AKLT 态二重简并的纠缠谱。反过来，通过用对偶算符作用在 AKLT 态上（或任何其他 $\ZtZt$ SPT 上），我们可以消除纠缠谱的简并。如下所述，这种观察可以转化为我们的优势。

类比于 KW 的例子，现在可以定义四个对称扭曲边界条件，对应于 $\ZtZt$ 的每个元素。我们请读者参阅文献~\cite{Li2023,Vancraeynest2025a}了解详细讨论。如其中所示，正如预期的那样，荷区被恒等映射，但是根据荷区，边界条件相互置换。

\bigskip\noindent 似乎总是存在唯一的对偶变换将一个模型映射到完全对称破缺的对偶模型，这一观察启示了模范畴与范畴对称有能隙物态之间的一一对应。事实上，这恰好是\emph{广义 Landau 范式}的核心内容。给定对称性 $\mc C$，有能隙物态与 $\mc C$-模范畴 $\mc P$ 一一对应~\cite{Thorngren2019,GarreRubio2022,Lootens2024}。该模范畴 $\mc P$ 既编码了基态 MPS 的不同基态或单射块（对应于其简单对象），又编码了 $\mc C$ 在基态子空间上的作用（通过模作用）。通过利用 parent Hamiltonian 构造，广义 Landau 范式在文献~\cite{GarreRubio2022}中于 MPS 框架下被显式建立。在此背景下，regular 模范畴 $\mc P=\mc C$ 代表完全对称破缺相，其中基态与对称算符一一双射。正如在第~\ref{sec:PhasesOfMatter}讲中已预示的那样，这推广了群对称性的情形，其中模范畴与物态均由配对 ($H\subseteq G,\psi$)，$[\psi]\in H^2(H,\rU(1))$ 所表征。

将广义 Landau 范式与对偶性框架相结合具有深远的影响。首先，注意到对称破缺相是唯一保证存在的有能隙相。\footnote{对于某些对称范畴，例如 $\mc C=\Fib$，它甚至是唯一的有能隙相。}事实上，完全对称相的存在等价于对称性是无反常的，因而完全可被 gauge。将两者相互转换的相应对偶算符的作用实际上是对对称破缺基态取平均~\cite{Lootens2024,Vancraeynest2025b}。令人着迷的是，从这个角度看，任何连续相变都可以被解释为一侧完全对称破缺相与另一侧（部分）对称相之间的转变。

值得注意的是，对称破缺对偶模型使基态纠缠熵最小化。这可以在 DMRG 算法中加以利用，因此对称破缺模型可以被称为\emph{最优}（对偶）模型。事实上，如 \cref{sec:SPT1d} 所示，纠缠谱可以表现出简并性，这是因为纠缠自由度在（射影）虚拟对称性的多重态中变换。在此情况下，基态的 MPS 描述中存在冗余：计算资源被消耗在微调某些变分参数以重现这些谱上。因此，模拟最大程度破缺（对偶）对称性以使得不再存在简并的最优对偶模型是有利的。反过来，原始模型的基态随后通过用对偶算符作用来获得，该算符恢复纠缠谱中的简并。文献~\cite{Lootens2024}中开展了一个案例研究，并证明了该方法的优越性。

确定了最优对偶模型还为任意有能隙相 $\mc P$ 中的初级激发提供了统一的视角。我们现在保留记号 $\mc Q$ 来表示对应于最优模型的关于 $\mc D$ 的右模范畴。正如在 \cref{sec:Excitations} 中所论证的，在这个最优且因此对称破缺相 $\mc D^\star_\mc Q$ 中的初级激发是拓扑畴壁激发。回顾在无限长链极限下，此类激发是通过用 $\mc D^\star_\mc Q$ 中的对称 MPO 作用在半条链上产生的，而其变分自由度编码在其端点处的张量中。在此情况下，准粒子激发组织在一个融合范畴 $\mc E$ 中，该范畴与对称性范畴重合，即 $\mc E=\mc D_\mc Q^\star$。给定此最优模型，其畴壁激发可以反过来映射回原始模型的激发，而后者的激发通常具有完全不同的性质，例如 Heisenberg 模型中的磁振子（magnon）激发（见 \cref{sec:magnon}）。这一情况描述如下：
\begin{equation}\label{eq:duality_excit}
    \inputtikz{Duality_excit} \,\, .
\end{equation}
在此，蓝色 MPO 代表对偶算符——巧合地由 $\mc P$ 中的一个对象标记——位于对应于 $\mc Q$ 的最优模型与感兴趣的模型 $\mc R$ 之间，如标签框所指示。畴壁激发随后描绘在底部，并由张量 $B$ 参数化，类似于 \cref{sec:Excitations}。注意该张量带有一个终止于对偶算符上的荷 $e\in\mc E$。此图像还揭示了 $\mc Q$ 的标签捕获了感兴趣基态的\emph{纠缠自由度}。

通过这种方式，最优对偶模型的畴壁激发为所有其对偶模型的激发提供了不变的刻画。任意模型中的基本激发都可以全部根据 \cref{eq:duality_excit} 进行参数化。

综合起来，这些与准粒子激发和模型相态相关的结构被封装在如下三角形图中，该图重现自文献~\cite{Lootens2024}：
\begin{equation}
    \inputtikz{TriangleDiagram_2}.
\end{equation}

\subsection{第 5 讲概要（Synopsis of lecture V）}
第 5 讲可以被看作是前面各讲的集大成：我们遇到了基本定理、MPO 代数、范畴论、对称破缺以及弦网模型。从张量网络的角度来看，拓扑序系统在其纠缠自由度中表现出对称破缺。这些对称性由构成\emph{融合范畴}表示的 MPO 代数描述。

一方面，这些对称性描述了弦网中的长程纠缠结构，而弦网本身很容易用张量网络表示。作为应用，我们引入了\emph{奇异关联函数}的概念。该方案通过一个由弦网与无特征短程纠缠态的重叠所构成的显式全息构造，将拓扑场论的点阵表示与共形场论联系起来。弦网的 MPO 对称性随后成为 CFT 中的拓扑缺陷，从而为 Fr\"ohlich-Fuchs-Runkel-Schweigert 关于 CFT 中拓扑缺陷的刻画提供了显式的点阵实现。

相同的 MPO 代数还描述了支配这些 (2+1) 维系统边缘模物理的有效哈密顿量的对称性。在一个令人着迷的转折中，这些对称性可以被表示的不同方式——在技术上用双模范畴描述——导出了一个完整的宏伟图像，即通过一组交织 MPO 来定义所有可能的边缘哈密顿量之间的对偶性。该构造最引人入胜的物理推论之一在于，每一个一维物态都对偶于一个（相对于对偶对称性）完全对称破缺的相。值得注意的是，无论该一维系统是定义在具有局域单点对称性的张量积结构上，还是定义在具有范畴非可逆对称性的受约束希尔伯特空间上，这一结论均成立。除其他意义之外，这意味着此类一维系统的每一个可能的二阶相变都等价于向完全对称破缺相的转变。

在这些讲义中，我们尚未涉及 MPO 对称性与 MPO 交织算符的高维版本。有趣的是，在更高维度中，对称性可以作用在点阵不同维度的自由可形变子流形上~\cite{Batista2004,Nussinov2006,Nussinov2009,Gaiotto2014,Freed2022,Lin2022,Roumpedakis2022,McGreevy2022}。这些\emph{高阶形式对称性}（higher-form symmetry）同样自然地用张量网络表示。在 (2+1) 维情形下，面算符用投影纠缠对算符（PEPO）实现，而位于 PEPO 之间界面上的 1-形式对称性则自然地由 MPO 描述。总体而言，该对称性结构及其重新耦合理论被封装在\emph{融合 2-范畴}（fusion 2-category）中~\cite{Douglas2018,Decoppet2024}。类似于一维情形，在点阵上表示融合 2-范畴对称性的不同方式由\emph{模 2-范畴}（module 2-category）分类~\cite{Decoppet2021,Delcamp2021}。

正如文献~\cite{Haegeman2014a}中所开创的，(2+1) 维点阵系统中的对偶性同样可以在点阵上由 PEPO 实现。基于这一张量网络方法，人们在理解高阶范畴对称性及其对偶性方面取得了重大进展~\cite{Delcamp2021,Delcamp2023,Inamura2023,Inamura2025,Eck2025,Vancraeynest2024,Vancraeynest2025a}。尽管取得了这些最新进展，仍有大量有趣的开放研究方向有待探索。

进一步的研究方向包括在该范畴框架中纳入连续对称性。正如文献~\cite{Lootens2024}中所述，\emph{Schur-Weyl 对偶}与用于描述对偶性的 MPO 框架完全兼容。它将具有连续 ${\rm SU}(N)$ 对称性的理论映射为具有离散 $\Rep({\rm SU}(N))$ 对称性的理论。由于 Mermin-Wagner 定理不再适用于后者，对偶映射可用于规避该定理。因此，Schur-Weyl 对偶映射交换了对称相和对称破缺相，从而使该对偶理论的 DMRG 模拟效率大幅提升。

\section*{致谢（Acknowledgements）}
我们衷心感谢组织者 Stephane Ouvry、Tomaž Prosen、Didina Serban 与 Masahito Yamazaki 成功举办了 2025 年 Les Houches 暑期学校 \emph{Exact Solvability and Quantum Information} 这一精彩盛会。此外，我们感谢 Dante Bosgoed、Lander Burgelman、Boris De Vos、Jutho Haegeman、Laurens Lootens、Anton Martin、Zohar Nussinov、Alex Turzillo 以及 Vic Vander Linden 对手稿提出的富有建设性的宝贵意见。

\paragraph{资助信息（Funding information）}
我们鸣谢英国研究与创新署（UKRI）资助项目 EP/Z003342/1、BOF-GOA（项目号 BOF23/GOA/021）、EOS（项目号 40007526）、IBOF（项目号 IBOF23/064），以及法兰德斯研究基金会（FWO）授予 BVDC 的博士奖学金（项目号 11O2423N）。

\begin{appendix}
\numberwithin{equation}{section}

\newpage\section{群上同调（Group cohomology）}\label{sec:GroupCohomology}
在本节中，我们阐述如何利用文献~\cite{Bulmash2020,Vancraeynest2022}中发展的简单线性代数方法，计算任意有限群 $G$（带有在 $\rU(1)$ 上的群作用 $\beta$）的上同调群 $H^n(G,U(1)^\beta)$。作为副产物，我们还将获得每个上同调类的代表元上循环（cocycle）。

\bigskip\noindent 首先介绍我们的记号与相关定义。我们将所有函数 $G^n\rightarrow \rU(1)$ 的集合记作 $C^n(G,\rU(1))$，这些函数称为 \emph{n-余链}（n-cochain）。\emph{余边界映射}（coboundary map）$\partial^{n+1}:C^n(G,\rU(1))\rightarrow C^{n+1}(G,\rU(1))$ 随后定义为：
\begin{equation}
\begin{split}
&(\partial^{n+1}\omega)(g_1,g_2,\dots,g_{n+1}) := \beta_{g_1}(\omega(g_2,g_3,\dots,g_{n+1})) \\
&\hfill\times \prod_{i=1}^n [\omega(g_1,g_2, \dots,g_{i-1},g_ig_{i+1},\dots,g_{n+1})]^{(-1)^i} \times  [\omega(g_1,g_2,\dots,g_n)]^{(-1)^{n + 1}} = 0.
\end{split}
\end{equation}
这里，$\beta$ 表示 $G$ 在（$G$-模）$\rU(1)$ 上的群作用。在对称群包含\emph{反幺正}（antiunitary）元素的情形下（正如正文中所涉及的时间反演与空间反射），当 $g$ 为反幺正时 $\beta_g(\bullet)$ 通过复共轭起作用，否则为恒等作用。容易验证这些映射是幂零的，即 $\partial^{n+1}\circ \partial^n = 1$。随后我们将 \emph{n-上循环}（n-cocycle）与 \emph{n-余边界}（n-coboundary）的阿贝尔群分别定义为：
\begin{align}
    Z^n(G,\rU(1)^\beta) &:= \ker(\partial^{n+1}),\\
    B^n(G,\rU(1)^\beta) &:= {\rm im}(\partial^n).
\end{align}
由余边界映射的幂零性可知，$B^n(G,\rU(1)^\beta)$ 是 $Z^n(G,\rU(1)^\beta)$ 的子群，因此可将 $G$ 的\emph{第 n 个上同调群}定义为商群
\begin{equation}
    H^n(G,\rU(1)^\beta) = Z^n(G,\rU(1)^\beta)/B^n(G,\rU(1)^\beta).
\end{equation}

\bigskip\noindent 计算 $H^n(G,\rU(1)^\beta)$ 并求解代表元上循环的核心思路如下。对\emph{上循环条件} $\partial^{n+1}\omega=1$ 两边取对数，并将未知相位 $\omega(g_1,\dots,g_n)$ 汇集在向量 $\bm{\omega}$ 中，该条件可写为在模 $2\pi$ 意义下必须满足的线性方程组：
\begin{equation}
	\Omega \bm{\omega} = \vec{0}, \mod 2\pi.
	\label{eq:System}
\end{equation}
这里，系数矩阵 $\Omega$ 源自余边界映射 $\partial^{n+1}$，且仅包含整数。$\bm{\omega}$ 的维数是 $|G|^n$，而 $\Omega$ 的维数是 $|G|^{n+1}\times|G|^n$。由于 $\Omega$ 仅包含整数，我们可将 $\Omega$ 写为其 \emph{Smith normal form}：
\begin{equation}
	\Omega = P\Lambda R.
\end{equation}
在此分解中，$P$ 与 $R$ 分别是维数为 $\left|G\right|^{n+1}\times \left|G\right|^{n+1}$ 和 $\left|G\right|^n\times \left|G\right|^n$ 的矩阵，它们仅包含整数且行列式为 1（因此具有整数矩阵逆）。$\Lambda$ 同样仅包含整数，维数为 $\left|G\right|^{n+1}\times \left|G\right|^n$，其形式为
\begin{equation}
	\Lambda
	=
	\begin{pmatrix}
		\begin{array}{c|c}
			\text{diag}\left(d_1,d_2,...,d_r\right) & 0_{r,n-r} \\
		      \hline
		      0_{n+1-r,r} & 0_{n+1-r,n-r}
		\end{array}
	\end{pmatrix}
	,
\end{equation}
其中对角线上的非零元素 $d_1,...,d_r$（其中部分可能为 1）按递增顺序排列 $d_1\leq d_2\leq...$，且每个元素均整除下一个元素 $d_i \mid d_{i+1}$。Smith normal form $\Lambda$ 是唯一的。将此分解代入方程组~(\ref{eq:System})，可得解
\begin{equation}
	\bm{\omega} = 2\pi R^{-1}\Lambda^+\bm{\nu}.
	\label{eq:Sol2Cocyc}
\end{equation}
在上述方程中，$\Lambda^+$ 表示 $\Lambda$ 的（唯一）Moore-Penrose 伪逆，满足 $\Lambda\Lambda^+\Lambda = \Lambda$，其显式形式为
\begin{equation}
	\Lambda^+
	=
	\begin{pmatrix}
		\begin{array}{c|c}    \text{diag}\left(d_1^{-1},d_2^{-1},...,d_r^{-1}\right) & 0_{r,n-r} \\
		\hline
		0_{n+1-r,r} & 0_{n+1-r,n-r}
		\end{array}
	\end{pmatrix}.
\end{equation}
向量 $\bm{\nu}$ 包含任意整数。将解~(\ref{eq:Sol2Cocyc})写为各分量形式为
\begin{equation}
	\omega_i = \sum_{j=1}^r 2\pi\left(R^{-1}\right)_{ij}\frac{\nu_j}{d_j}.
	\label{eq:Sol2CocyBis}
\end{equation}
由于 $\bm{\nu}$ 可自由选取，我们可以依次选取 $\nu_i = \delta_{1,i}, \delta_{2,i},...,\delta_{r,i}$，从而获得解空间的一组基，可写为
\begin{equation}
	(\bm{\omega}_j)_i = \frac{2\pi}{d_j}\left(R^{-1}\right)_{ij}, \quad \forall j\in\{1,...,r\}.
\end{equation}
因此，$n$-上循环 $\bm{\omega}$ 正是 $R^{-1}$ 的各列。由于 $R$ 是满秩的，所有这些解 $\bm{\omega}$ 均线性无关。特别地，（下文的）非平凡上循环不能通过余边界相联系：$\bm{\omega}'=\bm{\omega} + \Omega'\bm{\varphi}$，其中 $\Omega'$ 表示 $n-1$ 阶的余边界映射，$\bm{\varphi}$ 是包含任意实数的 $|G|$ 维向量。为了对所有可能的解进行分类，我们现在考察 $\Lambda$ 的对角元。

由~(\ref{eq:Sol2CocyBis})以及 $R^{-1}$ 仅包含整数的事实可知，对于每一个对角元 $d_i=1$，均得到平凡解 $\omega_i=0 \mod 2\pi$。非平凡解对应于那些 $d_i>1$ 的对角元。由~(\ref{eq:Sol2CocyBis})以及解空间是 $\mathbb{Z}$-线性的事实可知，对应于某个 $d_j$ 的上循环 $\bm{\omega}_j$ 生成一个循环群。注意到并非 $\bm{\omega}_j$ 的所有元素都能被 $d_j$ 或其任意素因子整除，否则将与 $R^{-1}$ 行列式为 1 相矛盾。因此，由 $\bm{\omega}_j$ 生成的循环群正是 $\mathbb{Z}_{d_j}$。最后，$\Lambda$ 中的零元在上同调中亦可舍弃，因为它们对应于上循环方程的平凡解（可乘以任意相位，因而对应于余边界）。

综上所述，我们得到如下图像。给定群 $G$，可以写出余边界映射 $\Omega$ 的矩阵表示，并将其化为 Smith normal form $P\Lambda R$。$\Lambda$ 的对角元 $d_1,d_2,...,d_r$ 决定了上同调群，其形式为 $\mathbb{Z}_{d_1}\times\mathbb{Z}_{d_2}\times ... \times \mathbb{Z}_{d_r}$，其中理解为 $\mathbb{Z}_1=\{e\}$ 表示平凡群，所有为零的对角元均可舍弃。在某种任意 gauge 下，非平凡的 n-上循环对应于 $R^{-1}$ 的各列，并按式~\cref{eq:Sol2CocyBis}乘上适当的权重因子 $2\pi/d_j$。
\newpage\section{范畴论（Category theory）}\label{sec:cat}
在本附录中，我们简要总结正文中用到的融合范畴（fusion category）以及（双）模范畴（(bi)module category）的相关概念。更详尽且技术严谨的阐述请参阅标准文献~\cite{Etingof2002,Etingof2015,Ostrik2001}与综述~\cite{Kitaev2005a,Beer2018,Kong2022}。

\subsection{融合范畴（Fusion categories）}
简要而言，一个（幺正）\emph{融合范畴} $\mc D$ 包含
\begin{enumerate}
    \item 包含\emph{简单对象}（拓扑缺陷）$a,b,c,\dots$ 的有限集合。
    \item 融合规则 $a\otimes b =\bigoplus_c N_{ab}^c c$，以及张量积的单位元 $\mbb 1$。
    \item 称为 \emph{F-符号}（F-symbol）或\emph{结合子}（associator）的幺正映射 $F:(\bullet \otimes\bullet)\otimes\bullet\overset{\sim}{\rightarrow}\bullet\otimes(\bullet\otimes\bullet)$。
\end{enumerate}
F-符号满足称为\emph{五边形方程}（pentagon equation）的自洽条件。展开而言，五边形方程归结为多元线性方程组，形式上示意为
\begin{equation}
    FF = \sum FFF,
\end{equation}
写出各分量形式则为
\begin{equation}
    \sum_o \left(F^{fcd}_e\right)^{h,no}_{g,lm} \left(F^{abh}_e\right)^{i,pq}_{f,ko} = \sum_{j,rst} \left(F^{abc}_g\right)^{j,rs}_{f,kl} \left(F^{ajd}_e\right)^{i,tq}_{g,sm} \left(F^{bcd}_i\right)^{h,np}_{j,rt}.
\end{equation}
值得注意的是，融合范畴不一定配有\emph{编织}（braiding），即幺正映射 $R_{ab}:a\otimes b\overset{\sim}{\rightarrow}b\otimes a$。然而，对于任意融合范畴，均可构造其 \emph{Drinfeld 中心} $\mc Z(\mc D)$，它是一个\emph{模}融合范畴（modular fusion category），特别意味着它配有\emph{非退化}的编织。该中心在 \cref{sec:SN} 的弦网（string-net）模型中可物理实现为模型的任意子（anyon）激发。它们的交换统计编码在 $\mc Z(\mc D)$ 的编织中。

\paragraph{（反常）有限群对称性} 对于任意\emph{有限}群 $G$，融合范畴 $\Vect_G$ 是有限维 $G$-分次向量空间的范畴。本质上，$\Vect_G$ 的简单对象对应于 $G$ 的群元素，融合规则由群乘法决定。此时 $F$-符号全部平凡。当 $G$ 具有非平凡的第三上同调群时（见 \cref{sec:GroupCohomology}），任意 3-上循环（3-cocycle）$\omega$ 均可用于构造 $\Vect_G^\omega$。该融合范畴与 $\Vect_G$ 类似地构造，但此时 $F$-符号的非零分量取值为 $\omega$。此上循环编码了 \emph{（'t Hooft）异常}（anomaly），即对该对称性实施 gauge 的阻碍；或等价地说，正如正文 \cref{sec:SPT1d} 中所讨论的，阻碍该对称性接纳平凡对称基态。

\bigskip\noindent 在这种意义下，融合范畴编码了对称性结构——对称性算符、其乘法规则与结合点——以及由 F-符号所编码的异常。

\paragraph{$\Rep(G)$ 对称性} 对于任意有限群 $G$，范畴 $\Rep(G)$ 将 $G$ 的表示作为对象。不可约表示则构成\emph{简单}对象。这些对象的融合由群表示的标准张量积给出。由此可以推导出，$\Rep(G)$ 的 $F$-符号正是 6j-符号。

\paragraph{Fibonacci 对称性}
Fibonacci 融合范畴具有两个简单对象 $\{\mbb 1,\tau\}$，其唯一的非平凡融合规则为 $\tau\otimes\tau=\mbb 1\oplus\tau$。其 F-符号在文献（如文献~\cite{Beer2018}）中有详细阐述。名称 \emph{Fibonacci} 指的是在 $\tau$ 的逐次张量幂中，$\tau$ 的重数生成了 Fibonacci 数列。同样可以证明，Fibonacci 融合范畴也不允许具有平凡对称基态的局域对称哈密顿量~\cite{GarreRubio2022}。

\subsection{模范畴（Module categories）}
设 $\mc C$ 为融合范畴。一个（幺正左）\emph{$\mc C$-模范畴 $\mc M$} 包含
\begin{enumerate}
    \item 简单对象 $A,B,C,\dots$ 的有限集合。
    \item $\mc C$ 在 $\mc M$ 上的\emph{作用}（action）$\act$：$a\act A =\bigoplus_{B} N_{aA}^{B}B$，其中 $B$ 为 $\mc M$ 中的简单对象。
    \item 称为\emph{模 $\mF$-符号}或（左）\emph{模结合子}的幺正映射 $\mF:(\bullet \otimes\bullet)\act\bullet\overset{\sim}{\rightarrow}\bullet\act(\bullet\act\bullet)$。
\end{enumerate}
这些 $\mF$-符号进而满足一个同时包含 $\mc C$ 的 $F$-符号与模 $\mF$-符号的自洽方程。示意形式为：
\begin{equation}
    \mF\mF = \sum F\mF\mF.
\end{equation}

\paragraph{$\Vect_G$ 的模范畴}
$\Vect_G$ 的不可分解模范畴由对偶对 $(H,\psi)$ 标记，其中 $H\subseteq G$ 为子群（在共轭意义下），$\psi$ 为 $H$ 的 2-上循环。在此情形下，简单对象由陪集 $r_1H,r_2H,\dots\in G/H$ 标记，$\Vect_G$ 在其上的作用定义为 $g\act rH := (gr)H$。模结合子 $\mF$ 可由 $\psi$ 构造，详见文献~\cite{Naidu2006}。

\paragraph{$\Rep(G)$ 的模范畴} 由于 $\Rep(G)$ 与 $\Vect_G$ 的 \emph{Morita 等价}（Morita equivalence），$\Rep(G)$ 上的模范畴同样由元组 $(H,\psi)$ 分类。对应的模范畴为 $\Rep^\psi(H)$，即 $H$ 的 $\psi$-射影表示范畴。$\Rep(G)$ 在 $\Rep^\psi(H)$ 上的作用由将 $G$-表示限制到子群 $H$，随后取（射影）$H$-表示的张量积给出。

\subsection{双模范畴（Bimodule categories）}
给定一对融合范畴 $\mc C,\mc D$，一个 \emph{$(\mc C,\mc D)$-双模范畴} $\mc M$ 是既为左 $\mc C$-模范畴又为右 $\mc D$-模范畴的范畴，并配有额外的幺正\emph{双模 $\rF$-符号}或\emph{双模结合子} $\rF:(\bullet \act \bullet)\cat\bullet \overset{\sim}{\rightarrow} \bullet\act(\bullet\cat\bullet)$。

\paragraph{$\Vect$} $\Vect$ 可解释为 $(\Vect_G,\Rep(G))$-双模范畴。在此情况下，$\rF$-符号的分量即为 $G$-表示的矩阵元。

\bigskip\noindent 在整篇讲义中，我们专门处理\emph{可逆} $(\mc C,\mc D)$-双模范畴。特别是，这保证了 $\mc C$ 与 $\mc D$ 是 \emph{Morita 等价}的，或者等价地说，它们的 Drinfeld 中心 $\mc Z(\mc C)$ 与 $\mc Z(\mc D)$ 是等价的模融合范畴~\cite{Mueger2002}。在文献~\cite{Bridgeman2022}中，可逆性被证明等价于 MPO 单射性（MPO injectivity）的概念，并推导出了用双模结合子给出的刻画。

给定融合范畴 $\mc C$ 与 $\mc C$-模范畴 $\mc M$，存在唯一的融合范畴——\emph{Morita 对偶}（Morita dual）$\mc C^\star_\mc M$，使得 $\mc M$ 成为可逆 $(\mc C,\mc C^\star_\mc M)$-双模范畴。例如：$(\Vect_G)^\star_\Vect=\Rep(G)$。关于 Morita 对偶融合范畴的严格定义，我们建议读者查阅相关文献。

\end{appendix}
\bibliography{refs.bib}

\end{document}